% Générateur et récepteur électriques — Physique 1ère C — Cours signé OnBuch+
% Programme officiel MINESEC. Compiler avec : tectonic N14-generateur-recepteur-electriques.tex
\newcommand{\DOCMATIERE}{Physique}
\newcommand{\DOCNIVEAU}{1ère C}
\newcommand{\DOCMODULE}{Module 4 : Aspects énergétiques des circuits électriques}
\newcommand{\DOCLECON}{N14}
\newcommand{\DOCTITRE}{Générateur et récepteur électriques}
% OnBuch+ — préambule « cours signés » ENRICHI (compilé avec tectonic / XeLaTeX)
% Système visuel « Pop » d'OnBuch+ : fond crème (jamais blanc), bordures encre
% épaisses, ombres « dures » décalées, titres Archivo Black en capitales.
% Version MATHÉMATIQUES : graphes (pgfplots/tikz), boîtes pédagogiques
% (définition, propriété, méthode, attention, le savais-tu, exercice résolu,
% exercice, TP, bilan), page de garde et signature OnBuch+ ancrée en bas.
%
% Macros attendues dans main.tex AVANT \input{preamble} :
%   \DOCMATIERE  \DOCNIVEAU  \DOCMODULE  \DOCLECON (numéro)  \DOCTITRE
\documentclass[11pt]{article}

\usepackage{fontspec}
\usepackage[french]{babel}
\usepackage[a4paper,margin=2cm,top=3.4cm,bottom=2.8cm,headheight=1.4cm,footskip=1.3cm]{geometry}
\usepackage{amsmath,amssymb,amsfonts}
\usepackage{mathtools} % \xmapsto, \coloneqq... souvent utilisés par les modèles
\usepackage{cancel} % simplifications barrées (bilans de forces)
% \abs{z} (module, valeur absolue) et \norm{u} (norme), utilisés par les modèles
\providecommand{\abs}{}\let\abs\relax\DeclarePairedDelimiter{\abs}{\lvert}{\rvert}
\providecommand{\norm}{}\let\norm\relax\DeclarePairedDelimiter{\norm}{\lVert}{\rVert}
\usepackage[table]{xcolor} % [table] : fournit aussi \rowcolors (alternance de lignes)
\usepackage{pagecolor}
\usepackage{tikz}
\usetikzlibrary{arrows.meta,positioning,calc,shapes.geometric,shapes.misc,shapes.arrows,decorations.pathmorphing,decorations.pathreplacing,decorations.markings,patterns,shadows,mindmap,trees,fit,backgrounds,matrix,intersections,angles,quotes}
\usepackage{pgfplots}
\usepackage[version=4]{mhchem} % équations nucléaires \ce{^{235}_{92}U}
\usepackage[european,siunitx]{circuitikz} % schémas électriques
\pgfplotsset{compat=1.18}
\usepgfplotslibrary{fillbetween,groupplots} % aires entre courbes (calcul intégral)
\usepackage{enumitem}
\usepackage{fancyhdr}
% [most] charge toutes les bibliothèques tcolorbox (listings, minted, raster,
% documentation...) et gonfle inutilement la mémoire TeX sur un document long
% (~300 boîtes) : on ne charge que ce qui est réellement utilisé.
\usepackage[skins,breakable]{tcolorbox}
\usepackage{graphicx}
\usepackage{colortbl}
\usepackage{booktabs}
\usepackage{tabularx}
\usepackage{diagbox} % case d'en-tête diagonale des tableaux à double entrée
% Colonnes extensibles centrée / gauche / droite (C, L, R), souvent utilisées par les modèles
\newcolumntype{C}{>{\centering\arraybackslash}X}
\newcolumntype{L}{>{\raggedright\arraybackslash}X}
\newcolumntype{R}{>{\raggedleft\arraybackslash}X}
\usepackage{array}
\usepackage{multicol}
\usepackage{siunitx}
% parse-numbers=false : accepte aussi \SI{2,5 \times 10^{-3}}{...}
\sisetup{locale=FR,output-decimal-marker={,},group-minimum-digits=4,parse-numbers=false}
\DeclareSIUnit\ohm{\ensuremath{\symup{\Omega}}} % U+2126 absent de la police
\DeclareSIUnit\division{div} % réglages d'oscilloscope (ms/div, V/div)
\DeclareSIUnit\tr{tr} % tours (vitesse de rotation, ex. \SI{1500}{\tr\per\minute})
\DeclareSIUnit\ch{ch} % cheval-vapeur (puissance moteur, unité usuelle hors SI)
\DeclareSIUnit\franccfa{FCFA} % franc CFA (coûts, contexte camerounais)
\DeclareSIUnit\dioptre{\ensuremath{\delta}} % vergence d'une lentille (dioptrie)
\usepackage{qrcode}
% \degree : commande fréquemment utilisée par les modèles pour un angle ou une
% température en dehors de \SI{}{} (non fournie par les paquets chargés).
\providecommand{\degree}{^{\circ}}
% \qed (fin de démonstration), utilisé par les modèles sans amsthm
\providecommand{\qed}{\hfill$\square$}
% \barre : double barre des valeurs interdites dans un tableau de variations (tkz-tab)
\providecommand{\barre}{\ensuremath{\|}}
% environnement proof (amsthm non chargé) utilisé par les modèles
\ifdefined\proof\else\newenvironment{proof}[1][Démonstration]{\par\noindent\textit{#1.}\ }{\qed\par}\fi
\usepackage{lastpage}
\usepackage{hyperref}
\hypersetup{hidelinks}

% ============================================================
% Palette « Pop » OnBuch+ — mêmes hex que la classe P (plus_ui.dart)
% ============================================================
\definecolor{poporange}{HTML}{F59321}
\definecolor{popdark}{HTML}{E07A0C}
\definecolor{popcream}{HTML}{FFF6E8}
\definecolor{popink}{HTML}{1C1714}
\definecolor{poppurple}{HTML}{7A5AE0}
\definecolor{popgreen}{HTML}{2E9E6A}
\definecolor{poppink}{HTML}{E0457B}
\definecolor{popblue}{HTML}{2F7DE1}
\definecolor{popgold}{HTML}{C98A12}
\definecolor{popmuted}{HTML}{8A7F76}
\definecolor{popline}{HTML}{EADFD2}
\definecolor{popgray}{HTML}{8A7F76}
\colorlet{popgrey}{popgray}
% Alias tolérants (le modèle nomme parfois les couleurs différemment)
\colorlet{popwhite}{white}
\colorlet{popblack}{popink}
\colorlet{popred}{poppink}
\colorlet{popyellow}{popgold}
% Teintes claires pour fonds de tableaux / zones
\colorlet{popblueL}{popblue!10!white}
\colorlet{poporangeL}{poporange!14!white}
\colorlet{popgreenL}{popgreen!12!white}
\colorlet{poppurpleL}{poppurple!12!white}
\colorlet{poppinkL}{poppink!12!white}
\colorlet{popgoldL}{popgold!14!white}

\pagecolor{popcream}

% ============================================================
% Typographie
% ============================================================
\defaultfontfeatures{Ligatures=TeX}
\setmainfont{PlusJakartaSans-Medium.ttf}[Path=../../../fonts/,BoldFont=PlusJakartaSans-Bold.ttf]
% Maths en sans-serif (Fira Math) avec chiffres et lettres droites de Plus
% Jakarta, pour que les formules se fondent dans le texte.
\usepackage[math-style=ISO,bold-style=ISO]{unicode-math}
\setmathfont{FiraMath-Regular.otf}[Path=../../../fonts/]
\setmathfont{PlusJakartaSans-Medium.ttf}[Path=../../../fonts/,range={up/num,up/{Latin,latin},\mathup}]
\usepackage{icomma} % virgule décimale sans espace en mode math
% Caractères mathématiques Unicode absents des polices de texte (Plus Jakarta,
% Archivo Black) : rendus par leur équivalent mathématique, en texte comme en
% formule — sinon ils s'affichent en carré vide (ex. « Arithmétique dans ℤ »).
\usepackage{newunicodechar}
\newunicodechar{ℝ}{\ensuremath{\mathbb{R}}}
\newunicodechar{ℤ}{\ensuremath{\mathbb{Z}}}
\newunicodechar{ℂ}{\ensuremath{\mathbb{C}}}
\newunicodechar{ℕ}{\ensuremath{\mathbb{N}}}
\newunicodechar{ℚ}{\ensuremath{\mathbb{Q}}}
\newunicodechar{𝔻}{\ensuremath{\mathbb{D}}}
\newunicodechar{α}{\ensuremath{\alpha}}
\newunicodechar{β}{\ensuremath{\beta}}
\newunicodechar{γ}{\ensuremath{\gamma}}
\newunicodechar{δ}{\ensuremath{\delta}}
\newunicodechar{ε}{\ensuremath{\varepsilon}}
\newunicodechar{θ}{\ensuremath{\theta}}
\newunicodechar{λ}{\ensuremath{\lambda}}
\newunicodechar{μ}{\ensuremath{\mu}}
\newunicodechar{σ}{\ensuremath{\sigma}}
\newunicodechar{τ}{\ensuremath{\tau}}
\newunicodechar{φ}{\ensuremath{\varphi}}
\newunicodechar{ψ}{\ensuremath{\psi}}
\newunicodechar{ω}{\ensuremath{\omega}}
\newunicodechar{Σ}{\ensuremath{\Sigma}}
% majuscules produites par \MakeUppercase dans les titres (α-aminés -> Α)
\newunicodechar{Α}{\ensuremath{\alpha}}
\newunicodechar{Β}{\ensuremath{\beta}}
\newunicodechar{Γ}{\ensuremath{\Gamma}}
\newunicodechar{Δ}{\ensuremath{\Delta}}
\newunicodechar{Ε}{\ensuremath{\varepsilon}}
\newunicodechar{Θ}{\ensuremath{\Theta}}
\newunicodechar{Λ}{\ensuremath{\Lambda}}
\newunicodechar{Μ}{\ensuremath{\mu}}
\newunicodechar{Τ}{\ensuremath{\tau}}
\newunicodechar{Φ}{\ensuremath{\Phi}}
\newunicodechar{Ψ}{\ensuremath{\Psi}}
\newunicodechar{Ω}{\ensuremath{\Omega}}
\newunicodechar{↦}{\ensuremath{\mapsto}}
\newunicodechar{∈}{\ensuremath{\in}}
\newunicodechar{∉}{\ensuremath{\notin}}
\newunicodechar{∧}{\ensuremath{\wedge}}
\newunicodechar{⇔}{\ensuremath{\Leftrightarrow}}
\newunicodechar{⟺}{\ensuremath{\Longleftrightarrow}}
\newunicodechar{⇒}{\ensuremath{\Rightarrow}}
\newunicodechar{⟹}{\ensuremath{\Longrightarrow}}
\newunicodechar{∘}{\ensuremath{\circ}}
\newunicodechar{∪}{\ensuremath{\cup}}
\newunicodechar{∩}{\ensuremath{\cap}}
\newunicodechar{≡}{\ensuremath{\equiv}}
\newunicodechar{⊂}{\ensuremath{\subset}}
\newunicodechar{⊥}{\ensuremath{\perp}}
\newunicodechar{∃}{\ensuremath{\exists}}
\newunicodechar{∀}{\ensuremath{\forall}}
\newunicodechar{∛}{\ensuremath{\sqrt[3]{\,}}}
\newunicodechar{∜}{\ensuremath{\sqrt[4]{\,}}}
\newunicodechar{⃗}{\ensuremath{{}^{\to}}}
\newunicodechar{ⁿ}{\ifmmode^{n}\else\textsuperscript{\ensuremath{n}}\fi}
\newunicodechar{ˣ}{\ifmmode^{x}\else\textsuperscript{\ensuremath{x}}\fi}
\newunicodechar{ᵇ}{\ifmmode^{b}\else\textsuperscript{\ensuremath{b}}\fi}
\newunicodechar{ʳ}{\ifmmode^{r}\else\textsuperscript{\ensuremath{r}}\fi}
\newunicodechar{ᵘ}{\ifmmode^{u}\else\textsuperscript{\ensuremath{u}}\fi}
\newunicodechar{ʸ}{\ifmmode^{y}\else\textsuperscript{\ensuremath{y}}\fi}
\newunicodechar{ᵃ}{\ifmmode^{a}\else\textsuperscript{\ensuremath{a}}\fi}
\newunicodechar{ᵉ}{\ifmmode^{e}\else\textsuperscript{\ensuremath{e}}\fi}
\newunicodechar{ᵏ}{\ifmmode^{k}\else\textsuperscript{\ensuremath{k}}\fi}
\newunicodechar{ᵖ}{\ifmmode^{p}\else\textsuperscript{\ensuremath{p}}\fi}
\newunicodechar{ᴺ}{\ifmmode^{N}\else\textsuperscript{\ensuremath{N}}\fi}
\newunicodechar{ᶜ}{\ifmmode^{c}\else\textsuperscript{\ensuremath{c}}\fi}
\newunicodechar{ʰ}{\ifmmode^{h}\else\textsuperscript{\ensuremath{h}}\fi}
\newunicodechar{ᵠ}{\ifmmode^{q}\else\textsuperscript{\ensuremath{q}}\fi}
\newunicodechar{ₙ}{\ifmmode_{n}\else\textsubscript{\ensuremath{n}}\fi}
\newunicodechar{ₐ}{\ifmmode_{a}\else\textsubscript{\ensuremath{a}}\fi}
\newunicodechar{ᵢ}{\ifmmode_{i}\else\textsubscript{\ensuremath{i}}\fi}
\newunicodechar{ᵣ}{\ifmmode_{r}\else\textsubscript{\ensuremath{r}}\fi}
\newunicodechar{ₖ}{\ifmmode_{k}\else\textsubscript{\ensuremath{k}}\fi}
\newunicodechar{ᵦ}{\ifmmode_{\beta}\else\textsubscript{\ensuremath{\beta}}\fi}
\newunicodechar{ₓ}{\ifmmode_{x}\else\textsubscript{\ensuremath{x}}\fi}
\newfontfamily\popdisplay{ArchivoBlack-Regular.ttf}[Path=../../../fonts/]
\newfontfamily\poplogo{SpaceGrotesk-Bold.ttf}[Path=../../../fonts/]
\newfontfamily\popsemi{PlusJakartaSans-SemiBold.ttf}[Path=../../../fonts/]
\newfontfamily\popxbold{PlusJakartaSans-ExtraBold.ttf}[Path=../../../fonts/]
\newfontfamily\popxboldls{PlusJakartaSans-ExtraBold.ttf}[Path=../../../fonts/,LetterSpace=3.0]

\setlist[enumerate]{leftmargin=1.7em,itemsep=5pt,topsep=4pt}
\setlist[itemize]{leftmargin=1.7em,itemsep=4pt,topsep=4pt,label={\color{poporange}\textbullet}}
\setlength{\parindent}{0pt}
\setlength{\parskip}{5pt}
\renewcommand{\arraystretch}{1.25}

% pgfplots : style Pop par défaut
\pgfplotsset{
  every axis/.append style={axis line style={popink,line width=1pt},tick style={popink},
    grid style={popline,line width=0.5pt},label style={font=\small},tick label style={font=\footnotesize}},
  popcurve/.style={poporange,line width=1.8pt,no markers,smooth},
}
% Même style disponible dans un \draw TikZ simple (hors pgfplots)
\tikzset{popcurve/.style={poporange,line width=1.8pt}}
\tikzset{popgrid/.style={popline,very thin}}
\colorlet{popcurve}{poporange} % le nom du style est parfois utilisé comme couleur

% ============================================================
% Pill / squiggle
% ============================================================
\newcommand{\poppill}[3]{%
  \tikz[baseline=-0.55ex]\node[draw=popink,line width=1.2pt,rounded corners=2.6pt,%
    fill=#2,inner xsep=6.5pt,inner ysep=3pt]{%
    \popxboldls\fontsize{7}{7}\selectfont\color{#3}\MakeUppercase{#1}};%
}
\newcommand{\popsquiggle}[1]{%
  \tikz[baseline=-0.4ex]{
    \draw[#1,line width=2.6pt,line cap=round]
      plot [smooth] coordinates {(0,0.09) (0.34,0.01) (0.68,0.21) (1.02,0.01) (1.36,0.10)};
  }%
}
% Mot-clé surligné (marqueur orange sous le texte)
\newcommand{\cle}[1]{{\popxbold\color{popdark}#1}}

% ============================================================
% En-tête / pied
% ============================================================
\pagestyle{fancy}
\fancyhf{}
\renewcommand{\headrulewidth}{0pt}
\renewcommand{\footrulewidth}{0pt}
\lhead{\raisebox{-0.15ex}{\poplogo\fontsize{13}{13}\selectfont\color{popink}OnBuch\color{poporange}+}}
\rhead{\poppill{\DOCMATIERE\ \textbullet\ \DOCNIVEAU\ \textbullet\ Leçon \DOCLECON}{white}{popink}}
\lfoot{\popxbold\fontsize{7.6}{7.6}\selectfont\color{popmuted}\MakeUppercase{Cours signé OnBuch+}\ \textbullet\ {\popsemi\DOCTITRE}}
\rfoot{\tikz[baseline=-0.6ex]\node[rounded corners=8.5pt,fill=popink,minimum height=17pt,minimum width=17pt,inner xsep=5pt,inner sep=0]{\popxbold\fontsize{8}{8}\selectfont\color{popcream}\thepage\,/\,\pageref*{LastPage}};}

% ============================================================
% Titres
% ============================================================
\newcommand{\chaptitre}[2]{%
  \begin{tcolorbox}[enhanced,colback=poporange,colframe=popink,boxrule=2.6pt,arc=11pt,
      drop shadow=popink,left=15pt,right=15pt,top=12pt,bottom=11pt,before skip=3pt,after skip=18pt]
    \poppill{#2}{popink}{popcream}\par\vspace{9pt}
    {\raggedright\hyphenpenalty=10000\exhyphenpenalty=10000\popdisplay\fontsize{22}{24}\selectfont\color{popcream}\MakeUppercase{#1}\par}
  \end{tcolorbox}
}
% Section numérotée : pastille orange avec le numéro + titre Archivo
\newcounter{popsec}
\newcommand{\coursec}[1]{%
  \stepcounter{popsec}\setcounter{popsub}{0}%
  \par\vspace{16pt}\noindent
  \begin{minipage}{\linewidth}
  \tikz[baseline=-0.7ex]\node[draw=popink,line width=1.6pt,fill=poporange,rounded corners=4pt,minimum size=22pt,inner sep=0]{\popdisplay\fontsize{12}{12}\selectfont\color{popcream}\thepopsec};%
  \hspace{8pt}{\popdisplay\fontsize{15}{17}\selectfont\color{popink}\MakeUppercase{#1}}\par
  \vspace{4pt}\noindent{\color{poporange}\rule{36pt}{3.6pt}}
  \end{minipage}\par\nopagebreak\vspace{8pt}
}
\newcounter{popsub}
\newcommand{\courssub}[1]{%
  \stepcounter{popsub}%
  \par\vspace{9pt}\noindent{\popxbold\fontsize{11.5}{13}\selectfont\color{popink}{\color{poporange}\thepopsec.\thepopsub}\hspace{6pt}#1}\par\nopagebreak\vspace{4pt}
}

% ============================================================
% Boîtes pédagogiques — même recette : fond blanc, bordure épaisse colorée,
% ombre dure encre, badge plein. Titre optionnel [#1].
% ============================================================
\tcbset{popbase/.style={breakable,enhanced,colback=white,boxrule=2.3pt,arc=9pt,
  shadow={3pt}{-3pt}{0pt}{popink},left=14pt,right=14pt,top=11pt,bottom=11pt,before skip=11pt,after skip=15pt,
  fontupper=\popsemi\fontsize{10.6}{14.5}\selectfont\color{popink},
  fontlower=\popsemi\fontsize{10.6}{14.5}\selectfont\color{popink}}}
% Coche dessinée (le glyphe ✓ n'existe pas dans les polices du document)
\newcommand{\popcheck}[1]{\tikz[baseline=-0.5ex]\draw[#1,line width=1.6pt,line cap=round,line join=round](-0.13,0.0)--(-0.04,-0.09)--(0.13,0.1);}
\providecommand{\coche}{\popcheck{popgreen}}
\providecommand{\resultat}[1]{\par\smallskip\noindent\fcolorbox{popgreen}{popgreenL}{\parbox{\dimexpr\linewidth-2\fboxsep-2\fboxrule\relax}{\textbf{Résultat.} #1}}\par\smallskip}
\newcommand{\popboxhead}[3]{\poppill{#1}{#2}{white}\ifx\relax#3\relax\else\hspace{6pt}{\popxbold\fontsize{10.6}{12}\selectfont\color{popink}#3}\fi\par\vspace{7pt}}

\newtcolorbox{aretenir}[1][]{popbase,colframe=popgreen,before upper={\popboxhead{À retenir}{popgreen}{#1}}}
\newtcolorbox{exemplebox}[1][]{popbase,colframe=poporange,before upper={\popboxhead{Exemple}{poporange}{#1}}}
\newtcolorbox{definition}[1][]{popbase,colframe=popblue,before upper={\popboxhead{Définition}{popblue}{#1}}}
\newtcolorbox{propriete}[1][]{popbase,colframe=poppurple,before upper={\popboxhead{Propriété}{poppurple}{#1}}}
\newtcolorbox{methode}[1][]{popbase,colframe=popgold,before upper={\popboxhead{Méthode}{popgold}{#1}}}
\newtcolorbox{attention}[1][]{popbase,colframe=poppink,before upper={\popboxhead{Attention}{poppink}{#1}}}
\newtcolorbox{experience}[1][]{popbase,colframe=popblue,colback=popblueL,before upper={\popboxhead{Expérience}{popblue}{#1}}}
% « Le savais-tu ? » : carte violette pleine (contexte, culture, Cameroun)
\newtcolorbox{savaistu}[1][]{popbase,colback=poppurple,colframe=popink,
  fontupper=\popsemi\fontsize{10.4}{14.2}\selectfont\color{white},
  before upper={\poppill{Le savais-tu\,?}{white}{poppurple}\ifx\relax#1\relax\else\hspace{6pt}{\popxbold\color{white}#1}\fi\par\vspace{7pt}}}
% Exercice résolu : énoncé en haut, \tcblower puis la solution sur fond crème
\newcounter{popexo}\newcounter{popexr}
\newtcolorbox{exoresolu}[1][]{popbase,colframe=poporange,colbacklower=popcream,
  segmentation style={popink,line width=1.2pt,dashed},
  before upper={\stepcounter{popexr}\popboxhead{Exercice résolu \thepopexr}{poporange}{#1}},
  before lower={\poppill{Solution}{popink}{popcream}\par\vspace{6pt}}}
% Exercice (non résolu en place) — la correction est dans la partie Corrigés
\newtcolorbox{exercice}[1][]{popbase,colframe=popink,
  before upper={\stepcounter{popexo}\popboxhead{Exercice \thepopexo}{popink}{#1}}}
% Corrigé
\newtcolorbox{corrige}[1][]{popbase,colframe=popgreen,colback=popgreenL,
  before upper={\popboxhead{Corrigé}{popgreen}{#1}}}
% Objectifs / prérequis / bilan
\newtcolorbox{objectifs}{popbase,colframe=popink,colback=poporangeL,
  before upper={\popboxhead{Objectifs de la leçon}{popink}{}}}
\newtcolorbox{prerequis}{popbase,colframe=popink,colback=popblueL,
  before upper={\popboxhead{Prérequis}{popblue}{}}}
\newtcolorbox{bilan}{popbase,colframe=popink,colback=white,boxrule=2.8pt,
  before upper={\popboxhead{Fiche bilan}{poporange}{L'essentiel à connaître pour l'examen}}}
% Figure encadrée avec légende
\newtcolorbox{popfigure}[1][]{enhanced,breakable=false,colback=white,colframe=popline,boxrule=1.4pt,arc=8pt,
  left=10pt,right=10pt,top=10pt,bottom=8pt,before skip=10pt,after skip=12pt,halign=center,
  lower separated=false,halign lower=center,
  fontlower=\popsemi\fontsize{9}{11.5}\selectfont\color{popmuted},
  #1}
\newcommand{\legende}[1]{\par\vspace{6pt}{\popsemi\fontsize{9}{11.5}\selectfont\color{popmuted}#1}}

% ============================================================
% Signature OnBuch+
% ============================================================
\newcommand{\poplogoinline}[1]{{\poplogo\fontsize{#1}{#1}\selectfont\color{popink}OnBuch\color{poporange}+}}
\newcommand{\signatureonbuch}{%
  \begin{tcolorbox}[enhanced,colback=popink,colframe=popink,boxrule=2.2pt,arc=10pt,
      drop shadow=poporange,left=14pt,right=14pt,top=10pt,bottom=10pt,before skip=0pt,after skip=0pt]
    \tikz[baseline=-0.7ex]\node[circle,fill=popgreen,draw=popcream,line width=1.3pt,minimum size=22pt,inner sep=0]{\popcheck{white}};%
    \hspace{9pt}\begin{minipage}[c]{0.62\linewidth}
      {\popxbold\fontsize{12}{13}\selectfont\color{popcream}Signé\ \textbullet\ {\poplogo OnBuch\color{poporange}+}}\par\vspace{2pt}
      {\raggedright\popsemi\fontsize{8}{10.5}\selectfont\color{popline}\MakeUppercase{Équipe pédagogique OnBuch+}\\\MakeUppercase{Conforme au programme officiel MINESEC}\par}
    \end{minipage}\hfill
    {\poplogo\fontsize{22}{22}\selectfont\color{popcream}OnBuch\color{poporange}+}
  \end{tcolorbox}%
}

% ============================================================
% Page de garde
% #1 titre · #2 matière · #3 niveau · #4 module · #5 n° leçon · #6 durée ·
% #7 description / objectifs (texte ou itemize)
% ============================================================
\newcommand{\pagegarde}[7]{%
  \thispagestyle{empty}
  \newgeometry{a4paper,margin=1.7cm,top=1.6cm,bottom=1.4cm}%
  \begin{tikzpicture}[remember picture,overlay]
    \fill[poporange!18] ([xshift=-3.2cm,yshift=-3.6cm]current page.north east) circle (4.6cm);
    \fill[poppurple!14] ([xshift=2.4cm,yshift=7.6cm]current page.south west) circle (3.8cm);
  \end{tikzpicture}%
  {\poplogo\fontsize{30}{30}\selectfont\color{popink}OnBuch\color{poporange}+}\hfill
  \raisebox{0.6ex}{\poppill{Cours signé \textbullet\ Premium}{popink}{popcream}}\par
  \vspace{1pt}\popsquiggle{poppurple}\par
  \vspace{8pt}
  \begin{tcolorbox}[enhanced,colback=poporange,colframe=popink,boxrule=2.8pt,arc=13pt,
      drop shadow=popink,left=18pt,right=18pt,top=12pt,bottom=13pt,after skip=10pt]
    \poppill{#2 \textbullet\ #3}{popink}{popcream}\hspace{5pt}\poppill{#4}{white}{popink}\par\vspace{8pt}
    {\popdisplay\fontsize{14}{14}\selectfont\color{popink}LEÇON #5}\par\vspace{4pt}
    {\raggedright\hyphenpenalty=10000\exhyphenpenalty=10000\popdisplay\fontsize{25}{27}\selectfont\color{popcream}\MakeUppercase{#1}\par}
    \vspace{8pt}
    {\popxbold\fontsize{9.5}{9.5}\selectfont\color{popink}\MakeUppercase{Durée conseillée : #6}}
  \end{tcolorbox}
  \begin{tcolorbox}[enhanced,colback=white,colframe=popink,boxrule=2.2pt,arc=10pt,
      drop shadow=popink,left=16pt,right=16pt,top=10pt,bottom=10pt,after skip=8pt]
    \poppill{Dans cette leçon}{poppurple}{white}\par\vspace{5pt}
    \popsemi\fontsize{9.3}{12.6}\selectfont\color{popink}#7
  \end{tcolorbox}
  \noindent\begin{minipage}[t]{0.487\linewidth}
    \begin{tcolorbox}[enhanced,colback=popgreen,colframe=popink,boxrule=2.2pt,arc=10pt,
        drop shadow=popink,left=11pt,right=11pt,top=10pt,bottom=10pt,height=3.1cm,valign=center]
      \tcbox[colback=white,colframe=popink,boxrule=1.3pt,arc=5pt,boxsep=0pt,nobeforeafter,
        left=3pt,right=3pt,top=3pt,bottom=3pt,valign=center]{\qrcode[height=1.9cm]{https://play.google.com/store/apps/details?id=com.luvvix.onbuch&hl=fr}}%
      \hspace{8pt}\begin{minipage}[c]{0.5\linewidth}
        {\popdisplay\fontsize{11}{12.5}\selectfont\color{white}\MakeUppercase{Télécharge OnBuch}}\par\vspace{4pt}
        {\raggedright\popsemi\fontsize{8.4}{11}\selectfont\color{white}Gratuit \textbullet\ Google Play\\Tuteur IA, annales, concours, résultats\par}
      \end{minipage}
    \end{tcolorbox}
  \end{minipage}\hfill
  \begin{minipage}[t]{0.487\linewidth}
    \begin{tcolorbox}[enhanced,colback=poppurple,colframe=popink,boxrule=2.2pt,arc=10pt,
        drop shadow=popink,left=11pt,right=11pt,top=10pt,bottom=10pt,height=3.1cm,valign=center]
      \tikz[baseline=-0.7ex]\node[circle,fill=white,draw=popink,line width=1.3pt,minimum size=1.6cm,inner sep=0]{\popdisplay\fontsize{24}{24}\selectfont\color{poppurple}+};%
      \hspace{8pt}\begin{minipage}[c]{0.6\linewidth}
        {\popdisplay\fontsize{11}{12.5}\selectfont\color{white}\MakeUppercase{Passe à OnBuch+}}\par\vspace{4pt}
        {\raggedright\popsemi\fontsize{8.4}{11}\selectfont\color{white}Tous les cours signés, Tuteur IA illimité, fiches PDF — onglet OnBuch+ de l'app\par}
      \end{minipage}
    \end{tcolorbox}
  \end{minipage}
  \vspace{8pt}
  \popcontenu
  \vfill
  \signatureonbuch
  \restoregeometry
  \newpage
}

% Rangée « Ce cours contient » (page de garde)
\newcommand{\poptile}[3]{% #1 couleur #2 gros texte #3 légende
  \begin{tcolorbox}[enhanced,colback=#1,colframe=popink,boxrule=2pt,arc=9pt,shadow={2.5pt}{-2.5pt}{0pt}{popink},
      left=6pt,right=6pt,top=9pt,bottom=9pt,halign=center,valign=center,height=1.75cm,width=\linewidth,nobeforeafter]
    {\popdisplay\fontsize{12.5}{13}\selectfont\color{white}\MakeUppercase{#2}}\par\vspace{3pt}
    {\centering\popsemi\fontsize{7.8}{9.5}\selectfont\color{white}#3\par}
  \end{tcolorbox}}
\newcommand{\popcontenu}{%
  \par\noindent{\popxboldls\fontsize{7.5}{7.5}\selectfont\color{popmuted}CE COURS SIGNÉ CONTIENT}\par\vspace{7pt}
  \noindent\begin{minipage}[t]{0.235\linewidth}\poptile{popblue}{Cours}{complet, illustré, pas à pas}\end{minipage}\hfill
  \begin{minipage}[t]{0.235\linewidth}\poptile{poporange}{Méthodes}{et exercices résolus}\end{minipage}\hfill
  \begin{minipage}[t]{0.235\linewidth}\poptile{poppink}{Exercices}{type examen + corrigés}\end{minipage}\hfill
  \begin{minipage}[t]{0.235\linewidth}\poptile{popgreen}{Fiche bilan}{l'essentiel à retenir}\end{minipage}\par}

% Sommaire
\newtcolorbox{sommaire}{enhanced,colback=white,colframe=popink,boxrule=2.2pt,arc=10pt,
  drop shadow=popink,left=18pt,right=18pt,top=13pt,bottom=13pt,before skip=6pt,after skip=20pt,
  before upper={\poppill{Sommaire}{poppurple}{white}\par\vspace{9pt}\popsemi\fontsize{10.6}{14.5}\selectfont\color{popink}}}

% Signature de fin de cours — poussée tout en bas de la dernière page
\newcommand{\findecours}{%
  \par\vfill\nopagebreak
  \begin{center}\popsquiggle{poporange}\end{center}\vspace{4pt}
  \signatureonbuch
}
\AtBeginDocument{\def\llbracket{\mathopen{[\mkern-3.5mu[}}\def\rrbracket{\mathclose{]\mkern-3.5mu]}}} % intervalles d'entiers (glyphe ⟦ absent de la police)
% Glyphes absents de Fira Math : redéfinis à partir de glyphes disponibles
\AtBeginDocument{%
  \def\setminus{\mathbin{\backslash}}\let\smallsetminus\setminus
  \def\vdots{\mathord{\vbox{\baselineskip3.2pt\lineskiplimit0pt\kern4pt\hbox{.}\hbox{.}\hbox{.}}}}%
  \def\ddots{\mathinner{\mkern1mu\raise6.5pt\hbox{.}\mkern2mu\raise3.6pt\hbox{.}\mkern2mu\raise0.7pt\hbox{.}\mkern1mu}}%
  \def\triangle{\mathord{\scalebox{0.9}{$\symup{\Delta}$}}}%
  \def\nabla{\mathord{\raisebox{\depth}{\scalebox{1}[-1]{$\symup{\Delta}$}}}}%
  \def\checkmark{\popcheck{popdark}}%
  \def\star{\mathbin{\ast}}\def\bigstar{\mathord{\ast}}%
  \def\diamond{\mathbin{\scalebox{0.6}{$\Diamond$}}}%
  \def\bigcup{\mathop{\mathchoice{\raisebox{-0.3em}{\scalebox{1.7}{$\cup$}}}{\scalebox{1.25}{$\cup$}}{\cup}{\cup}}\displaylimits}%
  \def\bigcap{\mathop{\mathchoice{\raisebox{-0.3em}{\scalebox{1.7}{$\cap$}}}{\scalebox{1.25}{$\cap$}}{\cap}{\cap}}\displaylimits}%
  \def\lesseqqgtr{\mathrel{\lessgtr}}\def\bigoplus{\mathop{\oplus}\displaylimits}%
}
\providecommand{\ssi}{si et seulement si} % abréviation fréquente
\AtBeginDocument{\providecommand{\wideparen}{\overparen}} % arc de cercle

%% FIGFIX-BEGIN v5 — correctifs globaux des figures (docs/onbuch-generation/tools/figfix)
\usepackage{adjustbox}\usepackage{etoolbox}\tcbuselibrary{raster}
% toute figure TikZ/pgfplots plus large que la ligne est réduite à la largeur disponible
\BeforeBeginEnvironment{tikzpicture}{\begin{adjustbox}{max width=\linewidth}}
\AfterEndEnvironment{tikzpicture}{\end{adjustbox}}
% légendes pgfplots lisibles par-dessus les courbes
\pgfplotsset{every axis/.append style={legend style={fill=white,fill opacity=0.92,text opacity=1,draw=black!15,inner sep=3pt}}}
% étiquettes d'arêtes (syntaxe edge["..."]) avec fond blanc
\tikzset{every edge quotes/.append style={fill=white,inner sep=1.2pt}}
%% FIGFIX-END
\begin{document}

\pagegarde{Générateur et récepteur électriques}{Physique}{1ère C}{Module 4 : Aspects énergétiques des circuits électriques}{N14}{5 h}{Cette leçon modélise le générateur électrique (f.é.m. E, résistance interne r) et le récepteur (f.c.é.m. E', résistance interne r') par leurs caractéristiques tension-intensité. Elle fournit les outils pour tracer, exploiter et exploiter ces caractéristiques, y compris pour des associations de générateurs, en vue des épreuves du Probatoire.
\vspace{4pt}
\begin{multicols}{2}\small
\begin{itemize}
\item À la fin de cette leçon, je sais définir un générateur et un récepteur et préciser leur rôle dans un circuit.
\item Schématiser un générateur et un récepteur en convention correcte.
\item Établir et écrire les relations U = E − rI (générateur) et U = E' + r'I (récepteur).
\item Réaliser le montage permettant de tracer la caractéristique d'un générateur et celle d'un récepteur.
\item Tracer et exploiter une caractéristique pour déterminer graphiquement E, r (ou E', r').
\item Déterminer les caractéristiques équivalentes d'une association de générateurs en série et en parallèle.
\end{itemize}\end{multicols}}

\begin{sommaire}
\begin{multicols}{2}
\begin{enumerate}
\item Le générateur électrique : définition, rôle et modélisation
\item Caractéristique U = f(I) d'un générateur
\item Association de générateurs
\item Le récepteur électrique : définition, types et rôle
\item Caractéristique U = f(I) d'un récepteur actif
\item Point de fonctionnement d'un circuit générateur–récepteur
\item Activité d'intégration
\item Méthodes et exercices résolus
\item Exercices
\item Corrigés des exercices
\item Fiche bilan
\end{enumerate}
\end{multicols}
\end{sommaire}

\begin{objectifs}
\begin{itemize}
\item À la fin de cette leçon, je sais définir un générateur et un récepteur et préciser leur rôle dans un circuit.
\item Schématiser un générateur et un récepteur en convention correcte.
\item Établir et écrire les relations U = E − rI (générateur) et U = E' + r'I (récepteur).
\item Réaliser le montage permettant de tracer la caractéristique d'un générateur et celle d'un récepteur.
\item Tracer et exploiter une caractéristique pour déterminer graphiquement E, r (ou E', r').
\item Déterminer les caractéristiques équivalentes d'une association de générateurs en série et en parallèle.
\item Distinguer les types de récepteurs (résistor, électrolyseur, moteur) et leur comportement.
\item Déterminer le point de fonctionnement d'un circuit générateur–récepteur.
\end{itemize}
\end{objectifs}

\begin{prerequis}
\begin{itemize}
\item Tension, intensité, loi d'Ohm et loi des mailles (leçons précédentes du module 4)
\item Conventions générateur et récepteur pour les flèches tension/courant
\item Lecture et exploitation d'un graphe affine (coefficient directeur, ordonnée à l'origine)
\item Utilisation du voltmètre, de l'ampèremètre et du rhéostat
\item Notion d'énergie et de puissance électrique (P = UI)
\end{itemize}
\end{prerequis}

\newpage

\coursec{Le générateur électrique : définition, rôle et modélisation}

À Yaoundé comme à Maroua, quand ENEO coupe le courant, on branche un groupe électrogène ou une batterie : la lampe brille, mais la tension aux bornes de la batterie chute dès qu'on allume plusieurs appareils. Pourquoi une pile neuve affiche 4,5 V à vide et seulement 4,2 V quand elle alimente une torche ? Et pourquoi le moteur du moulin à maïs chauffe-t-il en fonctionnement ? Comprendre ces phénomènes exige de modéliser les générateurs et les récepteurs par leurs caractéristiques : c'est l'objet de cette leçon.

\textbf{Problématique.} Comment modéliser un générateur réel pour prévoir la tension qu'il délivre en fonction du courant qu'il fournit ? C'est la question que nous allons résoudre dans cette première partie.

\courssub{Qu'est-ce qu'un générateur électrique ?}

\textbf{Observation.} Une pile plate de 4,5 V, la batterie d'une mototaxi, le panneau solaire posé sur le toit d'une maison d'électrification rurale dans l'Extrême-Nord, l'alternateur du barrage de Nachtigal : ces appareils n'ont rien en commun en apparence. Pourtant, branchés dans un circuit, ils font tous la même chose : ils \cle{mettent les charges électriques en mouvement} et maintiennent un courant dans le circuit. Pour cela, ils puisent dans une réserve d'énergie d'une autre nature : réactions chimiques dans la pile, énergie mécanique de l'eau pour l'alternateur, lumière du Soleil pour le panneau.

\begin{definition}[Générateur électrique]
Un \cle{générateur électrique} est un \textbf{dipôle} qui transforme une forme d'énergie non électrique (chimique, mécanique, lumineuse, thermique) en \textbf{énergie électrique}, et qui maintient une différence de potentiel entre ses bornes afin d'assurer la circulation d'un courant dans le circuit extérieur.
\end{definition}

Le rôle du générateur est donc double :
\begin{itemize}
\item \textbf{rôle énergétique} : il fournit l'énergie électrique que les récepteurs du circuit vont consommer ;
\item \textbf{rôle électrique} : il impose entre ses bornes une tension qui met en mouvement les porteurs de charge (les électrons dans les métaux) et entretient le courant malgré les résistances du circuit.
\end{itemize}

\begin{center}
\begin{tabularx}{\linewidth}{|l|X|X|}
\hline
\rowcolor{popblueL} \textbf{Générateur} & \textbf{Énergie source} & \textbf{Exemple d'usage au Cameroun} \\
\hline
Pile (saline, alcaline) & Énergie chimique & Télécommande, radio, torche \\
\hline
Batterie d'accumulateurs & Énergie chimique (réversible) & Démarrage des voitures et mototaxis, secours des antennes MTN/Orange \\
\hline
Dynamo, alternateur & Énergie mécanique & Éclairage de vélo, groupes électrogènes, barrages d'Edéa et de Lom Pangar \\
\hline
Panneau photovoltaïque & Énergie lumineuse (solaire) & Électrification rurale, bornes de recharge solaires \\
\hline
\end{tabularx}
\end{center}

\begin{popfigure}
\begin{tikzpicture}[font=\small]
% Pile
\draw[fill=popgoldL,draw=popdark] (0,0) rectangle (1.4,0.6);
\draw[fill=popdark] (1.4,0.2) rectangle (1.55,0.4);
\node[below] at (0.7,0) {\textbf{Pile}};
\node[below,popmuted] at (0.7,-0.35) {chimique};
% Batterie
\draw[fill=popgreenL,draw=popdark] (3,0) rectangle (4.6,0.7);
\draw[fill=popdark] (4.6,0.45) rectangle (4.75,0.6);
\draw[fill=popdark] (4.6,0.1) rectangle (4.75,0.25);
\node[below] at (3.8,0) {\textbf{Batterie}};
\node[below,popmuted] at (3.8,-0.35) {chimique};
% Panneau solaire
\draw[fill=popblueL,draw=popdark] (6.4,0) rectangle (8.2,0.8);
\draw[popline] (6.4,0.27) -- (8.2,0.27) (6.4,0.53) -- (8.2,0.53);
\draw[popline] (7,0) -- (7,0.8) (7.6,0) -- (7.6,0.8);
\draw[popgold,thick,->] (6.0,1.5) -- (6.6,0.9);
\draw[popgold,thick,->] (6.8,1.6) -- (7.3,0.95);
\node[popgold] at (6.2,1.7) {Soleil};
\node[below] at (7.3,0) {\textbf{Panneau solaire}};
\node[below,popmuted] at (7.3,-0.35) {lumineuse};
% Flèches énergie électrique
\draw[poporange,very thick,->] (0.7,0.9) -- (0.7,1.5);
\draw[poporange,very thick,->] (3.8,1.0) -- (3.8,1.5);
\draw[poporange,very thick,->] (7.3,1.1) -- (7.3,1.5);
\node[poporange,right] at (8.4,1.3) {\textbf{Énergie électrique}};
\draw[poporange,very thick] (0.7,1.5) -- (8.3,1.5);
\end{tikzpicture}
\legende{Trois générateurs, trois énergies sources différentes, un même produit : l'énergie électrique.}
\end{popfigure}

\courssub{Modélisation : le générateur idéal de tension}

Pour raisonner simplement, on commence par un modèle parfait.

\begin{definition}[Générateur idéal de tension]
Un \cle{générateur idéal de tension} est un dipôle qui maintient entre ses bornes une tension \textbf{constante}, égale à sa \textbf{force électromotrice} $E$, quelle que soit l'intensité du courant qu'il débite :
\[ U = E \quad \text{pour tout } I. \]
\end{definition}

Son symbole normalisé est un cercle traversé par la flèche de tension $E$, ou, pour les piles et batteries, le symbole à deux traits (trait long : pôle $+$ ; trait court et gras : pôle $-$).

\begin{attention}[Le générateur idéal n'existe pas]
Le générateur idéal est un \textbf{modèle limite} : aucun générateur réel ne délivre une tension rigoureusement constante. Dès qu'une pile débite du courant, sa tension chute. Utiliser le modèle idéal pour un circuit réel conduit à des erreurs systématiques : c'est précisément ce que constate l'élève qui mesure 4,5 V à vide et 4,2 V en charge.
\end{attention}

\courssub{Le générateur réel : force électromotrice $E$ et résistance interne $r$}

\textbf{Expérience décisive.} On mesure la tension aux bornes d'une pile de 4,5 V dans deux situations : circuit ouvert (la pile ne débite rien) puis circuit fermé sur une lampe. À vide, le voltmètre indique $U_0 = \SI{4,5}{V}$ ; la lampe branchée, il n'indique plus que $U = \SI{4,2}{V}$ et la pile chauffe légèrement. D'où viennent les \SI{0,3}{V} « perdus » ?

L'interprétation est physique : à l'intérieur de la pile, le courant traverse l'électrolyte et les électrodes, qui \textbf{s'opposent partiellement à son passage}. Le générateur réel possède donc une résistance propre, appelée \cle{résistance interne}, notée $r$. Le courant qui la traverse y provoque une chute de tension $rI$ et un échauffement par effet Joule ($P_J = rI^2$).

\begin{definition}[Force électromotrice et résistance interne]
\begin{itemize}
\item La \cle{force électromotrice} (f.é.m.) $E$ d'un générateur est la tension mesurée entre ses bornes lorsqu'il \textbf{ne débite aucun courant} ($I = 0$, circuit ouvert). Elle s'exprime en volts (V) et caractérise le générateur.
\item La \cle{résistance interne} $r$ est la résistance propre du générateur. Elle s'exprime en ohms ($\Omega$). Elle est responsable de la chute de tension en charge et de l'échauffement du générateur.
\end{itemize}
\end{definition}

\begin{propriete}[Modèle du générateur réel (modèle de Thévenin)]
Un générateur réel se modélise par l'association \textbf{en série} d'un générateur idéal de tension de f.é.m. $E$ et d'un résistor de résistance $r$. La tension entre ses bornes A et B, lorsqu'il débite un courant d'intensité $I$, est :
\[ U_{AB} = E - rI. \]
C'est l'\textbf{équation du générateur}. Vérification dimensionnelle : $[rI] = \Omega \times \text{A} = \text{V}$, homogène à une tension.
\end{propriete}

\textbf{Établissement par bilan énergétique.} Pour une charge $q$ qui traverse le générateur, l'énergie chimique (ou mécanique) convertie vaut $qE$. Une partie $q \cdot rI$ est dissipée par effet Joule dans la résistance interne ; le reste, $qU_{AB}$, est transmis au circuit extérieur. La conservation de l'énergie s'écrit $qE = qU_{AB} + qrI$, d'où, en divisant par $q \neq 0$ : $U_{AB} = E - rI$. En multipliant par $I$, on obtient le bilan de puissances :
\[ \underbrace{EI}_{\text{puissance convertie}} = \underbrace{U_{AB}\,I}_{\text{puissance utile fournie}} + \underbrace{rI^2}_{\text{puissance dissipée (Joule)}}. \]

\begin{popfigure}
\begin{tikzpicture}[european,font=\small]
\draw (0,0) to[V=$E$,invert] (0,2.2);
\draw (0,2.2) to[R=$r$,i={$I$}] (3.2,2.2);
\draw (3.2,2.2) -- (4.2,2.2) node[right]{A};
\draw (0,0) -- (4.2,0) node[right]{B};
\draw[poporange,very thick,->] (4.6,0.15) -- (4.6,2.05) node[midway,right]{$U_{AB}$};
\draw[dashed,popmuted] (-0.7,-0.5) rectangle (3.6,2.8);
\node[popmuted] at (1.5,-0.75) {générateur réel $(E,\,r)$};
\end{tikzpicture}
\legende{Modèle du générateur réel : source idéale $E$ en série avec la résistance interne $r$. En convention générateur, les flèches $U_{AB}$ et $I$ sont de sens opposés.}
\end{popfigure}

\courssub{La convention générateur}

\begin{definition}[Convention générateur]
Pour un générateur, on oriente le dipôle de sorte que les flèches représentant la tension $U$ et l'intensité $I$ soient de \textbf{sens opposés} : le courant sort par la borne de potentiel le plus élevé (borne $+$). On dit que le générateur est orienté en \cle{convention générateur}.
\end{definition}

Dans cette convention, le produit $P = UI$ représente la puissance électrique \textbf{fournie} par le générateur au reste du circuit (positive si le dipôle fonctionne effectivement en générateur). C'est l'inverse de la convention récepteur que nous rencontrerons plus loin : pour un récepteur, $U$ et $I$ sont fléchés dans le même sens et $P = UI$ est la puissance \textbf{reçue}.

\begin{methode}[Mesurer la f.é.m. $E$ au voltmètre]
\begin{enumerate}
\item Ouvrir le circuit : le générateur ne doit alimenter \textbf{aucun} récepteur ($I = 0$).
\item Brancher un voltmètre (calibre continu adapté, par exemple 20 V) directement aux bornes du générateur, borne V sur le pôle $+$.
\item Lire la tension : comme la résistance interne du voltmètre est très grande (de l'ordre de \SI{10}{M\Omega}), le courant qu'il prélève est quasi nul, donc la chute interne $rI$ est négligeable.
\item Conclure : la tension lue est $U_0 = E$, la f.é.m. du générateur.
\end{enumerate}
\end{methode}

\begin{attention}[Ne pas confondre $E$ et $U$]
\begin{itemize}
\item $E$ est une \textbf{caractéristique du générateur} : elle ne dépend que de lui (sa constitution chimique, sa température, son usure), pas du circuit extérieur.
\item $U = E - rI$ est la tension \textbf{réellement disponible} : elle dépend du courant débité, donc du circuit branché.
\item Dire « la tension de la pile est 4,5 V » n'est correct qu'à vide. En charge, il faut écrire $U < E$. Au Probatoire, confondre $E$ et $U$ dans l'équation du générateur est l'erreur la plus fréquente — et la plus sanctionnée.
\end{itemize}
\end{attention}

\begin{exoresolu}[Tension en charge d'une batterie]
Une batterie de voiture a pour f.é.m. $E = \SI{12,0}{V}$ et pour résistance interne $r = \SI{0,020}{\Omega}$. Elle alimente un phare qui absorbe un courant $I = \SI{5,0}{A}$. Calculer la tension $U$ à ses bornes et la puissance dissipée par effet Joule à l'intérieur de la batterie.
\tcblower
\textbf{Solution.} On applique l'équation du générateur :
\[ U = E - rI = 12,0 - 0,020 \times 5,0 = 12,0 - 0,10 = \SI{11,9}{V}. \]
La chute de tension interne vaut $rI = \SI{0,10}{V}$ : faible, car $r$ est très petite.
La puissance dissipée par effet Joule dans la batterie est :
\[ P_J = rI^2 = 0,020 \times (5,0)^2 = \SI{0,50}{W}. \]
La puissance utile fournie au phare vaut $P_u = UI = 11,9 \times 5,0 = \SI{59,5}{W}$, pour une puissance convertie $EI = \SI{60,0}{W}$ : le bilan $60,0 = 59,5 + 0,5$ est bien vérifié.
\end{exoresolu}

\begin{savaistu}[Pourquoi la batterie démarre le moteur]
Au démarrage, le démarreur d'une voiture réclame un courant énorme : plus de \SI{300}{A} ! Une batterie 12 V y parvient grâce à sa résistance interne extrêmement faible, de l'ordre de quelques milliohms ($r \approx \SI{5}{m\Omega}$). Même sous \SI{300}{A}, la chute interne ne vaut que $rI \approx \SI{1,5}{V}$. C'est aussi pourquoi une batterie usée, dont la résistance interne a augmenté, « n'arrive plus à démarrer » par matinée froide à Bafoussam : la tension s'effondre dès que le démarreur tire du courant. À l'échelle du pays, les alternateurs des barrages de Song Loulou et d'Edéa jouent le même rôle, avec des f.é.m. et des puissances des millions de fois plus grandes.
\end{savaistu}

\begin{aretenir}[L'essentiel de la section]
\begin{itemize}
\item Un générateur transforme une énergie (chimique, mécanique, lumineuse) en énergie électrique et maintient le courant dans le circuit.
\item Il se modélise par une source idéale $E$ en série avec une résistance interne $r$ : $U = E - rI$.
\item $E$ se mesure au voltmètre en circuit ouvert ; $r$ explique la chute de tension en charge et l'échauffement.
\item En convention générateur, les flèches $U$ et $I$ sont de sens opposés, et $P = UI$ est la puissance fournie au circuit.
\end{itemize}
\end{aretenir}

\coursec{Caractéristique U = f(I) d'un générateur}

Après avoir modélisé le générateur réel par un générateur idéal de tension $E$ en série avec une résistance interne $r$ (section précédente), nous devons maintenant \emph{valider ce modèle par l'expérience} et en tirer la relation fondamentale qui lie la tension $U$ aux bornes du générateur à l'intensité $I$ qu'il délivre. C'est l'objet de la \textbf{caractéristique $U = f(I)$}.

\courssub{Montage expérimental et protocole de mesure}

Pour tracer cette caractéristique, il faut faire varier l'intensité $I$ traversant le générateur et mesurer la tension $U$ à ses bornes pour chaque valeur de $I$. Le montage classique utilise un \textbf{rhéostat} (résistance variable) en série avec le générateur pour modifier le courant.

\begin{experience}[Montage pour tracer la caractéristique $U = f(I)$ d'un générateur]
\textbf{Matériel :} un générateur (pile, batterie, alimentation stabilisée), un rhéostat (ou une boîte à décennies), un ampèremètre, un voltmètre, fils de connexion.
\textbf{Schéma du montage :}
\begin{popfigure}
\begin{tikzpicture}[european, thick, scale=0.9]
% Nœuds principaux
\coordinate (N) at (0,-3); % Borne N (bas)
\coordinate (P) at (0,0);   % Borne P (haut)
\coordinate (A1) at (1.5,0);
\coordinate (B1) at (3.5,0);
\coordinate (B)  at (4,0);
\coordinate (C)  at (4,-3);
% Générateur : de N (-) vers P (+)
\draw (N) to[battery1, l={$E$}] (P);
% Ampèremètre en série sur la branche haute
\draw (P) to[ammeter, l={$A$}] (A1);
% Rhéostat
\draw (A1) to[R, l={$R_{rh}$}, -*] (B1);
\draw (B1) -- (B);
% Fil de retour
\draw (B) -- (C);
\draw (C) -- (N);
% Voltmètre en dérivation aux bornes du générateur (P et N)
\draw (P) to[voltmeter, l={$V$}, *-*] (N);
% Bornes P et N
\node[above left, popdark] at (P) {P ($+$)};
\node[below left, popdark] at (N) {N ($-$)};
% Flèches sens du courant conventionnel (sort par P)
\draw[->, poporange, thick] (0.75,0.3) -- (1.5,0.3) node[midway, above] {$I$};
\draw[->, poporange, thick] (0.3,-0.5) -- (0.3,-2.5) node[midway, left] {$I$};
\end{tikzpicture}
\legende{Montage expérimental : le rhéostat $R_{rh}$ permet de faire varier $I$ ; $A$ en série mesure $I$ ; $V$ en dérivation aux bornes P et N mesure $U$. Le générateur est orienté de N ($-$) vers P ($+$).}
\end{popfigure}

\textbf{Protocole :}
\begin{enumerate}
    \item Brancher le circuit en respectant les polarités (borne $+$ du générateur vers $+$ des appareils).
    \item Régler le rhéostat sur sa valeur \textbf{maximale} (courant minimal, tension aux bornes maximale).
    \item Relever $(U; I)$.
    \item Diminuer progressivement la résistance du rhéostat (augmenter $I$) et relever plusieurs couples $(U; I)$ (au moins 5 à 6 points bien répartis).
    \item \textbf{Attention :} ne pas descendre jusqu'au court-circuit ($R_{rh}=0$) pour ne pas endommager le générateur ni l'ampèremètre. On extrapolera la droite.
\end{enumerate}
\end{experience}

\begin{exemplebox}[Tableau de mesures typique (pile 4,5 V, $r \approx 2\,\Omega$)]
\begin{center}
\begin{tabularx}{\linewidth}{|c|*{6}{X|}}\hline
$R_{rh}\,(\Omega)$ & $\infty$ (circuit ouvert) & $10$ & $6,8$ & $4,7$ & $3,3$ & $2,2$ \\ \hline
$I\,(\text{mA})$ & $0$ & $370$ & $500$ & $670$ & $850$ & $1100$ \\ \hline
$U\,(\text{V})$ & $4,50$ & $3,75$ & $3,50$ & $3,15$ & $2,80$ & $2,30$ \\ \hline
\end{tabularx}
\end{center}
On observe bien que $U$ diminue quand $I$ augmente.
\end{exemplebox}

\courssub{Tracé de la caractéristique et établissement de la loi $U = E - rI$}

En reportant les points $(I; U)$ dans un repère $(O; \vec{\imath}, \vec{\jmath})$ avec $I$ en abscisse et $U$ en ordonnée, on obtient des points alignés sur une \textbf{droite décroissante}.

\begin{popfigure}
\begin{tikzpicture}
\begin{axis}[
    width=\linewidth, height=0.55\linewidth,
    axis lines=middle,
    xlabel=$I\,(\text{A})$, ylabel=$U\,(\text{V})$,
    xmin=-0.1, xmax=2.5,
    ymin=-0.2, ymax=5.0,
    xtick={0,0.5,1.0,1.5,2.0,2.25},
    ytick={0,1,2,3,4,4.5},
    grid=major, grid style={popline},
    clip=false,
    samples=2,
    domain=0:2.25,
]
% Droite caractéristique U = 4.5 - 2*I
\addplot[popcurve, thick, domain=0:2.25] {4.5 - 2*x};
% Point E (ordonnée à l'origine)
\node[popdark, anchor=south east] at (axis cs:0,4.5) {$E = 4,5\,\text{V}$};
\draw[dashed, popmuted] (axis cs:0,0) -- (axis cs:0,4.5);
% Point court-circuit (abscisse à l'origine)
\node[popdark, anchor=north west] at (axis cs:2.25,0) {$I_{CC} = 2,25\,\text{A}$};
\draw[dashed, popmuted] (axis cs:2.25,0) -- (axis cs:2.25,-0.2);
% Pente -r
\draw[poporange, thick, <->] (axis cs:0.5,3.5) -- node[above right, poporange] {$\Delta U = -1,0\,\text{V}$} (axis cs:1.0,3.5);
\draw[poporange, thick, <->] (axis cs:1.0,3.5) -- node[right, poporange] {$\Delta I = 0,5\,\text{A}$} (axis cs:1.0,3.0);
\node[poporange, anchor=south west] at (axis cs:1.1,3.3) {Pente $= \dfrac{\Delta U}{\Delta I} = -r = -2,0\,\Omega$};
% Points mesures
\addplot[only marks, mark=*, mark size=2.5, popblue] coordinates {
    (0,4.5) (0.37,3.75) (0.5,3.5) (0.67,3.15) (0.85,2.8) (1.1,2.3)
};
\end{axis}
\end{tikzpicture}
\legende{Caractéristique $U = f(I)$ d'un générateur réel : droite décroissante d'équation $U = E - rI$. $E$ est l'ordonnée à l'origine (circuit ouvert), $I_{CC}=E/r$ l'abscisse à l'origine (court-circuit), $r$ la valeur absolue de la pente.}
\end{popfigure}

\textbf{Établissement rigoureux par la loi des mailles (exigible au Probatoire).}
Considérons le modèle équivalent du générateur réel : un siège de force électromotrice (f.e.m.) $E$ en série avec une résistance $r$. Le circuit extérieur (rhéostat + ampèremètre) est traversé par le courant $I$ dans le sens $P \to N$ à l'extérieur du générateur (convention générateur).
Appliquons la \textbf{loi des mailles} (somme algébrique des tensions = 0) au maille unique, en parcourant le circuit dans le sens du courant $I$ (de $P$ vers $N$ à l'extérieur, puis de $N$ vers $P$ à l'intérieur) :
\begin{itemize}
    \item À l'extérieur (borne $P$ vers borne $N$) : la tension aux bornes du générateur est $U = V_P - V_N$. C'est une \emph{chute de tension} dans le sens du parcours $\to$ signe $-$ : $-U$.
    \item À l'intérieur (borne $N$ vers borne $P$) : on traverse la résistance $r$ (chute $rI$) puis le siège $E$ (hausse $E$). Somme : $-rI + E$.
\end{itemize}
La loi des mailles s'écrit :
\[
-U + (-rI + E) = 0 \quad \iff \quad U = E - rI
\]

\begin{propriete}[Loi de sortie d'un générateur réel — Démonstration exigible]
Pour un générateur réel modélisé par un générateur idéal de f.e.m. $E$ en série avec une résistance interne $r$, la tension $U$ à ses bornes (borne $P$ potentiel $+$, borne $N$ potentiel $-$) et l'intensité $I$ qu'il délivre (sortant par $P$) sont liées par :
\[
\boxed{U = E - rI}
\]
\begin{itemize}
    \item $E$ en volts ($\text{V}$) : force électromotrice (f.e.m.), tension à vide ($I=0$).
    \item $r$ en ohms ($\Omega$) : résistance interne.
    \item $I$ en ampères ($\text{A}$) : intensité du courant délivré (convention générateur : $I$ sort par $+$).
\end{itemize}
\emph{Analyse dimensionnelle :} $[E] = \text{V}$, $[rI] = \Omega \cdot \text{A} = \text{V}$. Homogène.
\end{propriete}

\courssub{Exploitation graphique : détermination de $E$ et $r$}

La caractéristique étant une droite $U = -r I + E$, on identifie :
\begin{itemize}
    \item L'\textbf{ordonnée à l'origine} ($I=0$) : $U = E$. C'est la tension à vide (circuit ouvert).
    \item Le \textbf{coefficient directeur} (pente) : $\dfrac{\Delta U}{\Delta I} = -r$. Comme la droite est décroissante, $r$ est la \textbf{valeur absolue de la pente}.
    \item L'\textbf{abscisse à l'origine} ($U=0$) : $0 = E - r I_{CC} \implies I_{CC} = \dfrac{E}{r}$. C'est l'intensité de court-circuit.
\end{itemize}

\begin{methode}[Déterminer $E$ et $r$ graphiquement à partir de la caractéristique]
\begin{enumerate}
    \item Tracer la droite caractéristique $U = f(I)$ au mieux des points expérimentaux (règle, nuage de points).
    \item Lire l'\textbf{ordonnée à l'origine} : $E = U(I=0)$ (en V).
    \item Choisir \textbf{deux points éloignés} $A(I_A; U_A)$ et $B(I_B; U_B)$ \textbf{sur la droite tracée} (pas sur les points de mesure bruts).
    \item Calculer la pente : $p = \dfrac{U_B - U_A}{I_B - I_A}$ (en V/A = $\Omega$).
    \item La résistance interne est $r = -p = |p|$ (en $\Omega$).
    \item \emph{(Vérification optionnelle) :} Lire l'abscisse à l'origine $I_{CC}$ et vérifier $r = E / I_{CC}$.
\end{enumerate}
\end{methode}

\begin{exoresolu}[Exploitation numérique : pile plate ?]
\textbf{Énoncé.} On a relevé la caractéristique d'une pile alcaline neuve (courbe bleue) et d'une pile usagée (courbe rouge). Les droites passent par les points suivants :
\begin{itemize}
    \item Pile neuve : $A_n(0; 1,55)$ et $B_n(0,50; 1,45)$.
    \item Pile usagée : $A_u(0; 1,50)$ et $B_u(0,10; 1,00)$.
\end{itemize}
Déterminer $E$ et $r$ pour chaque pile. Commenter.
\tcblower
\textbf{Solution.}
\textbf{Pile neuve :}
$E_n = 1,55\,\text{V}$ (ordonnée à l'origine).
Pente $p_n = \dfrac{1,45 - 1,55}{0,50 - 0} = \dfrac{-0,10}{0,50} = -0,20\,\Omega$.
$r_n = 0,20\,\Omega$.
$I_{CC,n} = E_n/r_n = 7,75\,\text{A}$.

\textbf{Pile usagée :}
$E_u = 1,50\,\text{V}$ (ordonnée à l'origine, peu changée).
Pente $p_u = \dfrac{1,00 - 1,50}{0,10 - 0} = \dfrac{-0,50}{0,10} = -5,0\,\Omega$.
$r_u = 5,0\,\Omega$.
$I_{CC,u} = 0,30\,\text{A}$.

\textbf{Commentaire.}
La f.e.m. $E$ (réaction chimique) baisse peu. En revanche, la résistance interne $r$ \textbf{augmente considérablement} (de $0,2\,\Omega$ à $5\,\Omega$) car l'électrolyte se dégrade et les électrodes se polarisent. La pile usagée ne peut plus fournir de courant intense ($I_{CC}$ divisé par 25) : sa tension s'effondre dès qu'on tire un peu de courant ($U = 1,5 - 5I$). C'est pourquoi une pile "plate" au voltmètre (haute impédance $\approx 10\,\text{M}\Omega$, $I\approx 0$, $U\approx 1,5\,\text{V}$) ne fait plus tourner un moteur (faible $R$, $I$ fort, $U$ chute).
\end{exoresolu}

\courssub{Cas limites : circuit ouvert et court-circuit}

Deux situations extrêmes éclairent le sens physique de $E$ et $r$.

\begin{definition}[Circuit ouvert ($I = 0$)]
Les bornes du générateur ne sont reliées à rien (ou à un voltmètre idéal de résistance infinie).
$I = 0 \implies U = E$.
La tension aux bornes \textbf{est égale à la f.e.m.} C'est la tension maximale que le générateur peut fournir. Aucune énergie n'est dissipée dans $r$ ($P_{\text{Joule}} = rI^2 = 0$).
\end{definition}

\begin{definition}[Court-circuit ($U = 0$)]
Les bornes $P$ et $N$ sont reliées directement par un fil de résistance nulle (ou un ampèremètre idéal).
$U = 0 \implies 0 = E - r I_{CC} \implies \boxed{I_{CC} = \dfrac{E}{r}}$.
L'intensité atteint sa \textbf{valeur maximale} $I_{CC}$. Toute l'énergie fournie par le siège $E$ est dissipée en chaleur \emph{à l'intérieur} du générateur : $P_{\text{dissipée}} = r I_{CC}^2 = \dfrac{E^2}{r}$.
\end{definition}

\begin{attention}[Danger du court-circuit — \textbf{Jamais faire cela !}]
\begin{itemize}
    \item \textbf{Échauffement violent :} La puissance $E^2/r$ est concentrée dans le volume minuscule de la pile/batterie. Risque d'explosion, de fuite d'électrolyte corrosif, d'incendie.
    \item \textbf{Destruction définitive :} La chimie interne est irréversiblement dégradée.
    \item \textbf{Risque pour l'utilisateur :} Projections, brûlures.
    \item \textbf{Mesure de $I_{CC}$ :} On ne mesure \textbf{JAMAIS} $I_{CC}$ en branchant un ampèremètre seul aux bornes. On déduit $I_{CC} = E/r$ à partir de mesures en charge (rhéostat).
\end{itemize}
\emph{Exemple concret :} Une batterie de voiture ($E=12\,\text{V}, r\approx 0,02\,\Omega$) en court-circuit délivre $I_{CC}=600\,\text{A}$ ! Puissance $7,2\,\text{kW}$ dans la batterie $\to$ explosion quasi certaine.
\end{attention}

\courssub{Générateur idéal de tension : modèle limite}

Si la résistance interne $r$ est \textbf{nulle} (ou négligeable devant la résistance de la charge), le générateur est dit \textbf{idéal de tension}.

\begin{propriete}[Générateur idéal de tension]
\begin{itemize}
    \item Modèle : $r = 0\,\Omega$.
    \item Loi de sortie : $\boxed{U = E = \text{constante}}$, quelle que soit l'intensité $I$ délivrée.
    \item Caractéristique $U = f(I)$ : \textbf{droite horizontale} d'ordonnée $E$.
    \item Il impose sa tension au circuit ; l'intensité est alors déterminée uniquement par le récepteur ($I = U/R_{\text{charge}}$).
\end{itemize}
\end{propriete}

\begin{attention}[Erreur fréquente : confusion $r = U/I$ vs $r = E/I_{CC}$]
\begin{itemize}
    \item \textbf{Faux :} $r = \dfrac{U}{I}$ pour un générateur en fonctionnement normal. \emph{Pourquoi ?} $U$ est la tension \emph{utile} aux bornes, $I$ le courant. $U/I$ est la résistance équivalente du \textbf{circuit extérieur} (récepteur), pas la résistance interne !
    \item \textbf{Vrai :} $r = \dfrac{E}{I_{CC}}$ (définition via court-circuit théorique) ou $r = -\dfrac{\Delta U}{\Delta I}$ (pente de la caractéristique) ou $r = \dfrac{E - U}{I}$ (chute de tension interne divisée par le courant).
\end{itemize}
\emph{Mnémotechnique :} La résistance interne $r$ se "voit" seulement quand le courant circule \emph{à l'intérieur}. Elle se mesure par la \textbf{chute de tension} $E-U$ quand $I$ circule.
\end{attention}

\begin{savaistu}[Le "test de la pile" avec la langue ? Une très mauvaise idée !]
Autrefois, on testait une pile 9 V en posant ses bornes sur la langue : un picotement signifiait "pile bonne". \textbf{Ne le faites jamais.}
\begin{itemize}
    \item La langue est un récepteur non linéaire (électrolyte + nerfs). La sensation dépend du courant $I$, pas de $E$ seul.
    \item Une pile neuve ($r$ faible) donne un courant fort $\to$ douleur, brûlure chimique possible.
    \item Une pile usagée ($r$ fort) donne un courant faible $\to$ pas de picotement, mais $E$ peut être encore 8 V.
    \item Risque court-circuit par la salive (conducteur) : échauffement de la pile dans la bouche.
\end{itemize}
Utilisez un \textbf{multimètre} (mode voltmètre, haute impédance $\approx 10\,\text{M}\Omega$) : il mesure $U \approx E$ (circuit quasi ouvert) sans danger.
\end{savaistu}

\begin{exemplebox}[Application : Alimentation d'un village isolé (Cameroun)]
Un groupe électrogène ($E = 230\,\text{V}$, $r = 0,5\,\Omega$) alimente un village via une ligne de résistance $R_{\text{ligne}} = 2\,\Omega$. La charge du village équivaut à $R_{\text{charge}} = 20\,\Omega$.
\begin{enumerate}
    \item Circuit total : $R_{\text{tot}} = r + R_{\text{ligne}} + R_{\text{charge}} = 22,5\,\Omega$.
    \item $I = \dfrac{E}{R_{\text{tot}}} = \dfrac{230}{22,5} \approx 10,2\,\text{A}$.
    \item Tension aux bornes du générateur : $U_{\text{gen}} = E - rI = 230 - 0,5 \times 10,2 = 224,9\,\text{V}$.
    \item Tension au village (bornes de $R_{\text{charge}}$) : $U_{\text{vill}} = U_{\text{gen}} - R_{\text{ligne}}I = 224,9 - 2 \times 10,2 = 204,5\,\text{V}$.
    \item Perte de tension dans la ligne : $\Delta U_{\text{ligne}} = 20,4\,\text{V}$ (9\%).
\end{enumerate}
Cela illustre pourquoi l'électrification rurale (SONATREL/ENEO) utilise le \textbf{Haute Tension} (HTA 30 kV) pour le transport : $I$ est faible, pertes $rI^2$ et chutes $rI$ réduites, puis un transformateur abaisse au 230 V au village.
\end{exemplebox}

\begin{aretenir}[Bilan : Caractéristique $U = f(I)$ d'un générateur]
\begin{itemize}
    \item \textbf{Montage :} Générateur + Rhéostat (série) + Ampèremètre (série) + Voltmètre (dérivation bornes P,N).
    \item \textbf{Loi (Loi des mailles) :} $\boxed{U = E - rI}$ (Convention générateur : $I$ sort par $+$).
    \item \textbf{Graphique :} Droite décroissante.
    \item \textbf{Paramètres :} $E = U(I=0)$ (ordonnée à l'origine) ; $r = \left|\dfrac{\Delta U}{\Delta I}\right|$ (valeur absolue de la pente) ; $I_{CC} = E/r$ (abscisse à l'origine).
    \item \textbf{Cas limites :} Circuit ouvert ($U=E, I=0$) ; Court-circuit ($U=0, I=I_{CC}$) \textbf{INTERDIT EXPÉRIMENTALEMENT}.
    \item \textbf{Générateur idéal :} $r=0 \implies U=E$ constant (droite horizontale).
\end{itemize}
\end{aretenir}

\coursec{Association de générateurs}

Dans la section précédente, nous avons établi la caractéristique $U = E - rI$ d'un générateur seul : sa tension aux bornes diminue quand le courant débité augmente, à cause de sa résistance interne $r$. Mais dans la vie quotidienne, on utilise rarement une pile seule. Ouvre le compartiment d'une torche : tu y trouveras deux, trois ou quatre piles placées bout à bout. Les batteries des installations solaires de nos villages sont couplées entre elles. Pourquoi ? Parce qu'une pile seule ne fournit qu'une tension limitée (\SI{1,5}{V} pour une pile cylindrique alcaline) et une énergie limitée. En \cle{associant} plusieurs générateurs, on peut obtenir la tension désirée ou augmenter l'autonomie du dispositif. Cette section répond à deux questions : comment associer les générateurs, et quelles sont les caractéristiques $(E_{\text{eq}}, r_{\text{eq}})$ du générateur équivalent ?

\courssub{Pourquoi associer des générateurs ?}

Une pile plate de \SI{4,5}{V} n'est autre que trois piles de \SI{1,5}{V} associées en série dans un même boîtier. Une radio portative exige souvent \SI{6}{V} : impossible avec une seule pile de \SI{1,5}{V}, il en faut quatre. Inversement, une lampe rechargeable de camping doit éclairer toute la nuit : on cherchera alors à augmenter la quantité d'énergie disponible plutôt que la tension.

\begin{definition}[Association de générateurs]
Associer des générateurs, c'est les connecter entre eux de manière à obtenir un ensemble qui se comporte comme \textbf{un seul générateur équivalent}, de force électromotrice $E_{\text{eq}}$ et de résistance interne $r_{\text{eq}}$, dont la caractéristique est :
\[ U = E_{\text{eq}} - r_{\text{eq}}\, I \]
Il existe deux associations fondamentales : l'association \cle{en série} (les générateurs se suivent sur la même branche) et l'association \cle{en parallèle} (les générateurs sont montés sur des branches distinctes entre les deux mêmes nœuds).
\end{definition}

\courssub{Association en série}

\textbf{Intuition.} Imagine deux pompes qui poussent l'eau l'une après l'autre dans la même conduite : chacune ajoute sa « pression » à celle de la précédente. De même, deux piles montées dans le même sens (le pôle $+$ de l'une relié au pôle $-$ de l'autre) ajoutent leurs forces électromotrices. En revanche, chaque pile apporte aussi sa propre résistance interne, qui s'ajoute aux autres.

\begin{propriete}[Générateurs en série — démonstration exigible au Probatoire]
Deux générateurs $(E_1, r_1)$ et $(E_2, r_2)$ montés en série (pôles de noms contraires reliés) équivalent à un générateur unique de caractéristiques :
\[ \boxed{E_{\text{eq}} = E_1 + E_2} \qquad \boxed{r_{\text{eq}} = r_1 + r_2} \]
Plus généralement, pour $n$ générateurs en série : $E_{\text{eq}} = \sum_{i=1}^n E_i$ et $r_{\text{eq}} = \sum_{i=1}^n r_i$.
\end{propriete}

\begin{experience}[Démonstration guidée]
Considérons les deux générateurs en série parcourus par le même courant $I$ (propriété des circuits en série), la tension aux bornes de l'ensemble étant $U$.
\begin{enumerate}
\item \textbf{Modélise chaque générateur} par son modèle de Thévenin : une f.é.m. idéale $E_i$ en série avec une résistance interne $r_i$.
\item \textbf{Écris la tension aux bornes de chacun} (convention générateur) :
\[ U_1 = E_1 - r_1 I \qquad U_2 = E_2 - r_2 I \]
\item \textbf{Applique l'additivité des tensions} (loi des mailles) : la tension aux bornes de l'association est la somme des tensions partielles :
\[ U = U_1 + U_2 = (E_1 + E_2) - (r_1 + r_2)\, I \]
\item \textbf{Identifie} avec la forme $U = E_{\text{eq}} - r_{\text{eq}} I$ : par unicité, $E_{\text{eq}} = E_1 + E_2$ et $r_{\text{eq}} = r_1 + r_2$. $\square$
\end{enumerate}
\textbf{Vérification dimensionnelle} : $E_1$ et $E_2$ sont des tensions (volts), leur somme est une tension ; $r_1$ et $r_2$ sont des résistances (ohms), leur somme est une résistance. La relation est homogène.
\end{experience}

\begin{propriete}[Cas de $n$ générateurs identiques en série]
Pour $n$ générateurs identiques $(E, r)$ montés en série :
\[ E_{\text{eq}} = nE \qquad r_{\text{eq}} = nr \]
\end{propriete}

\begin{attention}[Piles montées en opposition]
Si l'on monte deux piles « à l'envers » (pôles de même nom reliés), leurs f.é.m. se \textbf{retranchent} : $E_{\text{eq}} = |E_1 - E_2|$ (en prenant le sens de la plus grande), et les résistances internes s'ajoutent toujours ($r_{\text{eq}} = r_1 + r_2$). Si les deux piles sont identiques, $E_{\text{eq}} = 0$ : l'ensemble ne peut alimenter aucun circuit externe, et un courant de circulation interne use les piles inutilement (échauffement). Vérifie toujours le sens des piles dans une torche : le pôle $+$ de l'une doit toucher le pôle $-$ de la suivante.
\end{attention}

\begin{exemplebox}[Quatre piles en série pour une lampe]
Une lampe de poche est alimentée par 4 piles identiques ($E = \SI{1,5}{V}$ ; $r = \SI{0,5}{\ohm}$) montées en série. Le générateur équivalent a pour caractéristiques :
\[ E_{\text{eq}} = 4 \times \SI{1,5}{V} = \SI{6,0}{V} \qquad r_{\text{eq}} = 4 \times \SI{0,5}{\ohm} = \SI{2,0}{\ohm} \]
Si la lampe débite un courant $I = \SI{0,30}{A}$, la tension à ses bornes vaut :
\[ U = E_{\text{eq}} - r_{\text{eq}} I = \SI{6,0}{V} - \SI{2,0}{\ohm} \times \SI{0,30}{A} = \SI{5,4}{V} \]
Sans association série, une seule pile n'aurait fourni que $\SI{1,5}{V} - \SI{0,5}{\ohm} \times \SI{0,30}{A} = \SI{1,35}{V}$ : la lampe n'aurait pas éclairé.
\end{exemplebox}

\begin{popfigure}
\begin{center}
\begin{circuitikz}[european, scale=0.92, transform shape]
% --- Série concordante ---
\draw (0,0) to[battery1, l={$E_1$}] (1.6,0) to[battery1, l={$E_2$}] (3.2,0);
\draw (0,0) -- (0,-0.9) node[below]{\small A};
\draw (3.2,0) -- (3.2,-0.9) node[below]{\small B};
\node[popblue] at (1.6,0.75) {\small \textbf{Série concordante}};
\node at (1.6,-1.5) {\small $E_{\text{eq}} = E_1 + E_2$};
% --- Série en opposition ---
\draw (5.4,0) to[battery1, l={$E_1$}] (7.0,0);
\draw (7.0,0) to[battery1, l_={$E_2$}, invert] (8.6,0);
\draw (5.4,0) -- (5.4,-0.9) node[below]{\small A};
\draw (8.6,0) -- (8.6,-0.9) node[below]{\small B};
\node[poporange] at (7.0,0.75) {\small \textbf{Série en opposition}};
\node at (7.0,-1.5) {\small $E_{\text{eq}} = |E_1 - E_2|$};
% --- Parallèle ---
\draw (11.0,0.8) to[battery1, l={$E$}] (13.4,0.8);
\draw (11.0,-1.2) to[battery1, l_={$E$}] (13.4,-1.2);
\draw (11.0,0.8) -- (11.0,-1.2);
\draw (13.4,0.8) -- (13.4,-1.2);
\draw (11.0,-0.2) -- (10.4,-0.2) node[left]{\small A};
\draw (13.4,-0.2) -- (14.0,-0.2) node[right]{\small B};
\node[popgreen] at (12.2,1.6) {\small \textbf{Parallèle (identiques)}};
\node at (12.2,-2.1) {\small $E_{\text{eq}} = E$};
\end{circuitikz}
\end{center}
\legende{Associations de deux piles : en série concordante (les f.é.m. s'ajoutent), en série en opposition (les f.é.m. se retranchent), en parallèle (même f.é.m., courant partagé).}
\end{popfigure}

\courssub{Association en parallèle de générateurs identiques}

\textbf{Intuition.} Montées en parallèle, deux piles identiques se partagent le courant : chacune ne débite que la moitié du courant total. Elles « fatiguent » donc deux fois moins vite : l'autonomie double, et le courant maximal que l'ensemble peut fournir sans chute de tension excessive augmente.

\begin{propriete}[Générateurs identiques en parallèle — résultat admis]
Pour $n$ générateurs \textbf{identiques} $(E, r)$ montés en parallèle (pôles de même nom reliés ensemble) :
\[ \boxed{E_{\text{eq}} = E} \qquad \boxed{r_{\text{eq}} = \dfrac{r}{n}} \]
\end{propriete}

\textbf{Justification qualitative.} Tous les générateurs ont la même tension à leurs bornes (loi des nœuds), donc la f.é.m. de l'ensemble reste $E$ : le parallèle n'augmente pas la tension. En revanche, le courant total $I$ se répartit équitablement (symétrie) : chaque générateur ne débite que $I/n$. La chute de tension interne de chacun vaut alors $r \times I/n = (r/n) I$, ce qui montre que l'ensemble se comporte comme un générateur de résistance interne $r/n$. On retrouve d'ailleurs la loi d'association des résistances en parallèle ($1/r_{\text{eq}} = n/r$).

\begin{attention}[Condition impérative pour le montage en parallèle]
Le montage en parallèle n'est autorisé qu'avec des générateurs \textbf{identiques} (même $E$, même $r$) et de \textbf{même polarité}. Si deux piles de f.é.m. différentes sont reliées en parallèle, la pile de plus forte f.é.m. impose son courant dans l'autre, même circuit ouvert : un courant interne permanent circule, qui use les piles, les échauffe et peut les détériorer (risque de fuite ou d'explosion pour les accumulateurs). De même, une pile montée à l'envers est traversée par un courant destructeur. En pratique, on n'associe en parallèle que des éléments neufs, de même marque, de même technologie et de même état de charge.
\end{attention}

\begin{methode}[Tracer la caractéristique $U = f(I)$ d'une association de générateurs]
\begin{enumerate}
\item \textbf{Monte le circuit} : association des générateurs, rhéostat (résistance variable) en série, ampèremètre en série, voltmètre en parallèle sur les bornes de l'association.
\item \textbf{Varié le rhéostat} pour obtenir plusieurs couples de valeurs $(I, U)$ (au moins 5 points, de $I=0$ à $I_{\text{max}}$).
\item \textbf{Reporte les points} dans un repère $(I ; U)$.
\item \textbf{Trace la droite} de régression (ou la droite passant par les points extrêmes si alignement parfait).
\item \textbf{Déduis $E_{\text{eq}}$ et $r_{\text{eq}}$} : $E_{\text{eq}} = U(I=0)$ (ordonnée à l'origine) ; $r_{\text{eq}} = -\text{pente} = -\dfrac{\Delta U}{\Delta I}$.
\end{enumerate}
\end{methode}

\begin{exoresolu}[Série ou parallèle pour une radio ?]
Une radio portative fonctionne sous \SI{6}{V} et consomme un courant de \SI{0,20}{A}. On dispose de piles ($E = \SI{1,5}{V}$ ; $r = \SI{0,5}{\ohm}$).
\begin{enumerate}
\item Combien de piles faut-il et comment les monter ?
\item Calculer la tension réellement disponible aux bornes de la radio.
\item Que se passerait-il avec 4 piles en parallèle ?
\end{enumerate}
\tcblower
\textbf{1. Choix de l'association.} Il faut une tension de \SI{6}{V} alors que chaque pile fournit \SI{1,5}{V} : seule l'association série augmente la tension. Il faut $n = \dfrac{6}{1,5} = 4$ piles en série. Le générateur équivalent vaut :
\[ E_{\text{eq}} = 4 \times \SI{1,5}{V} = \SI{6,0}{V} \qquad r_{\text{eq}} = 4 \times \SI{0,5}{\ohm} = \SI{2,0}{\ohm} \]
\textbf{2. Tension disponible.} Avec $I = \SI{0,20}{A}$ :
\[ U = E_{\text{eq}} - r_{\text{eq}} I = \SI{6,0}{V} - \SI{2,0}{\ohm} \times \SI{0,20}{A} = \SI{5,6}{V} \]
La tension est légèrement inférieure à \SI{6}{V} à cause des résistances internes, mais reste acceptable pour la radio.
\textbf{3. Cas du parallèle.} Quatre piles en parallèle donneraient $E_{\text{eq}} = \SI{1,5}{V}$ et $r_{\text{eq}} = \SI{0,5}{\ohm}/4 = \SI{0,125}{\ohm}$. La tension de \SI{1,5}{V} est très insuffisante : la radio ne fonctionnerait pas. Le parallèle augmente l'autonomie, pas la tension : ce n'est pas la bonne solution ici.
\end{exoresolu}

\courssub{Bilan : quel montage choisir ?}

\begin{aretenir}[Choisir entre série et parallèle]
\begin{itemize}
\item \textbf{Série} : on augmente la \cle{tension} disponible ($E_{\text{eq}} = \sum E_i$), au prix d'une résistance interne plus grande ($r_{\text{eq}} = \sum r_i$). C'est le montage des torches, radios et télécommandes.
\item \textbf{Parallèle} (générateurs identiques uniquement) : la tension reste $E$, mais la résistance interne chute ($r_{\text{eq}} = r/n$) : on augmente l'\cle{autonomie} et le \cle{courant maximal} débité sans chute de tension excessive.
\end{itemize}
\end{aretenir}

\begin{center}
\begin{tabularx}{\linewidth}{|l|X|X|}
\hline
\rowcolor{popblueL} \textbf{Critère} & \textbf{Association série} & \textbf{Association parallèle} \\
\hline
F.é.m. équivalente & $E_{\text{eq}} = \sum E_i$ ($nE$ si identiques) & $E_{\text{eq}} = E$ \\
\hline
Résistance interne & $r_{\text{eq}} = \sum r_i$ ($nr$ si identiques) & $r_{\text{eq}} = \dfrac{r}{n}$ \\
\hline
Condition & Sens concordant souhaité & Générateurs identiques, même polarité \\
\hline
Avantage principal & Tension plus élevée & Autonomie et courant maximal accrus \\
\hline
Exemple d'usage & Torche 4 piles (\SI{6}{V}), batterie \SI{24}{V} & Lampe rechargeable longue durée \\
\hline
\end{tabularx}
\end{center}

\begin{savaistu}[Batteries couplées au Cameroun]
Les installations solaires photovoltaïques qui électrifient les villages et les centres de santé hors réseau (programme \emph{Électrification Rurale}) stockent l'énergie dans des batteries de \SI{12}{V}. Pour alimenter un onduleur de \SI{24}{V}, on monte deux batteries \textbf{en série} ; pour doubler l'autonomie à tension constante, on les monte \textbf{en parallèle}. De même, la batterie d'une moto-taxi (\SI{12}{V}) est constituée de six éléments de \SI{2}{V} associés en série dans un même boîtier. Les techniciens de SONATREL/ENEO ou des installateurs privés veillent toujours à ne coupler en parallèle que des batteries de même âge et de même capacité, sous peine de courants internes qui dégradent l'ensemble prématurément.
\end{savaistu}

\begin{exercice}[Pour t'entraîner]
On dispose de trois accumulateurs identiques ($E = \SI{2,0}{V}$ ; $r = \SI{0,10}{\ohm}$).
\begin{enumerate}
\item Calculer $E_{\text{eq}}$ et $r_{\text{eq}}$ lorsqu'ils sont montés en série, puis en parallèle.
\item Dans chaque cas, calculer la tension aux bornes de l'association lorsqu'elle débite un courant de \SI{1,0}{A}.
\item Un appareil exige \SI{6}{V} sous \SI{1}{A} : quel montage choisir ? La tension obtenue convient-elle ?
\end{enumerate}
\end{exercice}

\begin{corrige}[Exercice 1]
\textbf{1.} En série : $E_{\text{eq}} = 3 \times \SI{2,0}{V} = \SI{6,0}{V}$ et $r_{\text{eq}} = 3 \times \SI{0,10}{\ohm} = \SI{0,30}{\ohm}$. En parallèle : $E_{\text{eq}} = \SI{2,0}{V}$ et $r_{\text{eq}} = \SI{0,10}{\ohm}/3 \approx \SI{0,033}{\ohm}$.
\textbf{2.} Série : $U = \SI{6,0}{V} - \SI{0,30}{\ohm} \times \SI{1,0}{A} = \SI{5,7}{V}$. Parallèle : $U = \SI{2,0}{V} - \SI{0,033}{\ohm} \times \SI{1,0}{A} \approx \SI{1,97}{V}$.
\textbf{3.} Il faut le montage en série, seul capable d'approcher \SI{6}{V}. La tension réelle de \SI{5,7}{V} est légèrement inférieure à \SI{6}{V} à cause des résistances internes ; elle convient généralement, l'écart ne dépassant pas 5 \%. 
\end{corrige}

\coursec{Le récepteur électrique : définition, types et rôle}

Dans les sections précédentes, nous avons étudié le générateur, source de l'énergie électrique, puis appris à associer plusieurs générateurs pour obtenir la f.é.m. et la résistance interne désirées. Mais une question demeure : à quoi sert cette énergie une fois produite ? Elle est destinée à être \emph{utilisée}. Le ventilateur qui rafraîchit la classe à Yaoundé, le fer à repasser de la maison, le moteur de la mototaxi électrique, l'électrolyseur d'une usine de production d'aluminium à Edéa : tous ces appareils reçoivent de l'énergie électrique et la transforment. Ce sont des \cle{récepteurs}. Cette section leur est consacrée.

\courssub{Définition et rôle du récepteur dans un circuit}

\begin{experience}[Observation : que devient l'énergie électrique ?]
Branchons successivement trois appareils sur un même générateur : une lampe à incandescence, un petit moteur électrique entraînant une poulie qui soulève une masse, et une cuve à électrolyse contenant une solution de sulfate de cuivre. Observons :
\begin{itemize}
\item la lampe brille et chauffe : l'énergie électrique devient énergie thermique (chaleur) et rayonnement lumineux ;
\item le moteur soulève la masse : l'énergie électrique devient énergie \emph{mécanique} (et un peu de chaleur, le moteur est tiède) ;
\item dans la cuve, du cuivre se dépose sur une électrode : l'énergie électrique devient énergie \emph{chimique} (et un peu de chaleur).
\end{itemize}
Dans les trois cas, l'appareil \emph{consomme} l'énergie électrique fournie par le générateur et la convertit en d'autres formes d'énergie.
\end{experience}

\begin{definition}[Récepteur électrique]
Un \cle{récepteur} est un dipôle qui convertit l'énergie électrique qu'il reçoit en une ou plusieurs autres formes d'énergie (mécanique, chimique, thermique, rayonnement, etc.).
\end{definition}

\begin{aretenir}[Rôle du récepteur dans le circuit]
Dans un circuit électrique, le récepteur est l'\emph{utilisateur} de l'énergie : c'est pour lui que le générateur travaille. Le générateur fournit l'énergie électrique ; le récepteur la reçoit et la convertit. Pendant une durée $\Delta t$, un récepteur traversé par un courant d'intensité $I$ sous une tension $U$ reçoit l'énergie électrique :
\[ W_{\text{élec}} = U\,I\,\Delta t \]
Cette énergie est ensuite dégradée ou convertie selon la nature du récepteur.
\end{aretenir}

\courssub{Le récepteur passif : le résistor}

\begin{definition}[Récepteur passif]
Un récepteur est dit \cle{passif} lorsque l'énergie électrique qu'il reçoit est convertie uniquement en énergie thermique (effet Joule) et, le cas échéant, en rayonnement électromagnétique (lumière pour une lampe à incandescence), sans stockage ni conversion en énergie mécanique ou chimique macroscopique utile. Le récepteur passif par excellence est le \cle{résistor} (conducteur ohmique).
\end{definition}

Pour un résistor de résistance $R$, la loi d'Ohm s'écrit :
\[ U = R\,I \]
et l'énergie reçue pendant $\Delta t$ est entièrement dissipée en chaleur :
\[ W_J = R\,I^2\,\Delta t = \frac{U^2}{R}\,\Delta t \]

\begin{exemplebox}[Le radiateur électrique]
Un radiateur électrique de résistance $R = \SI{48.4}{\ohm}$ est branché sur le secteur ($U = \SI{220}{V}$). L'intensité qui le traverse est :
\[ I = \frac{U}{R} = \frac{220}{48.4} \approx \SI{4.5}{A} \]
La puissance convertie en chaleur vaut $P = UI = 220 \times 4,5 \approx \SI{1000}{W}$, soit 1 kW. Toute l'énergie électrique devient de l'énergie thermique : c'est un récepteur passif.
\end{exemplebox}

\courssub{Les récepteurs actifs : électrolyseur et moteur}

\begin{definition}[Récepteur actif]
Un récepteur est dit \cle{actif} lorsqu'une partie seulement de l'énergie électrique reçue est convertie en chaleur ; le reste est transformé en une autre forme d'énergie \emph{utile} :
\begin{itemize}
\item l'\cle{électrolyseur} convertit l'énergie électrique en énergie \emph{chimique} (réactions d'oxydo-réduction aux électrodes) et en chaleur ;
\item le \cle{moteur électrique} convertit l'énergie électrique en énergie \emph{mécanique} (rotation d'un arbre) et en chaleur.
\end{itemize}
\end{definition}

Pourquoi un moteur ne se comporte-t-il pas comme un simple résistor ? Parce que, lorsqu'il tourne, il développe à ses bornes une tension qui \emph{s'oppose} au passage du courant : c'est la \emph{force contre-électromotrice}.

\begin{definition}[Force contre-électromotrice et résistance interne]
Tout récepteur actif est caractérisé par deux grandeurs :
\begin{itemize}
\item sa \cle{force contre-électromotrice} (f.c.é.m.) $E'$, exprimée en volts (V) : c'est la part de la tension $U$ correspondant à la conversion de l'énergie électrique en énergie \emph{utile} (chimique ou mécanique) ;
\item sa \cle{résistance interne} $r'$, exprimée en ohms ($\Omega$) : elle rend compte de l'échauffement du récepteur (effet Joule interne).
\end{itemize}
La tension aux bornes d'un récepteur actif traversé par un courant $I$ s'écrit :
\[ U = E' + r'I \]
\end{definition}

\begin{attention}[Convention récepteur]
Pour un récepteur, on adopte la \cle{convention récepteur} : les flèches représentant la tension $U$ et le courant $I$ sont orientées \emph{dans le même sens} (le courant entre par la borne où arrive la flèche de $U$). C'est l'inverse de la convention générateur. Dans la relation $U = E' + r'I$, le signe $+$ devant $r'I$ est une conséquence directe de ce choix : le récepteur \emph{consomme} de la tension, il n'en fournit pas.
\end{attention}

\begin{popfigure}
\begin{center}
\begin{circuitikz}[european, scale=1.1]
\draw[thick, popblue] (0,0) to[R, l_={$r'$}, v={$u_{r'}$}] (3,0);
\draw[thick, popblue] (3,0) to[battery1, l_={$E'$}, *-*] (5.5,0);
\draw[->, thick, poporange] (-1.2,0.5) -- (0.2,0.5) node[midway, above]{$I$};
\draw[->, thick, popdark] (-1.2,-0.9) -- (5.5,-0.9) node[midway, below]{$U$};
\filldraw[popdark] (0,0) circle (1.5pt) node[above left]{A};
\filldraw[popdark] (5.5,0) circle (1.5pt) node[above right]{B};
\draw (-1.2,0) -- (0,0);
\end{circuitikz}
\end{center}
\legende{Modélisation d'un récepteur actif $(E', r')$ en convention récepteur : les flèches $U$ et $I$ ont le même sens. Le courant entre par la borne A, traverse la résistance interne $r'$ (chute de tension $r'I$) puis la f.c.é.m. $E'$ (borne + à gauche, borne - à droite).}
\end{popfigure}

\begin{propriete}[Bilan énergétique d'un récepteur actif]
En multipliant la relation $U = E' + r'I$ par $I\,\Delta t$, on obtient le bilan d'énergie pendant la durée $\Delta t$ :
\[ \underbrace{U\,I\,\Delta t}_{\text{énergie électrique reçue}} = \underbrace{E'\,I\,\Delta t}_{\text{énergie utile (chimique ou mécanique)}} + \underbrace{r'\,I^2\,\Delta t}_{\text{chaleur (effet Joule)}} \]
En divisant par $\Delta t$, on obtient le bilan de puissance : $P_{\text{reçue}} = P_{\text{utile}} + P_{\text{Joule}}$, soit $UI = E'I + r'I^2$.
\end{propriete}

\begin{popfigure}
\begin{center}
\begin{tikzpicture}[font=\small]
\node[draw=popblue, fill=popblueL, rounded corners, minimum width=4.6cm, minimum height=1.1cm, align=center] (rec) at (0,0) {\textbf{Énergie électrique reçue}\\ $W = U\,I\,\Delta t$};
\node[draw=popgreen, fill=popgreenL, rounded corners, minimum width=4.2cm, minimum height=1.1cm, align=center] (utile) at (-3.4,-2.6) {\textbf{Énergie utile}\\ $E'\,I\,\Delta t$\\ (chimique ou mécanique)};
\node[draw=poporange, fill=poporangeL, rounded corners, minimum width=4.2cm, minimum height=1.1cm, align=center] (joule) at (3.4,-2.6) {\textbf{Chaleur}\\ $r'\,I^2\,\Delta t$\\ (effet Joule)};
\draw[->, very thick, popgreen] (rec.south west) ++(0.4,0.1) -- (utile.north);
\draw[->, very thick, poporange] (rec.south east) ++(-0.4,0.1) -- (joule.north);
\node[draw=popdark, fill=popcream, rounded corners, align=center] (moteur) at (0,-5) {Récepteur actif : moteur ou électrolyseur $(E', r')$};
\draw[->, thick, popdark] (utile.south) |- (moteur.west);
\draw[->, thick, popdark] (joule.south) |- (moteur.east);
\end{tikzpicture}
\end{center}
\legende{Schéma-bilan énergétique d'un récepteur actif : l'énergie électrique reçue se partage entre une énergie utile (conversion non thermique, liée à $E'$) et de la chaleur dissipée par effet Joule dans la résistance interne $r'$.}
\end{popfigure}

\courssub{Condition de fonctionnement d'un récepteur actif}

\begin{propriete}[Condition de passage du courant]
Un récepteur actif ne laisse passer le courant dans le sens utile (sens de la conversion désirée) que si la tension appliquée à ses bornes est \emph{supérieure} à sa f.c.é.m. :
\[ U > E' \]
L'intensité du courant est alors :
\[ I = \frac{U - E'}{r'} \]
Si $U < E'$, le courant ne passe pas dans le sens utile : l'électrolyse ne démarre pas, le moteur ne tourne pas (il peut même freiner).
\end{propriete}

Vérifions l'homogénéité : $U$ et $E'$ sont en volts, $r'I$ en $\Omega \times \text{A} = \text{V}$. La relation est bien homogène.

\begin{exoresolu}[Moteur alimenté sous 10 V]
\textbf{Énoncé.} Un moteur électrique de f.c.é.m. $E' = \SI{6}{V}$ et de résistance interne $r' = \SI{2}{\ohm}$ est alimenté sous une tension $U = \SI{10}{V}$.
\begin{enumerate}
\item Vérifier que le moteur peut fonctionner et calculer l'intensité $I$ du courant.
\item Calculer la puissance électrique reçue, la puissance utile et la puissance dissipée par effet Joule.
\item Calculer le rendement du moteur.
\end{enumerate}
\tcblower
\textbf{Solution.}
\begin{enumerate}
\item On a $U = \SI{10}{V} > E' = \SI{6}{V}$ : la condition de fonctionnement est remplie. L'intensité vaut :
\[ I = \frac{U - E'}{r'} = \frac{10 - 6}{2} = \SI{2}{A} \]
\item Puissances mises en jeu :
\begin{align*}
P_{\text{reçue}} &= U\,I = 10 \times 2 = \SI{20}{W} \\
P_{\text{utile}} &= E'\,I = 6 \times 2 = \SI{12}{W} \\
P_{\text{Joule}} &= r'\,I^2 = 2 \times 2^2 = \SI{8}{W}
\end{align*}
On vérifie le bilan : $12 + 8 = \SI{20}{W}$. La conservation de l'énergie est satisfaite.
\item Le rendement est le rapport de la puissance utile à la puissance reçue :
\[ \eta = \frac{P_{\text{utile}}}{P_{\text{reçue}}} = \frac{E'}{U} = \frac{6}{10} = 0,60 \]
Le moteur convertit 60 \% de l'énergie électrique reçue en énergie mécanique ; les 40 \% restants chauffent le moteur.
\end{enumerate}
\end{exoresolu}

\begin{attention}[Le moteur bloqué : un danger réel]
Lorsqu'un moteur tourne, sa rotation engendre la f.c.é.m. $E'$ qui \emph{limite} le courant : $I = \dfrac{U - E'}{r'}$. Mais si l'on \emph{bloque} l'arbre du moteur (charge trop lourde, obstacle mécanique), la f.c.é.m. s'annule : $E' = 0$. Le moteur se comporte alors comme une simple résistance $r'$, et le courant devient :
\[ I_{\text{bloqué}} = \frac{U}{r'} \]
Pour le moteur de l'exemple précédent : $I_{\text{bloqué}} = \dfrac{10}{2} = \SI{5}{A}$, soit 2,5 fois le courant normal ! La puissance Joule passe de 8 W à $r'I^2 = 2 \times 25 = \SI{50}{W}$ : le moteur chauffe fortement et peut être détruit. C'est pourquoi les moteurs sont protégés par des fusibles ou des disjoncteurs thermiques.
\end{attention}

\begin{savaistu}[Pourquoi un moteur qui force consomme plus de courant]
C'est la f.c.é.m. qui explique un phénomène que tout mécanicien connaît : un moteur qui \emph{force} consomme davantage. Quand le moulin à maïs du quartier est fortement chargé, son moteur ralentit ; or la f.c.é.m. $E'$ est proportionnelle à la vitesse de rotation. Si le moteur ralentit, $E'$ diminue, donc $I = \dfrac{U - E'}{r'}$ augmente : le moteur tire plus de courant sur le réseau. C'est pour cela que, dans les quartiers de Douala ou de Yaoundé, on observe parfois une baisse de la luminosité des lampes au moment où un gros moteur démarre ou force : l'appel de courant fait chuter la tension disponible pour les autres usagers. À l'inverse, à vide, le moteur tourne vite, $E'$ est proche de $U$, et le courant consommé est faible.
\end{savaistu}

\begin{methode}[Exploiter la relation $U = E' + r'I$]
Pour traiter un exercice sur un récepteur actif :
\begin{enumerate}
\item Identifier $E'$, $r'$ et $U$ dans l'énoncé, en respectant la convention récepteur.
\item Vérifier la condition de fonctionnement $U > E'$.
\item Calculer l'intensité : $I = \dfrac{U - E'}{r'}$.
\item Effectuer le bilan de puissance : $UI = E'I + r'I^2$, et vérifier numériquement l'égalité.
\item Conclure par le rendement si demandé : $\eta = \dfrac{E'}{U}$ (pour un récepteur actif seul).
\end{enumerate}
\end{methode}

\begin{aretenir}[L'essentiel de la section]
\begin{itemize}
\item Un récepteur convertit l'énergie électrique reçue en d'autres formes d'énergie ; c'est l'utilisateur de l'énergie dans le circuit.
\item Récepteur \emph{passif} (résistor, lampe) : l'énergie devient chaleur (et lumière), $U = RI$ (linéaire ou non).
\item Récepteur \emph{actif} (moteur, électrolyseur) : modélisé par $(E', r')$, avec $U = E' + r'I$ en convention récepteur.
\item Condition de fonctionnement : $U > E'$, et alors $I = \dfrac{U - E'}{r'}$.
\item Bilan de puissance : $\underbrace{UI}_{\text{reçue}} = \underbrace{E'I}_{\text{utile}} + \underbrace{r'I^2}_{\text{Joule}}$.
\item Moteur bloqué : $E' = 0$, le courant devient $U/r'$, danger de destruction.
\end{itemize}
\end{aretenir}

Maintenant que le récepteur actif est modélisé par la relation $U = E' + r'I$, il nous reste à vérifier expérimentalement cette loi : c'est l'objet de la section suivante, où nous tracerons la caractéristique $U = f(I)$ d'un récepteur et apprendrons à en déduire $E'$ et $r'$.

\coursec{Caractéristique U = f(I) d'un récepteur actif}

Après avoir identifié les récepteurs électriques et distingué les récepteurs passifs (résistances pures) des récepteurs actifs (moteurs, électrolyseurs, batteries en charge), nous allons maintenant modéliser le comportement d'un \textbf{récepteur actif}. Tout comme pour le générateur, nous chercherons la relation mathématique qui lie la tension $U$ à ses bornes au courant $I$ qui le traverse. Cette relation, appelée \textbf{caractéristique courant-tension}, nous permettra de déterminer deux paramètres essentiels : la force électromotrice de contre-électromotrice $E'$ et la résistance interne $r'$.

\courssub{Montage expérimental et acquisition des données}

Pour tracer la caractéristique $U = f(I)$ d'un récepteur actif, nous devons faire varier le courant $I$ qui le traverse et mesurer la tension $U$ correspondante à ses bornes. Le montage classique utilise un générateur de tension variable (ou une pile avec un rhéostat en série) pour imposer le courant.

\begin{experience}[Tracé de la caractéristique d'un récepteur actif]
\textbf{Matériel :} Un générateur de tension continue variable (ou une pile \SI{9}{\volt} + rhéostat \SI{50}{\ohm}), un récepteur actif (moteur à courant continu, électrolyseur cuivre-cuivre, ou batterie rechargeable en charge), un ampèremètre (calibre \SI{1}{\ampere} ou \SI{2}{\ampere}), un voltmètre (calibre \SI{10}{\volt} ou \SI{20}{\volt}), fils de connexion.

\textbf{Montage :} L'ampèremètre est placé \textbf{en série} avec le récepteur pour mesurer l'intensité $I$. Le voltmètre est placé \textbf{en parallèle} aux bornes du récepteur pour mesurer la tension $U$. Le générateur variable permet de modifier $I$.

\begin{popfigure}
\begin{center}
\begin{tikzpicture}[european, thick, scale=0.9]
% Generateur variable
\draw (0,0) to[sV=$E_G$, *-*] (0,3);
\draw (0,3) -- (2,3);
% Rheostat
\draw (2,3) to[R=$R_{var}$, l_=Rhéostat] (4,3);
\draw (4,3) -- (5,3);
% Ammetre
\draw (5,3) to[ammeter, l_={$A$}, i={$I$}] (7,3);
\draw (7,3) -- (8,3);
% Recepteur Actif (modele Thevenin : E' + r')
\draw (8,3) to[R=$r'$, *-] (9,3);
\draw (9,3) -- (9.5,3);
\draw (9.5,3) to[battery1=$E'$, invert] (9.5,1.5); % Contre-EMF
\draw (9.5,1.5) -- (9,1.5);
\draw (9,1.5) to[short, -*] (8,1.5);
\draw (8,1.5) -- (8,0);
% Voltmeter aux bornes du recepteur (E' + r')
\draw (8,3) to[voltmeter, l_={$V$}, v={$U$}] (8,0);
\draw (8,0) -- (0,0);
% Current arrow
\draw[->, popblue, thick] (6,3.5) -- (6.5,3.5) node[above] {$I$};
\end{tikzpicture}
\end{center}
\legende{Montage pour tracer la caractéristique $U=f(I)$ d'un récepteur actif. Le voltmètre mesure la tension aux bornes de l'ensemble $E'+r'$.}
\end{popfigure}

\textbf{Protocole :}
\begin{enumerate}
\item Brancher le circuit. Régler le rhéostat (ou la tension du générateur) pour obtenir le courant maximal autorisé par le récepteur (sans l'endommager). Noter $I$ et $U$.
\item Diminuer progressivement le courant en augmentant la résistance du rhéostat (ou en baissant $E_G$). Relever au moins 6 à 8 couples de valeurs $(I, U)$ sur toute la plage de fonctionnement, y compris pour $I \approx 0$.
\item Inverser les bornes du générateur (ou du récepteur) si l'on souhaite explorer la zone de fonctionnement en générateur (optionnel pour le programme de Première).
\end{enumerate}

\textbf{Observations :} Le tableau de mesures montre que la tension $U$ \textbf{augmente} lorsque le courant $I$ augmente. Le nuage de points s'aligne sur une \textbf{droite croissante}.
\end{experience}

\begin{exemplebox}[Tableau de mesures typique — Moteur CC]
\begin{center}
\begin{tabular}{|c|c|c|c|c|c|c|c|}\hline
$I$ (\si{\ampere}) & \num{0,00} & \num{0,10} & \num{0,20} & \num{0,30} & \num{0,40} & \num{0,50} & \num{0,60} \\ \hline
$U$ (\si{\volt}) & \num{6,00} & \num{6,20} & \num{6,40} & \num{6,60} & \num{6,80} & \num{7,00} & \num{7,20} \\ \hline
\end{tabular}
\end{center}
\legende{Exemple de mesures sur un petit moteur à courant continu. On observe $U > 0$ même pour $I=0$.}
\end{exemplebox}

\courssub{Établissement de la loi $U = E' + r'I$ : Modélisation de Thévenin}

Un récepteur actif (moteur, électrolyseur, accumulateur en charge) se comporte comme une source de tension \emph{interne} qui s'oppose au courant imposé par le générateur externe. C'est la \textbf{contre-électromotrice} (c.e.m.), notée $E'$. De plus, les fils de bobinage, l'électrolyte ou les plaques de l'accumulateur possèdent une résistance propre, notée $r'$ (résistance interne du récepteur).

Le récepteur actif est donc modélisé par un \textbf{modèle de Thévenin récepteur} : un générateur de tension idéal $E'$ (la c.e.m.) \textbf{en série} avec une résistance $r'$.

\begin{popfigure}
\begin{center}
\begin{tikzpicture}[european, thick, scale=1]
% Model Thevenin Recepteur
\draw (0,0) to[R=$r'$, *-*] (2,0);
\draw (2,0) to[battery1=$E'$, invert] (2,-2); % E' oppose au courant entrant
\draw (2,-2) -- (0,-2);
\draw (0,-2) to[short, *-*] (0,0);
% Bornes A et B
\node[left] at (0,0) {A};
\node[left] at (0,-2) {B};
% Courant I entrant par A
\draw[->, popblue, thick] (-1,0) -- (0,0) node[midway, above] {$I$};
% Tension U_AB
\draw[<->, poporange, thick] (-0.5,0.3) -- (-0.5,-2.3) node[midway, left] {$U_{AB}$};
% Convention recepteur : I entre par A (+), U_AB = V_A - V_B
\node at (1, -1) [rotate=90, anchor=south] {Modèle de Thévenin};
\end{tikzpicture}
\end{center}
\legende{Modèle équivalent d'un récepteur actif : une c.e.m. $E'$ en série avec une résistance interne $r'$. Le courant $I$ entre par la borne A (potentiel le plus élevé).}
\end{popfigure}

Appliquons la \textbf{loi des mailles} (loi de Kirchhoff aux tensions) à ce modèle, en suivant le sens du courant $I$ (convention récepteur : le courant entre par la borne de potentiel le plus élevé A).

\begin{propriete}[Loi de tension aux bornes d'un récepteur actif — Démonstration exigible au Probatoire]
Soit un récepteur actif modélisé par une c.e.m. $E'$ en série avec une résistance interne $r'$. En convention récepteur ($I$ entrant par la borne $+$), la tension $U$ à ses bornes est donnée par :
\[
\boxed{U = E' + r'I}
\]
\textbf{Démonstration guidée :}
\begin{enumerate}
\item On parcourt la maille unique dans le sens du courant $I$ (de A vers B à l'intérieur du récepteur).
\item Aux bornes de la résistance $r'$ : la tension chute de $r'I$ (loi d'Ohm en convention récepteur : $U_{r'} = r'I$).
\item Aux bornes du générateur idéal $E'$ : on va du $-$ vers le $+$ (car $E'$ s'oppose au courant), la tension \textbf{augmente} de $E'$.
\item La tension totale $U = V_A - V_B$ est la somme des tensions rencontrées : $U = r'I + E'$.
\end{enumerate}
\emph{Analyse dimensionnelle :} $[E'] = \text{V}$, $[r'I] = \Omega \cdot \text{A} = \text{V}$. L'homogénéité est respectée.
\end{propriete}

\begin{attention}[Piège majeur : Confusion avec le générateur]
La caractéristique d'un \textbf{générateur} est $U = E - rI$ (droite \textbf{décroissante}).
La caractéristique d'un \textbf{récepteur actif} est $U = E' + r'I$ (droite \textbf{croissante}).
\emph{Erreur fréquente :} Écrire $U = E' - r'I$ pour un récepteur. Rappelle-toi : dans un récepteur, la tension aux bornes \textbf{doit dépasser} la c.e.m. pour faire passer le courant ($U > E'$). Le terme $r'I$ est une \textbf{chute de tension ohmique} qui s'\textbf{ajoute} à $E'$.
\end{attention}

\courssub{Exploitation graphique : Détermination de $E'$ et $r'$}

La relation $U = E' + r'I$ est l'équation d'une droite dans le repère $(I ; U)$ (abscisse $I$, ordonnée $U$). C'est une droite \textbf{croissante} (coefficient directeur $r' > 0$).

\begin{popfigure}
\begin{center}
\begin{tikzpicture}
\begin{axis}[
    width=\linewidth, height=0.6\linewidth,
    axis lines=middle,
    xlabel={$I$ (\si{\ampere})}, ylabel={$U$ (\si{\volt})},
    xmin=-0.05, xmax=0.7,
    ymin=5.5, ymax=7.5,
    xtick={0,0.2,0.4,0.6},
    ytick={6,6.5,7,7.2},
    grid=major, grid style={popmuted},
    tick label style={font=\small},
    label style={font=\normalsize},
    clip=false
]
% Droite U = 6 + 2I
\addplot[popcurve, domain=-0.05:0.65, samples=2, thick] {6 + 2*x};
% Points mesures
\addplot[only marks, mark=*, mark size=2.5, popblue] coordinates {
    (0,6) (0.1,6.2) (0.2,6.4) (0.3,6.6) (0.4,6.8) (0.5,7.0) (0.6,7.2)
};
% Droite horizontale E'
\addplot[dashed, poporange, domain=-0.05:0] {6};
\addplot[dashed, poporange, domain=0:0.65] {6};
% Fleches pour r'
\draw[popgreen, thick, <->] (axis cs:0.2, 6) -- (axis cs:0.2, 6.4) node[midway, right, popgreen] {$\Delta U = \SI{0,4}{\volt}$};
\draw[popgreen, thick, <->] (axis cs:0.2, 6) -- (axis cs:0.4, 6) node[midway, above, popgreen] {$\Delta I = \SI{0,2}{\ampere}$};
% Labels
\node[poporange, anchor=east] at (axis cs:-0.02, 6) {$E' = \SI{6,0}{\volt}$};
\node[popblue, anchor=south west] at (axis cs:0.6, 7.2) {$U = E' + r'I$};
\end{axis}
\end{tikzpicture}
\end{center}
\legende{Caractéristique $U=f(I)$ d'un récepteur actif. Droite croissante. $E'$ est l'ordonnée à l'origine ($I=0$). $r'$ est le coefficient directeur.}
\end{popfigure}

\begin{methode}[Déterminer $E'$ et $r'$ à partir du graphique]
\begin{enumerate}
\item \textbf{Tracer la droite} de régression (ou la droite qui passe au plus près des points expérimentaux) sur le graphique $U = f(I)$.
\item \textbf{Lire $E'$ :} C'est l'\textbf{ordonnée à l'origine} (intersection de la droite avec l'axe des ordonnées $U$, là où $I=0$). Physiquement, c'est la tension aux bornes du récepteur quand le courant est nul (moteur à l'arrêt, électrolyseur sans courant).
\item \textbf{Calculer $r'$ :} Choisir \textbf{deux points éloignés} $A(I_A; U_A)$ et $B(I_B; U_B)$ \textbf{sur la droite tracée} (pas nécessairement des points de mesure). Calculer le coefficient directeur :
\[
r' = \frac{U_B - U_A}{I_B - I_A} = \frac{\Delta U}{\Delta I}
\]
\item \textbf{Vérifier l'unité :} $r'$ s'exprime en ohms (\si{\ohm}).
\end{enumerate}
\end{methode}

\begin{exoresolu}[Détermination de $E'$ et $r'$ pour un moteur]
\textbf{Énoncé :} On a relevé les points suivants sur la caractéristique d'un moteur à courant continu : $A(\SI{0,20}{\ampere} ; \SI{6,4}{\volt})$ et $B(\SI{0,60}{\ampere} ; \SI{7,2}{\volt})$. Déterminer $E'$ et $r'$.

\textbf{Solution détaillée :}
\begin{enumerate}
\item \textbf{Ordonnée à l'origine $E'$ :} On utilise l'équation de la droite $U = E' + r'I$. On a deux équations :
\[
\begin{cases}
\num{6,4} = E' + r' \times \num{0,20} \\
\num{7,2} = E' + r' \times \num{0,60}
\end{cases}
\]
\item \textbf{Calcul de $r'$ (coefficient directeur) :}
\[
r' = \frac{\Delta U}{\Delta I} = \frac{\num{7,2} - \num{6,4}}{\num{0,60} - \num{0,20}} = \frac{\num{0,8}}{\num{0,40}} = \SI{2,0}{\ohm}
\]
\item \textbf{Calcul de $E'$ :} On reporte $r'$ dans la première équation :
\[
E' = \num{6,4} - \num{2,0} \times \num{0,20} = \num{6,4} - \num{0,4} = \SI{6,0}{\volt}
\]
\item \textbf{Vérification avec le second point :} $U = \num{6,0} + \num{2,0} \times \num{0,60} = \SI{7,2}{\volt}$. C'est correct.
\end{enumerate}
\textbf{Résultat :} $\boxed{E' = \SI{6,0}{\volt}}$ et $\boxed{r' = \SI{2,0}{\ohm}}$.
\end{exoresolu}

\courssub{Cas particulier : Le récepteur passif (Résistance ohmique)}

Un récepteur passif (résistance, filament de lampe à l'équilibre thermique) ne possède \textbf{pas} de source de tension interne. Sa c.e.m. est nulle : $E' = \SI{0}{\volt}$. Sa résistance interne $r'$ est simplement sa résistance $R$.

La loi $U = E' + r'I$ devient :
\[
U = 0 + RI \quad \Rightarrow \quad \boxed{U = RI}
\]
C'est la \textbf{loi d'Ohm}. La caractéristique est une droite \textbf{passant par l'origine} $(0,0)$, de coefficient directeur $R$.

\begin{popfigure}
\begin{center}
\begin{tikzpicture}
\begin{axis}[
    width=\linewidth, height=0.6\linewidth,
    axis lines=middle,
    xlabel={$I$ (\si{\ampere})}, ylabel={$U$ (\si{\volt})},
    xmin=-0.05, xmax=0.7,
    ymin=-0.2, ymax=7.5,
    xtick={0,0.2,0.4,0.6},
    ytick={0,2,4,6,7.2},
    grid=major, grid style={popmuted},
    tick label style={font=\small},
    label style={font=\normalsize},
    clip=false,
    legend style={at={(0.98,0.02)}, anchor=south east, font=\small, fill=white, draw=popline}
]
% Recepteur Actif
\addplot[popcurve, domain=0:0.65, samples=2, thick, popblue] {6 + 2*x};
\addlegendentry{Récepteur actif : $U = 6 + 2I$}
% Resistor (Ohm)
\addplot[popcurve, domain=0:0.65, samples=2, thick, poporange, dashed] {10*x};
\addlegendentry{Résistance $R=\SI{10}{\ohm}$ : $U = 10I$}
% Point origine
\addplot[only marks, mark=*, mark size=2, popdark] coordinates {(0,0)};
\end{axis}
\end{tikzpicture}
\end{center}
\legende{Comparaison : Récepteur actif (droite croissante ne passant pas par l'origine, $E' > 0$) vs Récepteur passif / Résistance (droite passant par l'origine, $E'=0$).}
\end{popfigure}

\courssub{Comparaison synthétique : Générateur vs Récepteur actif}

Il est crucial de ne pas confondre les deux modèles. Le tableau ci-dessous résume les différences fondamentales.

\begin{center}
\begin{tabularx}{\linewidth}{|l|X|X|}\hline
\rowcolor{popblueL}
\textbf{Caractéristique} & \textbf{Générateur (Source)} & \textbf{Récepteur Actif (Charge active)} \\ \hline
\textbf{Modèle de Thévenin} & $E$ (source) en série avec $r$ & $E'$ (c.e.m.) en série avec $r'$ \\ \hline
\textbf{Équation} & $U = E - rI$ & $U = E' + r'I$ \\ \hline
\textbf{Pente (coeff. directeur)} & $-r$ (\textbf{négative}) & $+r'$ (\textbf{positive}) \\ \hline
\textbf{Ordonnée à l'origine ($I=0$)} & $U = E$ (FEM) & $U = E'$ (c.e.m.) \\ \hline
\textbf{Abscisse à l'origine ($U=0$)} & $I_{cc} = E/r$ (court-circuit) & $I = -E'/r'$ (non physiquement atteignable en récepteur) \\ \hline
\textbf{Sens du courant (convention)} & Sort par la borne $+$ (convention générateur) & Entre par la borne $+$ (convention récepteur) \\ \hline
\textbf{Énergie} & Fournit de l'énergie ($P = EI > 0$) & Convertit l'énergie électrique en autre forme (mécanique, chimique) \\ \hline
\textbf{Exemples} & Pile, batterie, alternateur, panneau solaire & Moteur, électrolyseur, batterie en charge \\ \hline
\end{tabularx}
\end{center}

\begin{aretenir}[Bilan de la section]
\begin{itemize}
\item Un \textbf{récepteur actif} se modélise par une \textbf{contre-électromotrice $E'$} en série avec une \textbf{résistance interne $r'$}.
\item La caractéristique $U = f(I)$ est une \textbf{droite croissante} d'équation $\boldsymbol{U = E' + r'I}$ (loi des mailles, convention récepteur).
\item \textbf{$E'$} (en \si{\volt}) = ordonnée à l'origine ($I=0$). C'est la tension que le récepteur « oppose » au passage du courant.
\item \textbf{$r'$} (en \si{\ohm}) = coefficient directeur de la droite ($\Delta U / \Delta I$). C'est la résistance interne dissipative.
\item Le \textbf{récepteur passif} (résistance) est le cas particulier $E' = 0$ : la droite passe par l'origine ($U = RI$).
\item \textbf{Ne jamais confondre} : Générateur $\rightarrow$ droite \textbf{décroissante} ($U = E - rI$) ; Récepteur actif $\rightarrow$ droite \textbf{croissante} ($U = E' + r'I$).
\end{itemize}
\end{aretenir}

\begin{attention}[Domaine de validité de la caractéristique]
La relation $U = E' + r'I$ est valide \textbf{uniquement en convention récepteur} ($I \ge 0$, $U \ge E'$).
\begin{itemize}
\item Si $U < E'$ (ce qui impose $I < 0$), le composant fonctionne en \textbf{générateur} (ex: moteur entraîné de l'extérieur devient dynamo, batterie se décharge). La caractéristique se prolonge alors par la même droite dans le quadrant $I < 0$, mais le composant n'est plus un récepteur.
\item \textbf{Ne pas extrapoler} la droite pour $U < E'$ en parlant de « fonctionnement en récepteur ». C'est une erreur de rédaction sanctionnée.
\end{itemize}
\end{attention}

\begin{savaistu}[Le moteur du « Bend Skin » et la c.e.m.]
Au Cameroun, les mototaxis (\emph{bend skin}) utilisent souvent des moteurs à courant continu (importés ou récupérés). Quand le moteur démarre, $I$ est maximal (car $U = E' + r'I$ et $E' \approx 0$ au démarrage car vitesse nulle $\Rightarrow$ c.e.m. nulle). Le courant d'appel peut être 5 à 10 fois le courant nominal ! C'est pourquoi les contrôleurs électroniques (variateurs) limitent le courant au démarrage pour ne pas griller les bobines ($r'I^2$ effet Joule) ni faire chuter la tension de la batterie. Une fois la vitesse de croisière atteinte, $E'$ monte (proportionnelle à la vitesse), $I$ baisse, le moteur chauffe moins. C'est la physique de la c.e.m. en action dans nos rues !
\end{savaistu}

\courssub{Point de fonctionnement d'un circuit générateur–récepteur}

Lorsque l'on connecte un générateur réel (FEM $E$, résistance interne $r$) à un récepteur actif (c.e.m. $E'$, résistance interne $r'$), le circuit est fermé. Le point de fonctionnement $(I, U)$ est l'unique couple de valeurs qui satisfait simultanément les deux caractéristiques. C'est l'intersection de la droite générateur et de la droite récepteur dans le repère $(I, U)$.

\begin{popfigure}
\begin{center}
\begin{tikzpicture}
\begin{axis}[
    width=\linewidth, height=0.65\linewidth,
    axis lines=middle,
    xlabel={$I$ (\si{\ampere})}, ylabel={$U$ (\si{\volt})},
    xmin=-0.1, xmax=2.5,
    ymin=-0.5, ymax=13,
    xtick={0,0.5,1,1.5,2,2.5},
    ytick={0,2,4,6,8,10,12},
    grid=major, grid style={popmuted},
    tick label style={font=\small},
    label style={font=\normalsize},
    clip=false,
    legend style={at={(0.98,0.98)}, anchor=north east, font=\small, fill=white, draw=popline}
]
% Generateur: U = 12 - 2I
\addplot[popcurve, domain=0:6, samples=2, thick, popblue] {12 - 2*x};
\addlegendentry{Générateur : $U = E - rI$}
% Recepteur: U = 6 + 2I
\addplot[popcurve, domain=0:3, samples=2, thick, poporange] {6 + 2*x};
\addlegendentry{Récepteur : $U = E' + r'I$}
% Intersection point
\addplot[only marks, mark=*, mark size=3, popdark] coordinates {(1.5, 9)};
% Dashed lines to axes
\draw[dashed, popmuted] (axis cs:1.5, 0) -- (axis cs:1.5, 9);
\draw[dashed, popmuted] (axis cs:0, 9) -- (axis cs:1.5, 9);
% Labels
\node[popblue, anchor=south] at (axis cs:0, 12) {$E = \SI{12}{\volt}$};
\node[poporange, anchor=east] at (axis cs:0, 6) {$E' = \SI{6}{\volt}$};
\node[popdark, anchor=south west] at (axis cs:1.5, 9) {$P_f(I_f, U_f)$};
\node[popdark, anchor=north] at (axis cs:1.5, -0.3) {$I_f$};
\node[popdark, anchor=east] at (axis cs:-0.1, 9) {$U_f$};
% Short circuit current gen
\node[popblue, anchor=north] at (axis cs:6, -0.3) {$I_{cc}$};
\end{axis}
\end{tikzpicture}
\end{center}
\legende{Point de fonctionnement $P_f$ : intersection des caractéristiques du générateur ($E=\SI{12}{\volt}, r=\SI{2}{\ohm}$) et du récepteur actif ($E'=\SI{6}{\volt}, r'=\SI{2}{\ohm}$).}
\end{popfigure}

\begin{propriete}[Détermination analytique du point de fonctionnement]
Le point de fonctionnement $(I_f, U_f)$ est solution du système :
\[
\begin{cases}
U_f = E - r I_f & \text{(Générateur)} \\
U_f = E' + r' I_f & \text{(Récepteur actif)}
\end{cases}
\]
En égalant les deux expressions de $U_f$ :
\[
E - r I_f = E' + r' I_f \quad \Rightarrow \quad \boxed{I_f = \frac{E - E'}{r + r'}}
\]
La tension aux bornes du récepteur (et du générateur) est alors :
\[
\boxed{U_f = E - r I_f = E' + r' I_f = \frac{E r' + E' r}{r + r'}}
\]

\textbf{Condition de fonctionnement en récepteur :} Pour que $I_f > 0$ (courant entrant dans le récepteur), il faut \textbf{impérativement $E > E'$}. Si $E < E'$, le courant est négatif : le « récepteur » impose sa loi et fonctionne en générateur (ex: freinage régénératif, batterie qui se décharge dans un circuit).
\end{propriete}

\begin{methode}[Résoudre un circuit Générateur – Récepteur actif]
\begin{enumerate}
\item \textbf{Identifier} les paramètres : $E, r$ (générateur) et $E', r'$ (récepteur).
\item \textbf{Vérifier la condition} $E > E'$ (sinon, le récepteur agit en générateur).
\item \textbf{Calculer le courant} $I_f = \dfrac{E - E'}{r + r'}$.
\item \textbf{Calculer la tension} $U_f = E - r I_f$ (ou $U_f = E' + r' I_f$).
\item \textbf{Bilan de puissance (optionnel mais recommandé) :}
\begin{itemize}
\item Puissance fournie par le générateur : $P_{gen} = E I_f$.
\item Puissance absorbée par le récepteur : $P_{rec} = U_f I_f = E' I_f + r' I_f^2$ (énergie convertie + pertes Joule dans $r'$).
\item Puissance dissipée dans $r$ : $P_r = r I_f^2$.
\item Vérifier : $P_{gen} = P_{rec} + P_r$.
\end{itemize}
\end{enumerate}
\end{methode}

\begin{exoresolu}[Charge d'une batterie par un panneau solaire]
\textbf{Énoncé :} Un panneau solaire (modélisé par $E = \SI{18,0}{\volt}$, $r = \SI{0,50}{\ohm}$) charge une batterie \SI{12}{\volt} (modélisée en récepteur actif par $E' = \SI{12,0}{\volt}$, $r' = \SI{0,20}{\ohm}$). Déterminer le courant de charge $I_f$, la tension aux bornes de la batterie $U_f$ et la puissance convertie en énergie chimique dans la batterie.

\textbf{Solution détaillée :}
\begin{enumerate}
\item \textbf{Condition :} $E = \SI{18,0}{\volt} > E' = \SI{12,0}{\volt}$. Le fonctionnement en récepteur (charge) est possible.
\item \textbf{Courant de charge :}
\[
I_f = \frac{E - E'}{r + r'} = \frac{\SI{18,0}{\volt} - \SI{12,0}{\volt}}{\SI{0,50}{\ohm} + \SI{0,20}{\ohm}} = \frac{\SI{6,0}{\volt}}{\SI{0,70}{\ohm}} \approx \SI{8,57}{\ampere}
\]
\item \textbf{Tension aux bornes de la batterie :}
\[
U_f = E' + r' I_f = \SI{12,0}{\volt} + \SI{0,20}{\ohm} \times \SI{8,57}{\ampere} \approx \SI{12,0}{\volt} + \SI{1,71}{\volt} = \SI{13,71}{\volt}
\]
(Vérification : $U_f = E - r I_f = \SI{18,0}{\volt} - \SI{0,50}{\ohm} \times \SI{8,57}{\ampere} = \SI{13,71}{\volt}$).
\item \textbf{Puissance convertie (énergie chimique stockée) :}
\[
P_{chimique} = E' I_f = \SI{12,0}{\volt} \times \SI{8,57}{\ampere} \approx \SI{102,8}{\watt}
\]
\item \textbf{Bilan :} $P_{gen} = E I_f \approx \SI{154,3}{\watt}$. $P_{Joule, r} = r I_f^2 \approx \SI{36,7}{\watt}$. $P_{Joule, r'} = r' I_f^2 \approx \SI{14,7}{\watt}$. $P_{rec} = P_{chimique} + P_{Joule, r'} \approx \SI{117,5}{\watt}$. $P_{gen} = P_{rec} + P_{Joule, r}$. OK.
\end{enumerate}
\textbf{Résultat :} $\boxed{I_f \approx \SI{8,6}{\ampere}}$, $\boxed{U_f \approx \SI{13,7}{\volt}}$, $\boxed{P_{chimique} \approx \SI{103}{\watt}}$.
\end{exoresolu}

\begin{attention}[Erreurs fréquentes au Probatoire sur le point de fonctionnement]
\begin{itemize}
\item \textbf{Oublier la condition $E > E'$} et calculer un courant négatif sans l'interpréter (le composant n'est plus un récepteur).
\item \textbf{Confondre $U_f$ et $E'$ :} $U_f$ est la tension \emph{mesurée} aux bornes (borne + et - du récepteur). Elle est \textbf{toujours supérieure à $E'$} en fonctionnement récepteur ($U_f = E' + r'I_f$).
\item \textbf{Mélanger les conventions :} Utiliser $U = E - rI$ pour le générateur (convention générateur) et $U = E' + r'I$ pour le récepteur (convention récepteur). La tension $U$ et le courant $I$ sont les \textbf{mêmes} pour les deux (loi des nœuds et des mailles). Ne pas changer le signe de $I$ entre les deux équations.
\item \textbf{Puissance absorbée :} La puissance \emph{absorbée} par le récepteur est $P_{abs} = U_f I_f$. La puissance \emph{convertie} (utile) est $P_{conv} = E' I_f$. La différence $r' I_f^2$ est perdue en chaleur. Ne pas confondre les deux.
\end{itemize}
\end{attention}

\begin{exercice}[Application : Caractéristique d'un électrolyseur]
On réalise le montage de l'expérience pour tracer la caractéristique d'un électrolyseur cuivre-cuivre (solution de sulfate de cuivre). On obtient les points suivants :
$M_1(\SI{0,05}{\ampere} ; \SI{1,20}{\volt})$, $M_2(\SI{0,15}{\ampere} ; \SI{1,40}{\volt})$, $M_3(\SI{0,25}{\ampere} ; \SI{1,60}{\volt})$, $M_4(\SI{0,35}{\ampere} ; \SI{1,80}{\volt})$.
\begin{enumerate}
\item Tracer la caractéristique $U = f(I)$ (échelle suggérée : 1 cm pour \SI{0,05}{\ampere} en abscisse, 1 cm pour \SI{0,1}{\volt} en ordonnée).
\item Montrer que les points sont alignés. Déterminer graphiquement $E'$ et $r'$.
\item Vérifier par le calcul avec $M_1$ et $M_4$.
\item Quelle est la tension minimale à appliquer aux bornes de l'électrolyseur pour qu'un courant puisse le traverser ? (Interprétation physique de $E'$).
\end{enumerate}
\end{exercice}

\begin{corrige}[Exercice : Caractéristique d'un électrolyseur]
\begin{enumerate}
\item \textbf{Graphique :} Les points s'alignent sur une droite croissante.
\item \textbf{Graphique :} Lecture de l'ordonnée à l'origine : pour $I=0$, $U \approx \SI{1,10}{\volt}$. Donc $\boldsymbol{E' \approx \SI{1,10}{\volt}}$.
Pente : $r' = \frac{\SI{1,80}{\volt} - \SI{1,20}{\volt}}{\SI{0,35}{\ampere} - \SI{0,05}{\ampere}} = \frac{\SI{0,60}{\volt}}{\SI{0,30}{\ampere}} = \SI{2,0}{\ohm}$.
\item \textbf{Calcul :} $r' = \frac{U_4 - U_1}{I_4 - I_1} = \frac{\SI{1,80}{\volt} - \SI{1,20}{\volt}}{\SI{0,35}{\ampere} - \SI{0,05}{\ampere}} = \frac{\SI{0,60}{\volt}}{\SI{0,30}{\ampere}} = \SI{2,0}{\ohm}$.
$E' = U_1 - r'I_1 = \SI{1,20}{\volt} - \SI{2,0}{\ohm} \times \SI{0,05}{\ampere} = \SI{1,20}{\volt} - \SI{0,10}{\volt} = \SI{1,10}{\volt}$.
Vérification avec $M_4$ : $U = \SI{1,10}{\volt} + \SI{2,0}{\ohm} \times \SI{0,35}{\ampere} = \SI{1,80}{\volt}$. OK.
\item \textbf{Interprétation :} Il faut appliquer une tension $U > E' = \SI{1,10}{\volt}$ pour faire passer un courant. Cette tension $E'$ correspond à la \textbf{tension de décomposition} de l'électrolyte (seuil thermodynamique pour la réaction $\ce{Cu^{2+} + 2e- -> Cu}$ à la cathode et $\ce{Cu -> Cu^{2+} + 2e-}$ à l'anode). En dessous, l'électrolyseur se comporte comme un circuit ouvert (résistance infinie).
\end{enumerate}
\end{corrige}

\coursec{Point de fonctionnement d'un circuit générateur–récepteur}

\courssub{Transition : du modèle au circuit réel}

Dans les sections précédentes, nous avons modélisé un générateur réel par une source de tension idéale $E$ en série avec une résistance interne $r$, et un récepteur actif (moteur, batterie en charge, électrolyseur) par un générateur de tension inverse $E'$ en série avec une résistance $r'$. Nous savons tracer la caractéristique $U = f(I)$ d'un générateur (droite de pente $-r$, ordonnée à l'origine $E$) et celle d'un récepteur actif (droite de pente $+r'$, ordonnée à l'origine $E'$).

Nous allons maintenant \textbf{connecter} ces deux modèles. C'est l'étape fondamentale : un générateur ne sert à rien s'il n'alimente pas un récepteur. Le circuit le plus simple est une \textbf{liaison série} entre un générateur $(E, r)$ et un récepteur $(E', r')$. Le courant $I$ est unique dans la maille. La tension $U$ aux bornes du générateur est \emph{exactement} la tension aux bornes du récepteur (loi des mailles, les deux dipôles sont en parallèle aux bornes A et B). Le point de fonctionnement $(I_0, U_0)$ est l'unique couple $(I, U)$ qui satisfait simultanément les deux caractéristiques.

\begin{popfigure}
\begin{tikzpicture}[scale=1.1, european]
% Générateur
\draw (0,0) to[battery1, l_={$E$}, invert] (0,2)
      to[R=$r$] (2,2);
% Récepteur
\draw (2,2) to[R=$r'$] (4,2)
      to[battery1, l_={$E'$}, invert] (4,0);
% Fermeture
\draw (4,0) -- (0,0);
% Courant
\draw[->, thick, popblue] (1, 2.5) -- (3, 2.5) node[midway, above] {$I$};
% Tension aux bornes
\draw[<->, poporange] (0, -0.5) -- (4, -0.5) node[midway, below] {$U$};
% Bornes
\node[left] at (0,2) {A};
\node[left] at (0,0) {B};
\node[right] at (4,2) {A'};
\node[right] at (4,0) {B'};
\end{tikzpicture}
\legende{Circuit série générateur $(E, r)$ -- récepteur actif $(E', r')$. Le courant $I$ est unique ; la tension $U_{AB} = U_{A'B'}$.}
\end{popfigure}

\courssub{Établissement de l'intensité du courant : bilan de tensions}

Appliquons la \textbf{loi des mailles} (conséquence de la conservation de l'énergie par unité de charge) dans le sens du courant $I$ (convention générateur pour le générateur, convention récepteur pour le récepteur).

\begin{itemize}
    \item Aux bornes du générateur (convention générateur) : $U = E - rI$.
    \item Aux bornes du récepteur (convention récepteur) : $U = E' + r'I$.
\end{itemize}

Comme la tension $U$ est commune aux deux dipôles, nous égalons les deux expressions :

\[
E - rI = E' + r'I
\]

\begin{experience}[Démonstration guidée : détermination de $I$]
\textbf{Objectif :} Isoler $I$ dans l'équation $E - rI = E' + r'I$.

\begin{enumerate}
    \item Regrouper les termes en $I$ à droite : $E = E' + r'I + rI$.
    \item Factoriser par $I$ : $E = E' + (r + r')I$.
    \item Transposer $E'$ : $E - E' = (r + r')I$.
    \item Diviser par la résistance totale $r + r'$ (non nulle) :
    \[
    I = \frac{E - E'}{r + r'}
    \]
\end{enumerate}
\end{experience}

\begin{propriete}[Intensité dans un circuit générateur–récepteur série]
Dans un circuit série constitué d'un générateur $(E, r)$ et d'un récepteur actif $(E', r')$, l'intensité du courant est donnée par :
\[
\boxed{I = \frac{E - E'}{r + r'}}
\]
\emph{Condition d'existence d'un courant non nul :} $E \neq E'$. Si $E > E'$, $I > 0$ : le générateur charge le récepteur (mode charge). Si $E < E'$, $I < 0$ : le récepteur devient générateur et décharge dans l'autre (mode décharge). \textbf{Démonstration exigible au Probatoire.}
\end{propriete}

\courssub{Point de fonctionnement : intersection graphique}

Le point de fonctionnement $P_0(I_0, U_0)$ est l'intersection des deux droites caractéristiques dans le plan $(I, U)$ (abscisse $I$, ordonnée $U$).

\begin{itemize}
    \item Droite générateur : $U = E - rI$ (pente $-r$, ordonnée à l'origine $E$).
    \item Droite récepteur : $U = E' + r'I$ (pente $+r'$, ordonnée à l'origine $E'$).
\end{itemize}

L'abscisse de l'intersection est $I_0 = \dfrac{E - E'}{r + r'}$. L'ordonnée $U_0$ s'obtient en remplaçant $I_0$ dans l'une ou l'autre équation :
\[
U_0 = E - r I_0 = E' + r' I_0 = \frac{E r' + E' r}{r + r'}
\]

\begin{popfigure}
\begin{tikzpicture}
\begin{axis}[
    width=\linewidth, height=9cm,
    axis lines=middle,
    xlabel={$I$ (A)}, ylabel={$U$ (V)},
    xmin=-0.5, xmax=3.5,
    ymin=-0.5, ymax=13.5,
    xtick={0,1,2,3}, ytick={0,2,4,6,8,10,12},
    grid=major, grid style={popline},
    clip=false,
    legend style={at={(0.98,0.98)}, anchor=north east, fill=white, draw=popline},
    declare function={
        E=12; r=1; Ep=10.5; rp=0.5;
        I0 = (E-Ep)/(r+rp);
        U0 = E - r*I0;
    }
]
% Droite générateur
\addplot[popcurve, domain=0:3, samples=2, thick] ({x}, {E - r*x});
\addlegendentry{Générateur : $U = E - rI$}

% Droite récepteur
\addplot[poporange, domain=0:3, samples=2, thick] ({x}, {Ep + rp*x});
\addlegendentry{Récepteur : $U = E' + r'I$}

% Point de fonctionnement
\addplot[only marks, mark=*, mark size=3, popdark] coordinates {(1, 11)};
\node[above right, popdark, font=\bfseries] at (axis cs:1,11) {$P_0(1,0\,\text{A},\,11,0\,\text{V})$};

% Lignes pointillées vers axes
\draw[dashed, popmuted] (axis cs:1,0) -- (axis cs:1,11);
\draw[dashed, popmuted] (axis cs:0,11) -- (axis cs:1,11);
\node[below, popmuted] at (axis cs:1,0) {$I_0$};
\node[left, popmuted] at (axis cs:0,11) {$U_0$};

% Ordonnées à l'origine
\node[left, popblue] at (axis cs:0,12) {$E$};
\node[left, poporange] at (axis cs:0,10.5) {$E'$};
\end{axis}
\end{tikzpicture}
\legende{Détermination graphique du point de fonctionnement $P_0$. Exemple : Chargeur $E=12\,\text{V}, r=1\,\Omega$ ; Batterie $E'=10,5\,\text{V}, r'=0,5\,\Omega$. Intersection en $I_0=1,0\,\text{A}, U_0=11,0\,\text{V}$.}
\end{popfigure}

\courssub{Bilan de puissance et rendement}

Une fois le point de fonctionnement connu, l'analyse énergétique s'impose. C'est un classique du Probatoire.

\begin{itemize}
    \item \textbf{Puissance fournie par le générateur} (puissance électromagnétique) : $P_{\text{four}} = E I_0$.
    \item \textbf{Puissance utile absorbée par le récepteur} (puissance convertie en énergie chimique pour une batterie, mécanique pour un moteur) : $P_{\text{utile}} = E' I_0$.
    \item \textbf{Pertes par effet Joule} (dans les résistances internes $r$ et $r'$) : $P_{\text{Joule}} = r I_0^2 + r' I_0^2 = (r + r') I_0^2$.
\end{itemize}

Le bilan de puissance (conservation de l'énergie) s'écrit :
\[
P_{\text{four}} = P_{\text{utile}} + P_{\text{Joule}} \quad \Leftrightarrow \quad E I_0 = E' I_0 + (r + r') I_0^2
\]
On retrouve l'équation des mailles multipliée par $I_0$.

\begin{definition}[Rendement d'un circuit générateur–récepteur]
Le rendement $\eta$ (lettre grecqueêta) est le rapport de la puissance utile sur la puissance fournie :
\[
\eta = \frac{P_{\text{utile}}}{P_{\text{four}}} = \frac{E' I_0}{E I_0} = \frac{E'}{E}
\]
\textbf{Remarque cruciale :} Le rendement ne dépend \textbf{que} des f.é.m. $E$ et $E'$, pas des résistances internes $r, r'$ (pourvu que $I_0 \neq 0$). Les résistances internes influencent l'intensité $I_0$ et donc les puissances \emph{absolues}, mais pas leur \emph{rapport}.
\end{definition}

\begin{attention}[Pièges classiques au Probatoire]
\begin{enumerate}
    \item \textbf{Confondre $U_0$ et $E$ ou $E'$ :} $U_0$ est la tension \emph{réelle} aux bornes (mesurable au voltmètre). Elle vaut $E - rI_0$ (côté générateur) ou $E' + r'I_0$ (côté récepteur). Elle est \emph{toujours} comprise entre $E'$ et $E$ (si $E > E'$).
    \item \textbf{Signe du courant :} Vérifier toujours que $E > E'$ pour une charge. Si l'énoncé donne $E_{\text{chargeur}} < E'_{\text{batterie}}$, le courant est négatif : la batterie décharge dans le chargeur (danger !).
    \item \textbf{Rendement vs Puissance :} Un rendement élevé ($E' \approx E$) implique un courant faible ($E-E'$ petit) donc une puissance utile faible. Il y a un compromis entre rendement et puissance transférée (théorème de la puissance maximale : $r = r'$ pour puissance max, mais rendement = 50\,\%).
    \item \textbf{Unités :} Puissances en \si{\watt} (W), énergies en \si{\joule} (J) ou \si{\watt\hour} (Wh). $1\,\text{Wh} = 3600\,\text{J}$.
\end{enumerate}
\end{attention}

\courssub{Méthode : déterminer le point de fonctionnement}

\begin{methode}[Point de fonctionnement $(I_0, U_0)$]
\textbf{Données :} Générateur $(E, r)$, Récepteur $(E', r')$.

\textbf{Par le calcul (recommandé, plus précis) :}
\begin{enumerate}
    \item Calculer $I_0 = \dfrac{E - E'}{r + r'}$.
    \item Vérifier le signe : $I_0 > 0$ (charge) ou $I_0 < 0$ (décharge).
    \item Calculer $U_0 = E - r I_0$ (ou $E' + r' I_0$).
\end{enumerate}

\textbf{Graphiquement (sur papier millimétré) :}
\begin{enumerate}
    \item Tracer le repère $(I, U)$. Choisir des échelles adaptées (origines non nécessairement communes).
    \item Tracer la droite générateur : 2 points suffisent. Ex: $I=0 \to U=E$ ; $U=0 \to I=E/r$ (court-circuit).
    \item Tracer la droite récepteur : 2 points. Ex: $I=0 \to U=E'$ ; $U=0 \to I=-E'/r'$ (prolongement vers $I<0$).
    \item Lire les coordonnées $(I_0, U_0)$ de l'intersection.
\end{enumerate}
\end{methode}

\courssub{Applications numériques types}

\begin{exoresolu}[Charge d'une batterie par un chargeur secteur]
\textbf{Énoncé :} Un chargeur de batterie (modélisé par $E = 12,0\,\text{V}$, $r = 1,00\,\Omega$) est branché sur une batterie déchargée (modélisée par $E' = 10,5\,\text{V}$, $r' = 0,50\,\Omega$). On néglige la résistance des fils.
\begin{enumerate}
    \item Déterminer l'intensité $I_0$ et la tension $U_0$ aux bornes de la batterie pendant la charge.
    \item Calculer la puissance fournie par le chargeur, la puissance utile stockée dans la batterie, les pertes Joule et le rendement.
    \item La charge dure $5,0\,\text{h}$. Quelle énergie (en Wh et en J) a été stockée dans la batterie ?
\end{enumerate}
\tcblower
\textbf{Solution détaillée :}

\begin{enumerate}
    \item \textbf{Point de fonctionnement :}
    \[
    I_0 = \frac{E - E'}{r + r'} = \frac{12,0 - 10,5}{1,00 + 0,50} = \frac{1,5}{1,5} = \num{1,00}\,\text{A}
    \]
    $I_0 > 0$ : la charge a bien lieu (le chargeur impose son sens).
    \[
    U_0 = E - r I_0 = 12,0 - 1,00 \times 1,00 = \num{11,0}\,\text{V}
    \]
    (Vérification : $U_0 = E' + r' I_0 = 10,5 + 0,50 \times 1,00 = 11,0\,\text{V}$).
    La tension aux bornes de la batterie ($11,0\,\text{V}$) est supérieure à sa f.é.m. ($10,5\,\text{V}$) : c'est la signature de la charge.

    \item \textbf{Bilan de puissance :}
    \begin{align*}
        P_{\text{four}} &= E I_0 = 12,0 \times 1,00 = \num{12,0}\,\text{W} \\
        P_{\text{utile}} &= E' I_0 = 10,5 \times 1,00 = \num{10,5}\,\text{W} \\
        P_{\text{Joule}} &= (r + r') I_0^2 = 1,5 \times 1,00^2 = \num{1,50}\,\text{W} \\
        \text{Vérification :} & \quad 10,5 + 1,5 = 12,0\,\text{W} \quad \checkmark \\
        \eta &= \frac{E'}{E} = \frac{10,5}{12,0} = 0,875 = \text{\SI{87,5}{\percent}}
    \end{align*}

    \item \textbf{Énergie stockée :}
    \[
    W_{\text{utile}} = P_{\text{utile}} \times \Delta t = 10,5\,\text{W} \times 5,0\,\text{h} = \num{52,5}\,\text{Wh}
    \]
    \[
    W_{\text{utile}} = 52,5 \times 3600 = \num{189000}\,\text{J} = \num{1,89 \times 10^5}\,\text{J}
    \]
\end{enumerate}
\end{exoresolu}

\begin{exoresolu}[Moteur alimenté par une batterie]
\textbf{Énoncé :} Un moteur de modélisme (récepteur actif) est alimenté par une batterie Li-ion. Caractéristiques : Batterie $E = 11,1\,\text{V}$, $r = 0,15\,\Omega$ ; Moteur $E' = 9,0\,\text{V}$ (f.é.m. de rotation), $r' = 0,35\,\Omega$.
Déterminer $I_0$, $U_0$, la puissance mécanique développée et le rendement.
\tcblower
\textbf{Solution :}
\begin{align*}
    I_0 &= \frac{11,1 - 9,0}{0,15 + 0,35} = \frac{2,1}{0,50} = \num{4,2}\,\text{A} \\
    U_0 &= 11,1 - 0,15 \times 4,2 = 11,1 - 0,63 = \num{10,5}\,\text{V} \\
    P_{\text{méca}} &= E' I_0 = 9,0 \times 4,2 = \num{37,8}\,\text{W} \\
    \eta &= \frac{9,0}{11,1} \approx 0,811 = \text{\SI{81,1}{\percent}}
\end{align*}
\textit{Note :} La tension aux bornes du moteur ($10,5\,\text{V}$) est bien supérieure à sa f.é.m. de rotation ($9,0\,\text{V}$). La différence ($1,5\,\text{V}$) correspond à la chute de tension interne $r'I_0$. Les valeurs sont données avec 2 chiffres significatifs cohérents avec les données ($r, r', E'$ à 2 chiffres).
\end{exoresolu}

\courssub{Le savais-tu ? : L'électronique de puissance au service de la batterie}

\begin{savaistu}[Régulateurs solaires MPPT et BMS]
Au Cameroun, les kits solaires domestiques (panneau + batterie + régulateur) sont omniprésents. Le panneau photovoltaïque est un générateur dont la caractéristique $U=f(I)$ n'est \emph{pas} une droite (c'est une courbe non-linéaire avec un point de puissance maximale). La batterie est un récepteur actif $(E', r')$.

Le \textbf{régulateur de charge} (souvent technologie MPPT : \emph{Maximum Power Point Tracking}) est un convertisseur électronique (hacheur) qui adapte en permanence la tension vue par le panneau pour le faire fonctionner à son point de puissance maximale, tout en imposant à la batterie le courant de charge optimal (souvent $I_{\text{charge}} \approx C/10$ à $C/5$).

Sans cet électronique, si on branchait directement le panneau sur la batterie, le point de fonctionnement serait l'intersection de la courbe panneau et de la droite batterie. Ce point est rarement optimal : soit la batterie se charge trop lentement (panneau sous-exploité), soit le courant est trop fort (risque de surchauffe, dégazage). Le régulateur \textbf{découple} le point de fonctionnement du panneau de celui de la batterie. C'est de l'électronique de puissance pure, mais la base physique reste notre équation $E - rI = E' + r'I$ généralisée à des dipôles non-linéaires pilotés par des transistors commutés à haute fréquence (kHz--MHz).
\end{savaistu}

\courssub{Synthèse : les 3 régimes d'un circuit générateur–récepteur}

\begin{center}
\begin{tabularx}{\linewidth}{|l|X|X|X|}\hline
\rowcolor{popblueL}
\textbf{Régime} & \textbf{Condition} & \textbf{Sens de $I$} & \textbf{Interprétation physique} \\ \hline
\textbf{Charge} & $E > E'$ & $I > 0$ (Géné $\to$ Récep) & Le générateur fournit l'énergie. Le récepteur stocke (batterie) ou convertit (moteur). $U_0 \in [E', E]$. \\ \hline
\textbf{Décharge} & $E < E'$ & $I < 0$ (Récep $\to$ Géné) & Le "récepteur" (ex: batterie pleine) devient source. Il alimente le "générateur" (ex: chargeur non protégé) ou une charge externe. Danger si non prévu. \\ \hline
\textbf{Équilibre} & $E = E'$ & $I = 0$ & Tension aux bornes $U_0 = E = E'$. Aucun transfert d'énergie. Flottement (ex: batterie branchée sur chargeur éteint ou réglé à tension exacte). \\ \hline
\end{tabularx}
\end{center}

\begin{aretenir}[Essentiel pour le Probatoire]
\begin{itemize}
    \item Circuit série $(E, r)$ -- $(E', r')$ : \textbf{un seul courant $I$}, \textbf{une seule tension $U$} aux bornes.
    \item Équation fondamentale : $E - rI = E' + r'I \;\Rightarrow\; I = \dfrac{E - E'}{r + r'}$.
    \item Point de fonctionnement $P_0(I_0, U_0)$ = intersection des droites $U = E - rI$ et $U = E' + r'I$.
    \item Puissances : $P_{\text{four}} = EI_0$, $P_{\text{utile}} = E'I_0$, $P_{\text{Joule}} = (r+r')I_0^2$.
    \item Rendement : $\eta = \dfrac{E'}{E}$ (indépendant de $r, r'$).
    \item Condition de charge : $E_{\text{géné}} > E'_{\text{récep}}$.
\end{itemize}
\end{aretenir}

\begin{exercice}[Entraînement : Panneau solaire et batterie]
Un panneau solaire est modélisé par un générateur de f.é.m. $E = 18,0\,\text{V}$ et résistance interne $r = 2,0\,\Omega$. Il charge une batterie $12\,\text{V}$ de f.é.m. $E' = 12,5\,\text{V}$ et résistance interne $r' = 0,4\,\Omega$.
\begin{enumerate}
    \item Calculer l'intensité de charge $I_0$ et la tension $U_0$ aux bornes de la batterie.
    \item Calculer le rendement de la charge.
    \item On souhaite augmenter le courant de charge à $2,0\,\text{A}$ sans changer la batterie. Quelle doit être la nouvelle f.é.m. $E_{\text{nouv}}$ du générateur (en gardant $r = 2,0\,\Omega$) ?
\end{enumerate}
\end{exercice}

\begin{corrige}[Exercice : Panneau solaire et batterie]
\begin{enumerate}
    \item $I_0 = \dfrac{18,0 - 12,5}{2,0 + 0,4} = \dfrac{5,5}{2,4} \approx \num{2,29}\,\text{A}$.
    $U_0 = 18,0 - 2,0 \times 2,29 = 18,0 - 4,58 = \num{13,42}\,\text{V}$ (ou $12,5 + 0,4 \times 2,29 = 13,42\,\text{V}$).
    \item $\eta = \dfrac{E'}{E} = \dfrac{12,5}{18,0} \approx 0,694 = \text{\SI{69,4}{\percent}}$.
    \item On veut $I_0 = 2,0\,\text{A}$. $E_{\text{nouv}} = E' + (r+r')I_0 = 12,5 + 2,4 \times 2,0 = 12,5 + 4,8 = \num{17,3}\,\text{V}$.
    (Le panneau actuel $18\,\text{V}$ donne déjà $2,29\,\text{A} > 2,0\,\text{A}$, donc on pourrait même baisser légèrement la tension du panneau ou ajouter une résistance série pour limiter à $2,0\,\text{A}$ exactement).
\end{enumerate}
\end{corrige}

\newpage

\coursec{Activité d'intégration}

\coursec{Leçon 14 : Générateur et récepteur électriques}

\courssub{Objectifs de l'activité}
Cette activité d'intégration vise à mobiliser l'ensemble des savoirs et savoir-faire de la leçon sur les générateurs et récepteurs électriques à travers une situation concrète d'électrification rurale au Cameroun. À l'issue de ce travail, tu seras capable de :
\begin{itemize}
    \item Modéliser un panneau photovoltaïque réel par un générateur de Thévenin (\cle{$E$}, \cle{$r$}) et une batterie en charge par un récepteur actif (\cle{$E'$}, \cle{$r'$}).
    \item Exploiter des relevés de mesures \cle{$(U, I)$} pour tracer les caractéristiques \cle{$U = f(I)$} et en extraire les paramètres \cle{$E$}, \cle{$r$}, \cle{$E'$}, \cle{$r'$} par lecture graphique.
    \item Déterminer le \cle{$point de fonctionnement$} d'un circuit \emph{générateur--récepteur} par résolution graphique et analytique.
    \item Analyser l'effet d'une association de générateurs (série ou parallèle) sur le courant de charge et justifier le choix technique optimal.
\end{itemize}

\courssub{Situation : Alimenter une maison de quartier avec un kit solaire}
Dans le quartier \emph{Nkolbisson} à Yaoundé, une famille souhaite s'équiper d'un kit solaire autonome pour l'éclairage et la recharge de téléphones, en complément du réseau ENEO souvent instable. Le technicien leur propose un kit composé d'un \textbf{panneau photovoltaïque} de technologie silicium monocristallin et d'une \textbf{batterie au plomb} 12 V / 50 Ah.

Pour valider le bon fonctionnement de l'installation avant l'achat, le technicien fournit les relevés de mesures suivants, effectués en conditions standard d'ensoleillement (irradiance 1000 W/m², température 25 °C) pour le panneau, et en charge lente pour la batterie.

\medskip
\textbf{1. Caractéristique du panneau photovoltaïque (modélisé en générateur)}
Le panneau est modélisé par un générateur idéal de tension \cle{$E$} en série avec une résistance interne \cle{$r$}. La tension \cle{$U$} à ses bornes et l'intensité \cle{$I$} qui le traverse (convention générateur : \cle{$I$} sort par la borne \cle{$+$}) ont été mesurées :

\begin{center}
\begin{tabularx}{\linewidth}{|c|*{6}{X|}}\hline
\rowcolor{popblueL}
$I$ (\si{\ampere}) & 0,0 & 1,0 & 2,0 & 3,0 & 4,0 & 5,0 \\ \hline
$U$ (\si{\volt})   & 20,0 & 16,0 & 12,0 & 8,0 & 4,0 & 0,0 \\ \hline
\end{tabularx}
\end{center}

\medskip
\textbf{2. Caractéristique de la batterie en charge (modélisée en récepteur actif)}
La batterie est modélisée par un récepteur actif : une force électromotrice de contre-électromotrice \cle{$E'$} (tension de repos) en série avec une résistance interne \cle{$r'$}. La tension \cle{$U$} à ses bornes et l'intensité \cle{$I$} de charge (convention récepteur : \cle{$I$} entre par la borne \cle{$+$}) ont été mesurées :

\begin{center}
\begin{tabularx}{\linewidth}{|c|*{5}{X|}}\hline
\rowcolor{popgreenL}
$I$ (\si{\ampere}) & 0,0 & 1,0 & 2,0 & 3,0 & 4,0 \\ \hline
$U$ (\si{\volt})   & 12,0 & 12,5 & 13,0 & 13,5 & 14,0 \\ \hline
\end{tabularx}
\end{center}

\medskip
\textbf{3. Extension : Ajout d'un second panneau identique}
Pour réduire le temps de charge, la famille envisage d'acheter un second panneau \emph{strictement identique} au premier. Le technicien propose deux montages :
\begin{itemize}
    \item \textbf{Montage A :} Les deux panneaux en \textbf{série}.
    \item \textbf{Montage B :} Les deux panneaux en \textbf{parallèle}.
\end{itemize}

\courssub{Consignes}
\emph{Travail en groupes de 3 ou 4. Chaque groupe rend un rapport unique sur papier millimétré (ou logiciel de tracé) avec les réponses argumentées.}

\begin{enumerate}
    \item \textbf{Modélisation et schémas.}
    \begin{enumerate}
        \item Dessiner le schéma équivalent du panneau (générateur) et de la batterie (récepteur) avec les polarités des tensions \cle{$U$} et sens des courants \cle{$I$} conformes aux conventions générateur et récepteur.
        \item Écrire les équations caractéristiques \cle{$U = f(I)$} pour le générateur et pour le récepteur actif.
    \end{enumerate}

    \item \textbf{Exploitation graphique des caractéristiques (sur un même repère).}
    \begin{enumerate}
        \item Choisir une échelle adaptée (ex: 1 cm pour 1 A en abscisses, 1 cm pour 2 V en ordonnées). Tracer les points des deux tableaux et les droites caractéristiques \cle{$U_G(I)$} et \cle{$U_R(I)$}.
        \item Déterminer \textbf{graphiquement} les quatre paramètres : \cle{$E$}, \cle{$r$}, \cle{$E'$}, \cle{$r'$}. Indiquer la méthode de lecture sur le graphique (droites, intersections, triangles).
    \end{enumerate}

    \item \textbf{Point de fonctionnement du circuit simple (1 panneau + 1 batterie).}
    \begin{enumerate}
        \item Relier le générateur et le récepteur par des fils de résistance négligeable. Écrire l'équation de maille (loi des mailles) et en déduire l'expression analytique du courant de charge \cle{$I_0$} et de la tension aux bornes \cle{$U_0$}.
        \item Calculer les valeurs numériques de \cle{$I_0$} et \cle{$U_0$}.
        \item Repérer le point de fonctionnement \cle{$P_0(I_0, U_0)$} sur le graphique de la question 2. Vérifier la cohérence graphique/analytique.
    \end{enumerate}

    \item \textbf{Étude de l'extension : deux panneaux pour charger plus vite.}
    \begin{enumerate}
        \item Pour le \textbf{Montage A (série)} : Déterminer la f.é.m. équivalente \cle{$E_s$} et la résistance interne équivalente \cle{$r_s$}. Calculer le nouveau courant de charge \cle{$I_s$}.
        \item Pour le \textbf{Montage B (parallèle)} : Déterminer \cle{$E_p$} et \cle{$r_p$}. Calculer le nouveau courant de charge \cle{$I_p$}.
        \item Comparer \cle{$I_s$}, \cle{$I_p$} et \cle{$I_0$}. Quel montage permet de charger la batterie \textbf{le plus rapidement} ? Justifier ta réponse physiquement (énergie, puissance).
        \item \textit{(Question bonus)} : La tension aux bornes de la batterie dépasse-t-elle 14,4 V (tension de fin de charge recommandée pour du plomb 12 V) dans l'un des montages ? Que risque-t-il de se passer ?
    \end{enumerate}
\end{enumerate}

\courssub{Ressources à mobiliser}
\begin{itemize}
    \item \textbf{Modèle de Thévenin d'un générateur réel :} \cle{$U = E - r I$} (convention générateur : \cle{$U$} et \cle{$I$} de même sens). \cle{$E$} = tension à vide (\cle{$I=0$}) ; \cle{$r$} = \cle{$-\Delta U / \Delta I$} = \cle{$E / I_{cc}$} (\cle{$I_{cc}$} = courant de court-circuit).
    \item \textbf{Modèle d'un récepteur actif :} \cle{$U = E' + r' I$} (convention récepteur : \cle{$U$} et \cle{$I$} de sens opposés, \cle{$I$} entre par \cle{$+$}). \cle{$E'$} = tension à vide (\cle{$I=0$}) ; \cle{$r'$} = \cle{$\Delta U / \Delta I$}.
    \item \textbf{Loi des mailles (KVL) :} Dans un circuit série, la somme algébrique des tensions est nulle. \cle{$E - r I = E' + r' I$}.
    \item \textbf{Point de fonctionnement :} Intersection des deux droites caractéristiques \cle{$U_G(I)$} et \cle{$U_R(I)$}. Coordonnées \cle{$(I_0, U_0)$}.
    \item \textbf{Associations de générateurs identiques (E, r) :}
    \begin{itemize}
        \item Série : \cle{$E_{eq} = 2E$}, \cle{$r_{eq} = 2r$}.
        \item Parallèle : \cle{$E_{eq} = E$}, \cle{$r_{eq} = r/2$}.
    \end{itemize}
    \item \textbf{Puissance fournie par le générateur :} \cle{$P_G = U I = E I - r I^2$}. Puissance absorbée par le récepteur : \cle{$P_R = U I = E' I + r' I^2$}. Puissance utile (stockée) : \cle{$P_{stock} = E' I$}. Puissance perdue par effet Joule : \cle{$P_J = r I^2 + r' I^2$}.
\end{itemize}

\begin{exoresolu}{Corrigé de l'activité d'intégration}
\textbf{1. Modélisation et schémas}
\begin{itemize}
    \item \textbf{Panneau (Générateur) :} Schéma : source de tension \cle{$E$} (borne \cle{$+$} en haut) -- résistance \cle{$r$} -- borne \cle{$-$}. Flèche courant \cle{$I$} sortant de \cle{$+$}. Équation : \cle{$U = E - r I$}.
    \item \textbf{Batterie (Récepteur actif) :} Schéma : source de tension \cle{$E'$} (borne \cle{$+$} en haut) -- résistance \cle{$r'$} -- borne \cle{$-$}. Flèche courant \cle{$I$} entrant par \cle{$+$}. Équation : \cle{$U = E' + r' I$}.
\end{itemize}

\begin{popfigure}
\begin{tikzpicture}[european, scale=0.9, every node/.style={font=\small}]
    % Generateur
    \draw (0,0) node[left] {$+$} to[short] (0.5,0)
        to[battery1, l_={$E$}, invert] (0.5,-2)
        to[R, l={$r$}] (2.5,-2)
        to[short] (3,-2) node[right] {$-$};
    \draw[->, thick, popblue] (0.5,-1) -- (0.5,-0.5) node[above, pos=0.5, left] {$I$};
    \draw[<->, poporange] (0.5,0.3) -- (2.5,0.3) node[midway, above] {$U$};
    \node at (1.5, 0.8) {\textbf{Générateur (Panneau)}};

    % Recepteur
    \begin{scope}[xshift=6cm]
    \draw (0,0) node[left] {$+$} to[short] (0.5,0)
        to[R, l={$r'$}] (2.5,0)
        to[battery1, l={$E'$}] (2.5,-2)
        to[short] (3,-2) node[right] {$-$};
    \draw[->, thick, popgreen] (0.5,-1) -- (0.5,-1.5) node[below, pos=0.5, left] {$I$};
    \draw[<->, poporange] (0.5,0.3) -- (2.5,0.3) node[midway, above] {$U$};
    \node at (1.5, 0.8) {\textbf{Récepteur (Batterie)}};
    \end{scope}
\end{tikzpicture}
\legende{Schéma équivalent du panneau photovoltaïque (générateur) et de la batterie (récepteur actif) avec conventions de récepteur et générateur.}
\end{popfigure}

\textbf{2. Exploitation graphique}
\begin{itemize}
    \item \textbf{Échelle :} Abscisses : 1 cm $\leftrightarrow$ 1 A. Ordonnées : 1 cm $\leftrightarrow$ 2 V. Origine (0,0) en bas à gauche.
    \item \textbf{Droite Générateur \cle{$U_G(I)$} :} Passe par \cle{$A(0; 20)$} et \cle{$B(5; 0)$}. Droite décroissante.
    \item \textbf{Droite Récepteur \cle{$U_R(I)$} :} Passe par \cle{$C(0; 12)$} et \cle{$D(4; 14)$}. Droite croissante.
    \item \textbf{Lecture des paramètres :}
    \begin{itemize}
        \item \cle{$E$} = Ordonnée à l'origine de \cle{$U_G$} = \textbf{20,0 V}. (Tension à vide).
        \item \cle{$r$} = Pente (en valeur absolue) de \cle{$U_G$} = \cle{$\frac{20-0}{0-5}$} = \textbf{4,0 \si{\ohm}}. (Ou \cle{$E/I_{cc}$} = 20/5 = 4 \si{\ohm}).
        \item \cle{$E'$} = Ordonnée à l'origine de \cle{$U_R$} = \textbf{12,0 V}. (Tension de repos / contre-électromotrice).
        \item \cle{$r'$} = Pente de \cle{$U_R$} = \cle{$\frac{14-12}{4-0}$} = $\frac{2}{4}$ = \textbf{0,5 \si{\ohm}}.
    \end{itemize}
\end{itemize}

\begin{popfigure}
\begin{tikzpicture}[scale=0.7]
\begin{axis}[
    width=\linewidth, height=10cm,
    axis lines=middle,
    xlabel={$I$ (\si{\ampere})}, ylabel={$U$ (\si{\volt})},
    xmin=-0.5, xmax=6, ymin=-1, ymax=22,
    xtick={0,1,2,3,4,5,6}, ytick={0,4,8,12,16,20},
    grid=major, grid style={popline},
    legend style={at={(0.98,0.98)}, anchor=north east, fill=white, draw=popdark},
    clip=false
]
% Generateur line
\addplot[popcurve, domain=0:5, samples=2, thick] {20 - 4*x};
\addlegendentry{$U_G = E - rI$ (Panneau)}
% Recepteur line
\addplot[popcurve, domain=0:5, samples=2, thick, popgreen] {12 + 0.5*x};
\addlegendentry{$U_R = E' + r'I$ (Batterie)}
% Points mesures Gen
\addplot[only marks, mark=*, popblue] coordinates {(0,20) (1,16) (2,12) (3,8) (4,4) (5,0)};
% Points mesures Rec
\addplot[only marks, mark=square*, popgreen] coordinates {(0,12) (1,12.5) (2,13) (3,13.5) (4,14)};
% Operating point
\addplot[only marks, mark=*, mark size=3, poporange] coordinates {(1.777, 12.888)};
\node[poporange, anchor=south west] at (axis cs:1.777, 12.888) {$P_0$};
% Dashed lines for parameters
\draw[dashed, popmuted] (axis cs:0,20) -- (axis cs:0,0) node[below, pos=0.5, left] {$E$};
\draw[dashed, popmuted] (axis cs:0,12) -- (axis cs:0,0) node[below, pos=0.5, left] {$E'$};
\end{axis}
\end{tikzpicture}
\legende{Caractéristiques \cle{$U = f(I)$} du générateur (bleu) et du récepteur (vert). Point de fonctionnement \cle{$P_0$} (orange).}
\end{popfigure}

\textbf{3. Point de fonctionnement (1 panneau + 1 batterie)}
\begin{itemize}
    \item \textbf{Maille :} Sens horaire. \cle{$U_{gén} + U_{récep} = 0$} $\Rightarrow$ \cle{$(E - r I) - (E' + r' I) = 0$} (car \cle{$U$} aux bornes du récepteur est une chute de tension dans le sens du courant).
    \item \cle{$E - E' = (r + r') I_0$} $\Rightarrow$ \cle{$I_0 = \frac{E - E'}{r + r'}$}.
    \item \textbf{Calcul :} \cle{$I_0 = \frac{20,0 - 12,0}{4,0 + 0,5} = \frac{8,0}{4,5} = 1,78 \si{\ampere}$} (arrondi à \num{1,78} A).
    \item \textbf{Tension \cle{$U_0$} :} \cle{$U_0 = E - r I_0 = 20,0 - 4,0 \times 1,78 = 20,0 - 7,11 = 12,89 \si{\volt}$}.
    \item \textit{Vérification côté récepteur :} \cle{$U_0 = E' + r' I_0 = 12,0 + 0,5 \times 1,78 = 12,0 + 0,89 = 12,89 \si{\volt}$}. \checkmark
    \item \textbf{Graphique :} Le point \cle{$P_0(1,78; 12,89)$} est bien l'intersection des deux droites tracées.
\end{itemize}

\textbf{4. Extension : Deuxième panneau identique}
\textit{Rappel : Panneau unique : \cle{$E = 20 V$}, \cle{$r = 4 \si{\ohm}$}.}

\begin{enumerate}
    \item \textbf{Montage A (Série) :}
    \cle{$E_s = 2E = 40,0 V$}, \cle{$r_s = 2r = 8,0 \si{\ohm}$}.
    \cle{$I_s = \frac{E_s - E'}{r_s + r'} = \frac{40,0 - 12,0}{8,0 + 0,5} = \frac{28,0}{8,5} = \textbf{3,29 \si{\ampere}}$}.
    
    \item \textbf{Montage B (Parallèle) :}
    \cle{$E_p = E = 20,0 V$}, \cle{$r_p = \frac{r}{2} = 2,0 \si{\ohm}$}.
    \cle{$I_p = \frac{E_p - E'}{r_p + r'} = \frac{20,0 - 12,0}{2,0 + 0,5} = \frac{8,0}{2,5} = \textbf{3,20 \si{\ampere}}$}.
    
    \item \textbf{Comparaison et choix :}
    \cle{$I_s \approx 3,29 A > I_p \approx 3,20 A > I_0 \approx 1,78 A$}.
    Les deux montages augmentent fortement le courant de charge par rapport au panneau unique (presque doublé). Le \textbf{montage série (A) donne un courant très légèrement supérieur} (environ 3\% de plus) pour ce modèle linéaire précis.
    
    \item \textbf{Analyse énergétique (Puissances) :}
    \begin{center}
    \begin{tabular}{|c|c|c|c|c|}\hline
    \rowcolor{popblueL}
    Montage & \cle{$I$} (A) & \cle{$U_0$} (V) & \cle{$P_{fournie}=UI$} (W) & \cle{$P_{stock}=E'I$} (W) \\ \hline
    1 Panneau & 1,78 & 12,89 & 22,9 & 21,3 \\ \hline
    Série (A) & 3,29 & 14,65 & 48,3 & 39,5 \\ \hline
    Parallèle (B) & 3,20 & 13,60 & 43,5 & 38,4 \\ \hline
    \end{tabular}
    \end{center}
    Le montage série fournit plus de puissance électrique totale (\cle{$P_{fournie}$}) et plus de puissance stockée (\cle{$P_{stock} = E'I$}). Il charge donc \textbf{plus vite}.
    
    \item \textbf{Question bonus (Tension de fin de charge) :}
    \begin{itemize}
        \item Montage Série : \cle{$U_0 = E' + r' I_s = 12,0 + 0,5 \times 3,29 = 13,65 V$}. (Ou \cle{$E_s - r_s I_s = 40 - 8 \times 3,29 = 13,68 V$}). \textbf{< 14,4 V}. Pas de surcharge immédiate.
        \item Montage Parallèle : \cle{$U_0 = 12,0 + 0,5 \times 3,20 = 13,60 V$}. \textbf{< 14,4 V}.
    \end{itemize}
    Dans les deux cas, la tension reste inférieure à 14,4 V pour ces courants. \textbf{Attention cependant :} Si la batterie se charge, \cle{$E'$} augmente (ex: passe à 13,5 V). En série, la tension aux bornes monterait très vite au-delà de 14,4 V (risque de surcharge, dégagement d'hydrogène, destruction). En pratique, un \textbf{régulateur de charge (MPPT ou PWM)} est \textbf{indispensable} pour couper le courant ou abaisser la tension quand \cle{$U$} atteint 14,4 V. Le montage série, avec sa haute tension à vide (40 V), impose un régulateur MPPT capable de descendre la tension ; le montage parallèle (20 V) est plus simple à réguler pour une batterie 12 V.
\end{enumerate}
\tcblower
\textbf{Commentaires pédagogiques pour l'enseignant :}
\begin{itemize}
    \item \textbf{Conventions de signes :} Insister sur le fait que pour le récepteur, \cle{$U = E' + r'I$} car \cle{$U$} et \cle{$I$} sont de sens opposés (convention récepteur). L'équation de maille s'écrit \cle{$E - rI = E' + r'I$} (somme des f.é.m. = somme des chutes de tension ohmiques).
    \item \textbf{Précision graphique :} Le tracé sur papier millimétré doit être soigné. L'intersection détermine \cle{$I_0$}. Une erreur de 0,5 mm sur \cle{$I$} donne ~0,05 A d'erreur. La méthode analytique sert de référence exacte.
    \item \textbf{Analyse dimensionnelle :} \cle{$[E] = V$}, \cle{$[r] = \Omega = V/A$}. \cle{$I = V/\Omega = A$}. Cohérent.
    \item \textbf{Choix Série vs Parallèle :} Contre-intuitivement pour des élèves qui pensent "parallèle = plus de courant", le série gagne ici car \cle{$E$} double alors que \cle{$r$} double aussi, mais \cle{$E'$} est fixe. Le "gain" en tension (numérateur) l'emporte sur la "perte" en résistance (dénominateur) car \cle{$E \gg E'$}. Si \cle{$E'$} était proche de \cle{$E$} (ex: charge batterie 24V avec panneaux 12V), le parallèle serait meilleur. C'est une excellente discussion sur l'adaptation de la source à la charge.
    \item \textbf{Contexte réel :} Au Cameroun, les kits solaires "Plug \& Play" (type \emph{Sun King}, \emph{d.light}, \emph{Bboxx}) intègrent le panneau, la batterie LiFePO4 ou Plomb, et \textbf{l'électronique de régulation (MPPT)} dans un boîtier unique. Cette activité montre la physique "brute" que l'électronique gère automatiquement.
\end{itemize}
\end{exoresolu}

\courssub{Bilan de l'activité}
Cette activité a permis de mettre en évidence les points clés suivants, essentiels pour le Probatoire :
\begin{enumerate}
    \item \textbf{Modélisation unifiée :} Un appareil réel (panneau, batterie) se modélise par un dipôle linéaire (Thévenin) caractérisé par deux paramètres seulement : une tension caractéristique (\cle{$E$} ou \cle{$E'$}) et une résistance interne (\cle{$r$} ou \cle{$r'$}). La nature (générateur ou récepteur) dicte le signe devant \cle{$rI$} dans l'équation \cle{$U = f(I)$}.
    \item \textbf{Caractéristique \cle{$U = f(I)$} :} C'est l'outil graphique central. \cle{$E$} (ou \cle{$E'$}) se lit à l'ordonnée à l'origine (\cle{$I=0$}). \cle{$r$} (ou \cle{$r'$}) est la pente (en valeur absolue pour le générateur). La droite est décroissante pour un générateur, croissante pour un récepteur actif.
    \item \textbf{Point de fonctionnement :} Il est l'unique solution du système des deux équations caractéristiques. Il se trouve à l'intersection des deux droites. Son détermination analytique (\cle{$I_0 = \frac{E-E'}{r+r'}$}) est plus précise que la lecture graphique mais cette dernière donne une excellente intuition visuelle (existence, unicité, influence des paramètres).
    \item \textbf{Association de générateurs :} Le choix entre série et parallèle n'est pas anodin. Il modifie \cle{$E_{eq}$} et \cle{$r_{eq}$} différemment. Pour charger une batterie de tension \cle{$E'$} fixe avec des panneaux de tension \cle{$E > E'$}, le montage série augmente la "force motrice" nette (\cle{$E_{eq}-E'$}) plus vite qu'il n'augmente la résistance totale, ce qui maximise le courant (et la puissance utile \cle{$E'I$}) \emph{dans le cadre de ce modèle linéaire simple}.
    \item \textbf{Limites du modèle et sécurité :} Le modèle linéaire ne tient pas aux fortes intensités (échauffement, non-linéarité chimique de la batterie). La tension aux bornes de la batterie ne doit pas dépasser la tension de fin de charge (14,4 V pour plomb 12 V). \textbf{Toute installation réelle impose un régulateur de charge} entre le panneau et la batterie, ce qui sort du cadre du cours mais en justifie la nécessité physique.
\end{enumerate}
\textbf{Compétence validée :} \og Analyser le fonctionnement énergétique d'un circuit électrique réel en modélisant ses éléments par des dipôles linéaires et en déterminant leur point de fonctionnement. \fg

\newpage

\coursec{Méthodes et exercices résolus}

\courssub{Méthodes et exercices résolus}

\begin{methode}[Schématiser un générateur réel et écrire sa loi de tension]
\textbf{Objectif :} Représenter un générateur réel par son modèle de Thévenin et écrire la relation entre la tension aux bornes $U$ et l'intensité $I$ traversante.
\begin{enumerate}
    \item \textbf{Identifier les bornes} du générateur (A et B) et le sens du courant \emph{reçu} par le générateur (convention récepteur : $I$ entre par la borne $+$).
    \item \textbf{Tracer le modèle équivalent :} Un générateur de tension idéal de f.e.m. $E$ (flèche à l'intérieur du symbole, du $-$ vers le $+$) en série avec une résistance interne $r$ (flèche de courant $I$ dans le même sens).
    \item \textbf{Ecrire la loi des mailles (Loi de Kirchhoff) :} En partant de la borne $A$ vers $B$ dans le sens du courant $I$ :
    \[ U_{AB} = E - r I \]
    \item \textbf{Vérifier l'homogénéité :} $[E] = \si{V}$, $[rI] = \si{\ohm}\cdot\si{A} = \si{V}$. Correct.
    \item \textbf{Noter les cas limites :} Circuit ouvert ($I=0 \Rightarrow U=E$) ; Court-circuit ($U=0 \Rightarrow I_{cc}=E/r$).
\end{enumerate}
\cle{Convention récepteur} : $I$ entre par la borne $+$ du dipôle. \cle{Modèle de Thévenin} : $E$ en série avec $r$.
\end{methode}

\begin{exoresolu}[Modélisation d'une pile alcaline usagée]
Une pile alcaline de référence \texttt{LR6} (format AA) neuve possède une f.e.m. $E_0 = \SI{1,50}{V}$ et une résistance interne $r_0 = \SI{0,15}{\ohm}$. Après une longue utilisation dans une lampe de poche, sa f.e.m. a chuté à $E = \SI{1,20}{V}$ et sa résistance interne a augmenté à $r = \SI{1,2}{\ohm}$.
\begin{enumerate}
    \item Schématiser le modèle équivalent de cette pile usagée en indiquant les polarités des bornes $A$ ($+$) et $B$ ($-$), le sens du courant $I$ (convention récepteur) et les valeurs de $E$ et $r$.
    \item Écrire la relation $U_{AB} = f(I)$.
    \item Calculer la tension aux bornes $U_{AB}$ lorsque la pile alimente une résistance de charge $R = \SI{10}{\ohm}$.
    \item Calculer l'intensité de court-circuit $I_{cc}$ et la puissance dissipée \emph{à l'intérieur} de la pile en court-circuit. Commenter le danger.
\end{enumerate}
\tcblower
\textbf{Solution détaillée}
\begin{enumerate}
    \item \textbf{Schéma du modèle de Thévenin :}
    \begin{popfigure}
    \begin{tikzpicture}[european, thick]
        \draw (0,0) node[anchor=east] {$A\ (+\!)$} to[short] (1,0)
            to[battery1, l_={$E=\SI{1,20}{V}$}] (3,0)
            to[R, l_={$r=\SI{1,2}{\ohm}$}] (5,0)
            to[short] (6,0) node[anchor=west] {$B\ (-\!)$};
        \draw[->, thick, popblue] (2.5,0.5) -- (2.5,1.2) node[above] {$I$};
        \draw[<->] (1,-0.8) -- node[below] {$U_{AB}$} (5,-0.8);
    \end{tikzpicture}
    \legende{Modèle de Thévenin de la pile usagée (convention récepteur).}
    \end{popfigure}
    \item \textbf{Loi de tension :} En appliquant la loi des mailles dans le sens du courant $I$ :
    \[ U_{AB} = E - r I = 1,20 - 1,2\,I \quad (U \text{ en V}, I \text{ en A}) \]
    \item \textbf{Fonctionnement en charge ($R=\SI{10}{\ohm}$) :}
    Le circuit est une maille unique : $E = r I + R I = (r+R)I$.
    \[ I = \frac{E}{r+R} = \frac{1,20}{1,2 + 10} = \frac{1,20}{11,2} \approx \SI{0,107}{A} \]
    Tension aux bornes (tension utile) :
    \[ U_{AB} = E - r I = 1,20 - 1,2 \times 0,1071 \approx 1,20 - 0,1286 = \SI{1,07}{V} \]
    (Vérification : $U_{AB} = R I = 10 \times 0,1071 = \SI{1,07}{V}$).
    \item \textbf{Court-circuit ($R=0, U_{AB}=0$) :}
    \[ I_{cc} = \frac{E}{r} = \frac{1,20}{1,2} = \SI{1,00}{A} \]
    Puissance dissipée \emph{internement} (effet Joule dans $r$) :
    \[ P_{r} = r I_{cc}^2 = 1,2 \times (1,00)^2 = \SI{1,20}{W} \]
    \textbf{Commentaire :} Une puissance de \SI{1,2}{W} dissipée dans le volume minuscule de la pile (quelques cm$^3$) entraîne un échauffement rapide, risque de fuite d'électrolyte, de gonflement ou d'explosion. \textbf{Ne jamais court-circuiter une pile, même usagée.}
\end{enumerate}
\cle{Conclusion} : La pile usagée délivre une tension utile de \SI{1,07}{V} sous \SI{10}{\ohm} (contre \SI{1,48}{V} pour une neuve), et son courant de court-circuit de \SI{1}{A} reste dangereux.
\end{exoresolu}

\begin{methode}[Exploiter la caractéristique $U=f(I)$ d'un générateur pour déterminer $E$ et $r$]
\textbf{Objectif :} À partir du tracé expérimental $U_{AB} = f(I)$ (droite descendante), extraire la f.e.m. $E$ et la résistance interne $r$.
\begin{enumerate}
    \item \textbf{Vérifier la linéarité :} Les points doivent être alignés. La relation est $U = E - r I$ (équation de droite $y = ax+b$).
    \item \textbf{Déterminer $E$ (ordonnée à l'origine) :} Lire la tension $U$ pour $I=0$ (circuit ouvert). $E = U(I=0)$.
    \item \textbf{Déterminer $r$ (pente) :} Deux méthodes rigoureuses :
        \begin{itemize}
            \item \textbf{Graphique :} Choisir deux points \emph{loin l'un de l'autre} sur la droite tracée $(I_1, U_1)$ et $(I_2, U_2)$. Calculer $r = -\frac{\Delta U}{\Delta I} = \frac{U_1 - U_2}{I_2 - I_1}$ (valeur positive).
            \item \textbf{Par l'abscisse à l'origine :} Lire $I_{cc}$ (court-circuit, $U=0$). $r = \frac{E}{I_{cc}}$.
        \end{itemize}
    \item \textbf{Unités et chiffres significatifs :} $E$ en \si{V}, $r$ en \si{\ohm}. Conserver 2 à 3 chiffres significatifs cohérents avec la précision graphique.
    \item \textbf{Rédiger :} « La caractéristique est une droite, le générateur est linéaire. $E = \dots\,\si{V}$, $r = \dots\,\si{\ohm}$. »
\end{enumerate}
\cle{Réflexe} : Toujours prendre les points sur la \emph{droite tracée}, pas sur les points de mesure bruts (qui ont des incertitudes).
\end{methode}

\begin{exoresolu}[Détermination de $E$ et $r$ d'un panneau solaire par sa caractéristique]
On souhaite caractériser un petit panneau photovoltaïque (type chargeur téléphone) sous un éclairage constant. On réalise le montage variateur + ampèremètre + voltmètre et on relève les couples $(I ; U)$ suivants :
\begin{center}
\begin{tabular}{|c|c|c|c|c|c|c|}\hline
$I$ (\si{mA}) & 0 & 10 & 25 & 40 & 55 & 65 \\ \hline
$U$ (\si{V}) & 5,80 & 5,20 & 4,20 & 3,10 & 2,00 & 0,90 \\ \hline
\end{tabular}
\end{center}
\begin{enumerate}
    \item Tracer la caractéristique $U = f(I)$ dans un repère orthonormé (échelle : 1 cm pour 10 mA en abscisses, 1 cm pour 0,5 V en ordonnées).
    \item Montrer que le panneau se comporte comme un générateur linéaire dans cette plage. Déterminer $E$ et $r$ par la méthode graphique (deux points sur la droite).
    \item Calculer la puissance maximale $P_{max}$ que ce panneau peut fournir à une charge adaptée (théorème de la puissance maximale : $R_{opt} = r$).
\end{enumerate}
\tcblower
\textbf{Solution détaillée}
\begin{enumerate}
    \item \textbf{Tracé graphique :}
    \begin{popfigure}
    \begin{tikzpicture}
        \begin{axis}[
            width=\linewidth, height=8cm,
            xlabel={$I$ (\si{mA})}, ylabel={$U$ (\si{V})},
            xmin=0, xmax=70, ymin=0, ymax=6.2,
            grid=major, minor tick num=1,
            axis lines=middle,
            tick label style={font=\small},
            label style={font=\small},
        ]
        \addplot[only marks, mark=*, mark size=2, popblue] coordinates {
            (0,5.80) (10,5.20) (25,4.20) (40,3.10) (55,2.00) (65,0.90)
        };
        % Droite de régression / tendance visuelle
        \addplot[poporange, thick, domain=0:70, samples=2] {-0.0753*x + 5.80};
        \end{axis}
    \end{tikzpicture}
    \legende{Caractéristique $U=f(I)$ du panneau solaire. Les points sont alignés.}
    \end{popfigure}
    Les points expérimentaux sont parfaitement alignés. Le panneau se comporte comme un \textbf{générateur linéaire} (modèle de Thévenin valide).
    \item \textbf{Détermination de $E$ et $r$ :}
    \textbf{Ordonnée à l'origine ($I=0$) :} $E = U(0) = \SI{5,80}{V}$.
    \textbf{Pente (deux points sur la droite tracée) :} On choisit l'origine $(0; 5,80)$ et le point à $I=60\,\si{mA}$ lu sur la droite : $U(60) \approx 1,28\,\si{V}$.
    \[ r = -\frac{\Delta U}{\Delta I} = \frac{5,80 - 1,28}{60 - 0} = \frac{4,52}{60} = 0,0753\,\si{V/mA} = \SI{75,3}{\ohm} \]
    \textit{Vérification par $I_{cc}$ :} Abscisse à l'origine $\approx 77\,\si{mA}$. $r = E/I_{cc} = 5,80 / 0,077 \approx \SI{75,3}{\ohm}$. Cohérent.
    \item \textbf{Puissance maximale :}
    Au point de fonctionnement optimal (adaptation d'impédance), $R_{charge} = r = \SI{75,3}{\ohm}$.
    Courant optimal : $I_{opt} = \frac{E}{2r} = \frac{5,80}{2 \times 75,3} = \frac{5,80}{150,6} \approx \SI{38,5}{mA}$.
    Tension aux bornes : $U_{opt} = \frac{E}{2} = \SI{2,90}{V}$.
    Puissance maximale :
    \[ P_{max} = U_{opt} I_{opt} = \frac{E^2}{4r} = \frac{(5,80)^2}{4 \times 75,3} = \frac{33,64}{301,2} \approx \SI{0,112}{W} = \SI{112}{mW} \]
\end{enumerate}
\cle{Conclusion} : Le panneau équivaut à un générateur $E=\SI{5,80}{V}$, $r=\SI{75,3}{\ohm}$. Sa puissance max est \SI{112}{mW} pour $R=\SI{75,3}{\ohm}$ ($U=\SI{2,90}{V}, I=\SI{38,5}{mA}$).
\end{exoresolu}

\begin{methode}[Déterminer les caractéristiques $E_{eq}$ et $r_{eq}$ d'une association de générateurs]
\textbf{Objectif :} Remplacer une association de générateurs réels par un unique générateur de Thévenin équivalent ($E_{eq}, r_{eq}$).
\begin{enumerate}
    \item \textbf{Identifier le type d'association :}
        \begin{itemize}
            \item \textbf{Série :} Bornes $+$ du 1er reliée au $-$ du 2nd (courant $I$ identique). Bornes libres = bornes de l'association.
            \item \textbf{Parallèle :} Bornes $+$ ensemble, bornes $-$ ensemble (tension $U$ identique aux bornes de l'association).
            \item \textbf{Association mixte :} Réduire par étapes (série puis parallèle ou inverse).
        \end{itemize}
    \item \textbf{Appliquer les formules de groupement (démonstration exigible au Probatoire) :}
        \begin{itemize}
            \item \textbf{Série :} $E_{eq} = \sum E_i$ ; $r_{eq} = \sum r_i$.
            \item \textbf{Parallèle (générateurs de même f.e.m. $E$) :} $E_{eq} = E$ ; $\frac{1}{r_{eq}} = \sum \frac{1}{r_i}$.
            \item \textbf{Parallèle (f.e.m. différentes $E_1 \neq E_2$) :} \textbf{Interdit en régime permanent continu} (courants de boucle destructifs). Si imposé (ex: charge de batterie), utiliser les lois de Kirchhoff maille par maille, pas de formule simple.
        \end{itemize}
    \item \textbf{Vérifier les unités :} $E_{eq}$ en \si{V}, $r_{eq}$ en \si{\ohm}.
    \item \textbf{Schématiser l'équivalent :} Un seul générateur $E_{eq}$ en série avec $r_{eq}$.
\end{enumerate}
\cle{Attention} : Parallèle de f.e.m. différentes $\Rightarrow$ court-circuit interne. Formules série/parallèle valides \textbf{uniquement} pour générateurs linéaires (modèle Thévenin).
\end{methode}

\begin{exoresolu}[Association série-parallèle de piles pour motoréducteur]
Un motoréducteur de voiture télécommandée (type \emph{brushless}) nécessite une tension nominale de \SI{7,2}{V} et un courant de démarrage de \SI{2,5}{A}. On dispose de piles NiMH rechargeables (format AA) de caractéristiques : $E_0 = \SI{1,20}{V}$, $r_0 = \SI{0,050}{\ohm}$.
On réalise un pack « 6S2P » : 6 piles en série (6S), et 2 branches identiques en parallèle (2P).
\begin{enumerate}
    \item Schématiser l'association complète (12 piles) en indiquant les polarités.
    \item Déterminer la f.e.m. équivalente $E_{eq}$ et la résistance interne équivalente $r_{eq}$ du pack.
    \item Calculer la tension aux bornes du pack $U_{pack}$ au démarrage du moteur ($I = \SI{2,5}{A}$).
    \item Le moteur démarre-t-il correctement ? (Tension minimale de fonctionnement : \SI{6,0}{V}).
\end{enumerate}
\tcblower
\textbf{Solution détaillée}
\begin{enumerate}
    \item \textbf{Schéma :} 2 branches parallèles. Chaque branche = 6 piles en série.
    \begin{popfigure}
    \begin{tikzpicture}[european, thick, scale=0.8]
        % Branche haute
        \draw (0,2) node[left] {$+$} to[short] (1,2)
            foreach \x in {1,...,6} {
                to[battery1, l_={$E_0, r_0$}] ++(1.5,0)
            } node[right] (A) {};
        \draw (10,2) -- (11,2) node[right] {$A$};
        % Branche basse
        \draw (0,-2) node[left] {$-$} to[short] (1,-2)
            foreach \x in {1,...,6} {
                to[battery1, l_={$E_0, r_0$}] ++(1.5,0)
            } node[right] (B) {};
        \draw (10,-2) -- (11,-2) node[right] {$B$};
        % Connexions
        \draw (0,2) -- (0,-2);
        \draw (10,2) -- (10,-2);
    \end{tikzpicture}
    \legende{Pack 6S2P : 2 branches parallèles de 6 piles série.}
    \end{popfigure}
    \item \textbf{Calcul équivalent :}
    \textbf{Étape 1 : Une branche série (6 piles).}
    $E_{branche} = 6 \times E_0 = 6 \times 1,20 = \SI{7,20}{V}$.
    $r_{branche} = 6 \times r_0 = 6 \times 0,050 = \SI{0,300}{\ohm}$.
    \textbf{Étape 2 : Parallèle de 2 branches identiques ($E_{branche}$ same).}
    $E_{eq} = E_{branche} = \SI{7,20}{V}$.
    $\frac{1}{r_{eq}} = \frac{1}{r_{branche}} + \frac{1}{r_{branche}} = \frac{2}{0,300} \Rightarrow r_{eq} = \frac{0,300}{2} = \SI{0,150}{\ohm}$.
    \item \textbf{Tension aux bornes sous $I = \SI{2,5}{A}$ :}
    $U_{pack} = E_{eq} - r_{eq} I = 7,20 - 0,150 \times 2,5 = 7,20 - 0,375 = \SI{6,825}{V} \approx \SI{6,83}{V}$.
    \item \textbf{Analyse :} $U_{pack} = \SI{6,83}{V} > \SI{6,0}{V}$ (seuil min).
    \textbf{Oui, le moteur démarre correctement.} La chute de tension interne ($0,375\,\si{V}$) est acceptable grâce au montage parallèle qui a divisé $r$ par 2.
\end{enumerate}
\cle{Conclusion} : Le pack 6S2P fournit $E_{eq}=\SI{7,20}{V}$, $r_{eq}=\SI{0,150}{\ohm}$. Sous \SI{2,5}{A}, $U=\SI{6,83}{V} > \SI{6,0}{V}$. Le choix 2P (au lieu de 1P qui donnerait $r=0,30\,\si{\ohm}$ et $U=6,45\,\si{V}$) améliore la tenue en charge.
\end{exoresolu}

\begin{methode}[Réaliser le montage et tracer la caractéristique $U=f(I)$ d'un générateur]
\textbf{Objectif :} Monter le circuit expérimental, choisir les appareils de mesure, varier le courant, relever les valeurs, tracer le graphique.
\begin{enumerate}
    \item \textbf{Montage (Schéma) :}
        \begin{itemize}
            \item Générateur $G$ à caractériser.
            \item Rhéostat (résistance variable) $R_v$ en série pour faire varier $I$.
            \item Ampèremètre $A$ \textbf{en série} (calibre adapté, borne $+$ vers $+$ générateur).
            \item Voltmètre $V$ \textbf{en parallèle} aux bornes du générateur (bornes $+$ sur $+$ générateur).
            \item \textbf{Interrupteur} pour couper entre mesures (éviter décharge/échauffement).
        \end{itemize}
    \item \textbf{Choix des calibres :} Estimer $E$ et $I_{max} \approx E/r$. Choisir calibre $I$ juste au-dessus de $I_{max}$, calibre $U$ juste au-dessus de $E$.
    \item \textbf{Protocole :}
        \begin{enumerate}
            \item Régler $R_v$ au maximum (courant min, $I \approx 0$) $\rightarrow$ relever $U_{max} \approx E$.
            \item Diminuer $R_v$ par paliers réguliers (ex: 5 à 10 points) jusqu'au minimum (court-circuit quasi $U \approx 0$).
            \item Noter $(I_k ; U_k)$ à chaque palier. Ouvrir l'interrupteur entre chaque mesure.
        \end{enumerate}
    \item \textbf{Tracé graphique :} Repère orthonormé. $I$ en abscisses, $U$ en ordonnées. Échelle adaptée (origine $(0,0)$ si possible). Nuage de points $\rightarrow$ droite de régression (règle, pas « point à point »).
    \item \textbf{Exploitation :} Voir Méthode 2 (détermination $E, r$).
\end{enumerate}
\cle{Sécurité} : Ne pas maintenir le circuit fermé en court-circuit ($R_v=0$) plus que le temps de la lecture. Risque d'échauffement du générateur et du rhéostat.
\end{methode}

\begin{exoresolu}[Montage et tracé pour une pile bouton au lithium (CR2032)]
On veut tracer la caractéristique d'une pile bouton CR2032 ($E \approx \SI{3}{V}, r \approx \SI{15}{\ohm}$) au laboratoire.
\begin{enumerate}
    \item Proposer un schéma de montage complet avec les appareils de mesure (ampèremètre, voltmètre, rhéostat, interrupteur). Justifier le choix des calibres pour l'ampèremètre ($I_{max}$) et le voltmètre.
    \item On obtient le tableau de mesures suivant :
    \begin{center}
    \begin{tabular}{|c|c|c|c|c|c|c|}\hline
    $I$ (\si{mA}) & 0,0 & 5,0 & 10,0 & 15,0 & 20,0 & 22,0 \\ \hline
    $U$ (\si{V}) & 3,00 & 2,25 & 1,50 & 0,75 & 0,10 & 0,00 \\ \hline
    \end{tabular}
    \end{center}
    Tracer la caractéristique sur papier millimétré (échelle proposée : 1 cm $\rightarrow$ 5 mA ; 1 cm $\rightarrow$ 0,5 V). Déterminer $E$ et $r$ graphiquement.
    \item La pile est-elle encore bonne pour alimenter une carte mère CMOS (consommation $\approx \SI{5}{\micro A}$, tension min \SI{2,5}{V}) ?
\end{enumerate}
\tcblower
\textbf{Solution détaillée}
\begin{enumerate}
    \item \textbf{Schéma de montage :}
    \begin{popfigure}
    \begin{tikzpicture}[european, thick]
        \draw (0,0) node[anchor=east] {$+$} to[short] (1,0)
            to[battery1, l_=CR2032] (3,0)
            to[short] (4,0)
            to[ammeter, l={$A$}] (6,0)
            to[short] (7,0)
            to[R, l_={$R_v$}] (9,0)
            to[nos, l={$K$}] (10,0) node[anchor=west] {$-$};
        % Voltmètre en parallèle sur la pile
        \draw (1,1) to[voltmeter, l={$V$}] (3,1);
        \draw (1,0) -- (1,1);
        \draw (3,0) -- (3,1);
    \end{tikzpicture}
    \legende{Montage de caractérisation : $A$ en série, $V$ aux bornes du générateur, $K$ ouvert au repos.}
    \end{popfigure}
    \textbf{Choix des calibres :}
    $I_{cc} \approx E/r = 3,0 / 15 = \SI{0,20}{A} = \SI{200}{mA}$.
    $\rightarrow$ Ampèremètre : calibre \textbf{200 mA} (ou 500 mA si 200 indisponible).
    $E \approx \SI{3}{V}$.
    $\rightarrow$ Voltmètre : calibre \textbf{10 V} (ou 20 V).
    \item \textbf{Tracé et détermination graphique :}
    \begin{popfigure}
    \begin{tikzpicture}
        \begin{axis}[
            width=\linewidth, height=8cm,
            xlabel={$I$ (\si{mA})}, ylabel={$U$ (\si{V})},
            xmin=0, xmax=25, ymin=0, ymax=3.2,
            grid=major, axis lines=middle,
            tick label style={font=\small}, label style={font=\small},
        ]
        \addplot[only marks, mark=*, mark size=2.5, popblue] coordinates {
            (0,3.00) (5,2.25) (10,1.50) (15,0.75) (20,0.10) (22,0.00)
        };
        % Droite de régression passant par (0,3) et (22,0) -> pente -3/22 = -0.136 V/mA
        \addplot[poporange, thick, domain=0:22, samples=2] {3 - 0.136*x};
        \end{axis}
    \end{tikzpicture}
    \legende{Caractéristique $U=f(I)$ de la pile CR2032. Droite de régression tracée (modèle linéaire global).}
    \end{popfigure}
    Les points sont globalement alignés pour les faibles courants ; une courbure apparaît à fort courant (polarisation). On trace la droite de régression passant par l'origine $(0;3,00)$ et le point de court-circuit $(22;0)$.
    \textbf{Lecture graphique :}
    $E = U(I=0) = \SI{3,00}{V}$.
    $I_{cc} \approx \SI{22,0}{mA}$ (abscisse à l'origine).
    $r = \frac{E}{I_{cc}} = \frac{3,00}{22,0 \times 10^{-3}} = \SI{136}{\ohm}$.
    \textit{Note :} La valeur théorique attendue pour une pile neuve était $\sim \SI{15}{\ohm}$. Ici $r=\SI{136}{\ohm}$ indique une pile \textbf{très usagée} (résistance interne multipliée par 9).
    \item \textbf{Utilisation CMOS ($I=\SI{5}{\micro A}$) :}
    $U = E - r I = 3,00 - 136 \times 5 \times 10^{-6} = 3,00 - 0,00068 = \SI{2,999}{V} \approx \SI{3,00}{V}$.
    $U > \SI{2,5}{V}$. \textbf{Oui, la pile convient} pour cette application très faible courant, malgré sa résistance interne élevée.
\end{enumerate}
\cle{Conclusion} : Caractéristique tracée, $E=\SI{3,00}{V}$, $r=\SI{136}{\ohm}$. Pile inutilisable pour fort courant, mais OK pour CMOS.
\end{exoresolu}

\begin{methode}[Schématiser un récepteur actif et exploiter sa caractéristique $U=f(I)$]
\textbf{Objectif :} Modéliser un dipôle récepteur (moteur, batterie en charge, LED...) par un modèle de Thévenin « récepteur » : une f.e.m. de contre-électromotrice $E'$ (ou $E_{rec}$) en série avec une résistance $r'$, et déterminer $E', r'$ à partir de la caractéristique $U=f(I)$.
\begin{enumerate}
    \item \textbf{Convention récepteur :} $I$ entre par la borne $+$ du dipôle. $U = U_{AB} = V_A - V_B$.
    \item \textbf{Modèle équivalent (Générateur « inverse ») :} Le récepteur actif s'oppose au courant. Schéma : Une source de tension $E'$ (flèche du $-$ vers le $+$ \emph{à l'intérieur} du dipôle, c'est-à-dire dans le sens $A \to B$ si $A$ est $+$) en série avec $r'$. Le courant $I$ traverse $r'$ puis $E'$ de $+$ vers $-$ à l'intérieur de $E'$.
    \item \textbf{Loi de tension (Loi des mailles) :}
    \[ U = E' + r' I \]
    (Contrairement au générateur : $U = E - rI$). La droite a une \textbf{pente positive} $r'$.
    \item \textbf{Exploitation graphique :}
        \begin{itemize}
            \item $E' = U(I=0)$ (ordonnée à l'origine). C'est la tension de « seuil » ou f.e.m. de contre-électromotrice.
            \item $r' = \frac{\Delta U}{\Delta I}$ (pente positive). Choisir deux points sur la droite tracée.
        \end{itemize}
    \item \textbf{Types de récepteurs :}
        \begin{itemize}
            \item \textbf{Passif (ohmique) :} $E'=0$, $U=r'I$ (droite passant par l'origine).
            \item \textbf{Actif (moteur, batterie en charge, pile en charge, électrolyse) :} $E' > 0$.
        \end{itemize}
\end{enumerate}
\cle{Différence clé} : Générateur : $U = E - rI$ (pente $-r$). Récepteur actif : $U = E' + r'I$ (pente $+r'$).
\end{methode}

\begin{exoresolu}[Caractéristique d'un moteur à courant continu en charge]
Un petit moteur CC (type \emph{drone} ou \emph{voiture RC}) est alimenté par une source de tension variable. On mesure la tension $U$ à ses bornes en fonction du courant $I$ absorbé (convention récepteur). Résultats :
\begin{center}
\begin{tabular}{|c|c|c|c|c|c|}\hline
$I$ (\si{A}) & 0,0 & 0,5 & 1,0 & 1,5 & 2,0 \\ \hline
$U$ (\si{V}) & 1,80 & 2,10 & 2,40 & 2,70 & 3,00 \\ \hline
\end{tabular}
\end{center}
\begin{enumerate}
    \item Tracer la caractéristique $U=f(I)$. Quelle est la nature de ce récepteur ? Justifier.
    \item Déterminer $E'$ (f.e.m. de contre-électromotrice à vide) et $r'$ (résistance interne du bobinage).
    \item Le moteur est alimenté sous $U = \SI{5,0}{V}$. Calculer l'intensité $I$ absorbée et la puissance mécanique utile $P_{mec}$ développée.
    \item Calculer le rendement $\eta$ à ce point de fonctionnement.
\end{enumerate}
\tcblower
\textbf{Solution détaillée}
\begin{enumerate}
    \item \textbf{Tracé et nature :}
    \begin{popfigure}
    \begin{tikzpicture}
        \begin{axis}[
            width=\linewidth, height=8cm,
            xlabel={$I$ (\si{A})}, ylabel={$U$ (\si{V})},
            xmin=0, xmax=2.2, ymin=1.5, ymax=3.2,
            grid=major, axis lines=middle,
            tick label style={font=\small}, label style={font=\small},
        ]
        \addplot[only marks, mark=*, mark size=2.5, popblue] coordinates {
            (0,1.80) (0.5,2.10) (1.0,2.40) (1.5,2.70) (2.0,3.00)
        };
        \addplot[poporange, thick, domain=0:2.2, samples=2] {1.8 + 0.6*x};
        \end{axis}
    \end{tikzpicture}
    \legende{Caractéristique $U=f(I)$ du moteur CC. Droite à pente positive.}
    \end{popfigure}
    Les points sont alignés sur une droite à \textbf{pente positive}. $U(I=0) = \SI{1,80}{V} \neq 0$.
    $\rightarrow$ C'est un \textbf{récepteur actif} (modèle : $E'$ en série avec $r'$).
    \item \textbf{Paramètres du modèle :}
    $E' = U(I=0) = \SI{1,80}{V}$.
    Pente $r' = \frac{\Delta U}{\Delta I} = \frac{3,00 - 1,80}{2,0 - 0,0} = \frac{1,20}{2,0} = \SI{0,60}{\ohm}$.
    Modèle : $U = 1,80 + 0,60\,I$.
    \item \textbf{Fonctionnement sous $U = \SI{5,0}{V}$ :}
    $U = E' + r' I \Rightarrow 5,0 = 1,80 + 0,60 I$
    $0,60 I = 3,20 \Rightarrow I = \frac{3,20}{0,60} = \SI{5,33}{A}$.
    Puissance mécanique (puissance développée par la f.e.m. $E'$) :
    $P_{mec} = E' I = 1,80 \times 5,33 = \SI{9,60}{W}$.
    \item \textbf{Rendement :}
    Puissance absorbée (électrique) : $P_{abs} = U I = 5,0 \times 5,33 = \SI{26,67}{W}$.
    Pertes Joule dans $r'$ : $P_J = r' I^2 = 0,60 \times (5,33)^2 = \SI{17,07}{W}$.
    Vérification : $P_{abs} = P_{mec} + P_J = 9,60 + 17,07 = 26,67\,\si{W}$. OK.
    Rendement : $\eta = \frac{P_{mec}}{P_{abs}} = \frac{9,60}{26,67} = 0,360 = \textbf{36,0 \%}$.
    (Ou $\eta = \frac{E'}{U} = \frac{1,80}{5,0} = 0,36$).
\end{enumerate}
\cle{Conclusion} : Moteur actif : $E'=\SI{1,80}{V}, r'=\SI{0,60}{\ohm}$. Sous \SI{5}{V} : $I=\SI{5,33}{A}, P_{mec}=\SI{9,60}{W}, \eta=36\%$. Rendement faible car $E' \ll U$ (beaucoup de pertes Joule).
\end{exoresolu}

\begin{methode}[Réaliser le montage et tracer la caractéristique $U=f(I)$ d'un récepteur]
\textbf{Objectif :} Monter le circuit pour caractériser un récepteur (passif ou actif), relever $(I, U)$, tracer.
\begin{enumerate}
    \item \textbf{Montage :}
        \begin{itemize}
            \item Générateur de tension variable $G$ (alimentation stabilisée ou pile + rhéostat en série) pour imposer $U$ ou $I$.
            \item Récepteur $D$ à caractériser.
            \item Ampèremètre $A$ \textbf{en série} avec $D$.
            \item Voltmètre $V$ \textbf{en parallèle} aux bornes de $D$ (bornes $+$ sur $+$ de $D$).
            \item Interrupteur $K$.
        \end{itemize}
    \item \textbf{Deux stratégies de variation :}
        \begin{itemize}
            \item \textbf{Imposer $U$ (source de tension idéale) :} Varier $U_G$, lire $I$ et $U$. Facile pour récepteurs passifs.
            \item \textbf{Imposer $I$ (source de courant / générateur réel + rhéostat série) :} Varier $R_v$, lire $I$ et $U$. \textbf{Indispensable pour récepteurs actifs} (moteur, batterie en charge) car leur caractéristique a une pente positive : une source de tension idéale imposerait $U$, mais le point de fonctionnement serait instable ou inaccessible sans résistance série.
        \end{itemize}
    \item \textbf{Protocole :} Varier le paramètre de commande ($U_G$ ou $R_v$) par paliers. Relever $(I_k, U_k)$. Ouvrir $K$ entre mesures.
    \item \textbf{Tracé :} $I$ abscisses, $U$ ordonnées. Droite de régression.
    \item \textbf{Exploitation :} $E' = U(I=0)$ (extrapolation si $I=0$ non atteint). $r' = \frac{\Delta U}{\Delta I}$ (pente).
\end{enumerate}
\cle{Astuce} : Pour un récepteur actif (moteur), utiliser un générateur réel (pile) + rhéostat en série permet de balayer la caractéristique de $I_{max}$ (court-circuit sur le générateur) vers $I=0$ (moteur à vide).
\end{methode}

\begin{exoresolu}[Caractérisation d'une batterie NiMH en charge (récepteur actif)]
On charge une batterie NiMH ($E_{bat} \approx \SI{1,2}{V}, r_{bat} \approx \SI{0,05}{\ohm}$) à l'aide d'une alimentation de laboratoire $U_G = \SI{5,0}{V}$ et d'une résistance de limitation $R = \SI{10}{\ohm}$ en série. On mesure $U_{bat}$ et $I_{charge}$ pour différentes valeurs de $U_G$ (ou en agissant sur $R$).
On obtient les points : $(I=0,0\,\si{A}; U=1,25\,\si{V})$, $(0,1\,\si{A}; 1,30\,\si{V})$, $(0,2\,\si{A}; 1,35\,\si{V})$, $(0,3\,\si{A}; 1,40\,\si{V})$, $(0,35\,\si{A}; 1,42\,\si{V})$.
\begin{enumerate}
    \item Schématiser le montage de caractérisation (Alim $U_G$, $R$, $A$, $V$, Batterie).
    \item Tracer $U_{bat}=f(I_{charge})$. Déterminer $E'$ et $r'$ de la batterie \emph{en charge}.
    \item Interpréter physiquement $E'$ et $r'$ ici. Pourquoi $E' > E_{nominal}$ ?
    \item Si on veut charger à $I = \SI{0,5}{A}$ (courant nominal), quelle tension $U_G$ faut-il imposer avec $R=\SI{10}{\ohm}$ ?
\end{enumerate}
\tcblower
\textbf{Solution détaillée}
\begin{enumerate}
    \item \textbf{Montage :}
    \begin{popfigure}
    \begin{tikzpicture}[european, thick]
        \draw (0,0) node[anchor=east] {$U_G$} to[sV, l_={$U_G$}] (0,-2) -- (2,-2) node[anchor=north] {$B(-)$};
        \draw (0,0) to[short] (2,0) node[anchor=south] {$A(+)$};
        \draw (2,0) to[ammeter, l={$A$}] (4,0)
            to[R, l={$R=\SI{10}{\ohm}$}] (6,0)
            to[battery1, l_={$E_{bat}, r_{bat}$}] (6,-2)
            to[short] (4,-2)
            to[voltmeter, l={$V$}] (2,-2);
        \draw (4,0) -- (4,0.5) node[above] {$I_{charge}$};
    \end{tikzpicture}
    \legende{Montage charge batterie : $U_G$ + $R$ limite le courant. $V$ aux bornes batterie.}
    \end{popfigure}
    \item \textbf{Tracé et paramètres :}
    \begin{popfigure}
    \begin{tikzpicture}
        \begin{axis}[
            width=\linewidth, height=8cm,
            xlabel={$I$ (\si{A})}, ylabel={$U$ (\si{V})},
            xmin=0, xmax=0.4, ymin=1.2, ymax=1.45,
            grid=major, axis lines=middle,
            tick label style={font=\small}, label style={font=\small},
        ]
        \addplot[only marks, mark=*, mark size=2.5, popblue] coordinates {
            (0,1.25) (0.1,1.30) (0.2,1.35) (0.3,1.40) (0.35,1.42)
        };
        \addplot[poporange, thick, domain=0:0.4, samples=2] {1.25 + 0.5*x};
        \end{axis}
    \end{tikzpicture}
    \legende{Caractéristique charge batterie. Droite $U = E' + r'I$.}
    \end{popfigure}
    Droite tracée. $E' = U(I=0) = \SI{1,25}{V}$.
    Pente $r' = \frac{1,42 - 1,25}{0,35 - 0} = \frac{0,17}{0,35} \approx \SI{0,486}{\ohm}$.
    \item \textbf{Interprétation :}
    $E' = \SI{1,25}{V}$ est la \textbf{tension de charge} (f.e.m. de contre-électromotrice chimique) à vide. Elle est $> E_{nominal}=\SI{1,2}{V}$ car la réaction chimique inverse (charge) impose un potentiel plus élevé que l'équilibre thermodynamique (sur-tension de charge, polarisation).
    $r' = \SI{0,486}{\ohm}$ est la \textbf{résistance interne apparente en charge}. Elle est $> r_{bat}=\SI{0,05}{\ohm}$ car elle inclut les effets de polarisation de concentration et d'activation (cinétique électrochimique) qui s'opposent au passage du courant.
    \item \textbf{Calcul $U_G$ pour $I=\SI{0,5}{A}$ :}
    Tension aux bornes batterie nécessaire : $U_{bat} = E' + r' I = 1,25 + 0,486 \times 0,5 = 1,25 + 0,243 = \SI{1,493}{V}$.
    Chute de tension sur $R$ : $U_R = R I = 10 \times 0,5 = \SI{5,0}{V}$.
    Tension générateur requise : $U_G = U_{bat} + U_R = 1,493 + 5,0 = \SI{6,49}{V}$.
    \textbf{Problème :} $U_G = \SI{6,49}{V} > \SI{5,0}{V}$ (alim max). \textbf{Impossible avec $R=\SI{10}{\ohm}$ et $U_G=\SI{5}{V}$.}
    Il faut baisser $R$ : $R_{new} = \frac{U_G - U_{bat}}{I} = \frac{5,0 - 1,493}{0,5} = \frac{3,507}{0,5} = \SI{7,01}{\ohm}$.
\end{enumerate}
\cle{Conclusion} : Batterie en charge : $E'=\SI{1,25}{V}, r'=\SI{0,49}{\ohm}$. Pour charger à \SI{0,5}{A} avec alim \SI{5}{V}, il faut $R \approx \SI{7}{\ohm}$ (au lieu de \SI{10}{\ohm}).
\end{exoresolu}

\begin{methode}[Déterminer le point de fonctionnement d'un circuit Générateur–Récepteur]
\textbf{Objectif :} Trouver le couple $(I_{PF}, U_{PF})$ commun au générateur et au récepteur connectés ensemble.
\begin{enumerate}
    \item \textbf{Écrire les deux équations caractéristiques :}
        \begin{align*}
            \text{Générateur :} \quad & U = E - r I \\
            \text{Récepteur :} \quad & U = E' + r' I \quad (\text{si actif}) \quad \text{ou} \quad U = R I \quad (\text{si passif})
        \end{align*}
    \item \textbf{Méthode 1 : Analytique (calcul exact).} Égaliser les deux expressions de $U$ :
    \[ E - r I = E' + r' I \quad \Rightarrow \quad I_{PF} = \frac{E - E'}{r + r'} \]
    Puis $U_{PF} = E - r I_{PF}$ (ou $E' + r' I_{PF}$).
    \textbf{Condition d'existence :} $E > E'$ (sinon $I<0$, le « générateur » se fait charger).
    \item \textbf{Méthode 2 : Graphique (construction).} Tracer les deux droites $U_G(I)$ et $U_R(I)$ dans le \textbf{même repère} ($I$ abscisses, $U$ ordonnées). Le point d'intersection est le point de fonctionnement.
    \item \textbf{Vérifications obligatoires :}
        \begin{itemize}
            \item $I_{PF} \ge 0$ (convention récepteur respectée).
            \item $U_{PF} \ge 0$.
            \item Puissances cohérentes : $P_{fournie} = E I_{PF}$, $P_{absorbée} = U_{PF} I_{PF}$, $P_{perdues} = (r+r') I_{PF}^2$. $P_{fournie} = P_{absorbée} + P_{perdues}$.
        \end{itemize}
    \item \textbf{Rédiger la conclusion :} « Le point de fonctionnement est $I_{PF} = \dots\,\si{A}$, $U_{PF} = \dots\,\si{V}$. »
\end{enumerate}
\cle{Réflexe Probatoire} : Toujours vérifier $E > E'$ avant de calculer. Si $E < E'$, le générateur ne peut pas imposer le courant dans ce sens ; le système s'inverse (charge).
\end{methode}

\begin{exoresolu}[Point de fonctionnement : Panneau solaire + Moteur CC]
On connecte directement (sans électronique de puissance) le panneau solaire de l'exercice 2 ($E_p = \SI{5,80}{V}, r_p = \SI{75,3}{\ohm}$) au moteur CC de l'exercice 5 ($E_m = \SI{1,80}{V}, r_m = \SI{0,60}{\ohm}$).
\begin{enumerate}
    \item Schématiser le circuit complet (générateur + récepteur). Indiquer le sens du courant $I$ et la tension $U$ aux bornes du moteur.
    \item Déterminer analytiquement le point de fonctionnement $(I_{PF}, U_{PF})$.
    \item Calculer la puissance mécanique $P_{mec}$ fournie par le moteur et le rendement global de la chaîne (puissance mécanique / puissance solaire incidente supposée \SI{500}{mW}).
    \item Tracer graphiquement les deux caractéristiques dans le même repère et montrer le point de fonctionnement.
    \item Le moteur tourne-t-il à sa vitesse nominale ? (Rappel : $U_{nom} = \SI{5,0}{V}$ pour ce moteur).
\end{enumerate}
\tcblower
\textbf{Solution détaillée}
\begin{enumerate}
    \item \textbf{Schéma :}
    \begin{popfigure}
    \begin{tikzpicture}[european, thick]
        \draw (0,0) node[anchor=east] {$+$} to[short] (1,0)
            to[battery1, l_={$E_p=\SI{5,80}{V}$}] (3,0)
            to[R, l={$r_p=\SI{75,3}{\ohm}$}] (5,0)
            to[short] (6,0) node[above] {A};
        \draw (6,0) to[R, l={$r_m=\SI{0,60}{\ohm}$}] (8,0)
            to[battery1, l_={$E_m=\SI{1,80}{V}$}] (10,0)
            to[short] (11,0) node[above] {B};
        \draw (11,0) -- (11,-1) -- (0,-1) -- (0,0);
        \draw[->, thick, popblue] (5.5,0.3) -- (5.5,0.8) node[above] {$I$};
        \draw[<->] (6,-0.5) -- node[below] {$U_{moteur}$} (10,-0.5);
    \end{tikzpicture}
    \legende{Circuit Panneau (G) -- Moteur (R). Convention récepteur sur le moteur.}
    \end{popfigure}
    \item \textbf{Calcul analytique :}
    Condition : $E_p = 5,80 > E_m = 1,80$. OK.
    \[ I_{PF} = \frac{E_p - E_m}{r_p + r_m} = \frac{5,80 - 1,80}{75,3 + 0,60} = \frac{4,00}{75,9} = \SI{0,0527}{A} = \SI{52,7}{mA} \]
    \[ U_{PF} = E_m + r_m I_{PF} = 1,80 + 0,60 \times 0,0527 = 1,80 + 0,0316 = \SI{1,83}{V} \]
    (Vérif : $U_{PF} = E_p - r_p I_{PF} = 5,80 - 75,3 \times 0,0527 = 5,80 - 3,97 = \SI{1,83}{V}$).
    \item \textbf{Puissances et rendement :}
    Puissance mécanique : $P_{mec} = E_m I_{PF} = 1,80 \times 0,0527 = \SI{0,0949}{W} = \SI{94,9}{mW}$.
    Puissance fournie par panneau : $P_{pan} = E_p I_{PF} = 5,80 \times 0,0527 = \SI{0,306}{W} = \SI{306}{mW}$.
    Pertes totales (Joule) : $P_J = (r_p + r_m) I_{PF}^2 = 75,9 \times (0,0527)^2 = \SI{0,211}{W}$.
    Bilan : $306 = 94,9 + 211 \approx 306$. OK.
    Rendement global (électrique $\to$ mécanique) : $\eta_{glob} = \frac{P_{mec}}{P_{pan}} = \frac{94,9}{306} = 0,310 = \textbf{31,0 \%}$.
    Rendement par rapport au soleil ($P_{sun}=\SI{500}{mW}$) : $\eta_{tot} = \frac{94,9}{500} = \textbf{19,0 \%}$.
    \item \textbf{Tracé graphique :}
    \begin{popfigure}
    \begin{tikzpicture}
        \begin{axis}[
            width=\linewidth, height=9cm,
            xlabel={$I$ (\si{mA})}, ylabel={$U$ (\si{V})},
            xmin=0, xmax=80, ymin=0, ymax=6.5,
            grid=major, axis lines=middle,
            tick label style={font=\small}, label style={font=\small},
            legend pos=north east,
        ]
        % Panneau : U = 5.8 - 0.0753 I (I en mA)
        \addplot[popblue, thick, domain=0:80, samples=2] {5.8 - 0.0753*x};
        \addlegendentry{Panneau : $U = 5,80 - 0,0753 I$}
        % Moteur : U = 1.8 + 0.0006 I (I en mA -> 0.6 ohm = 0.0006 V/mA)
        \addplot[poporange, thick, domain=0:80, samples=2] {1.8 + 0.0006*x};
        \addlegendentry{Moteur : $U = 1,80 + 0,0006 I$}
        % Point intersection
        \addplot[only marks, mark=*, mark size=3, popgreen] coordinates {(52.7, 1.83)};
        \addlegendentry{Point de fonctionnement}
        \end{axis}
    \end{tikzpicture}
    \legende{Intersection des caractéristiques : Point de fonctionnement $I_{PF}=52,7\,\si{mA}, U_{PF}=1,83\,\si{V}$.}
    \end{popfigure}
    \item \textbf{Analyse vitesse :} Tension aux bornes moteur $U_{PF} = \SI{1,83}{V} \ll U_{nom} = \SI{5,0}{V}$.
    Le moteur tourne \textbf{très lentement} (vitesse $\propto U - E' \approx 0,03\,\si{V}$). Le panneau est \textbf{sous-dimensionné} (résistance interne trop élevée $75\,\si{\ohm}$) par rapport à la charge moteur ($r' \approx 0,6\,\si{\ohm}$). Il faut un convertisseur (MPPT / hacheur) pour adapter l'impédance.
\end{enumerate}
\cle{Conclusion} : $I_{PF}=\SI{52,7}{mA}, U_{PF}=\SI{1,83}{V}$. Moteur sous-alimenté ($U \ll U_{nom}$). Adaptation d'impédance nécessaire.
\end{exoresolu}

\begin{attention}[Les 5 erreurs qui coûtent des points au Probatoire]
\begin{enumerate}
    \item \textbf{Confondre les conventions Générateur / Récepteur :} Écrire $U = E + rI$ pour un générateur ou $U = E' - r'I$ pour un récepteur actif.
        \cle{Règle} : Générateur (source) : $U = E - rI$ (pente $-$). Récepteur (charge) : $U = E' + r'I$ (pente $+$). Toujours définir le sens de $I$ (entrant par $+$).
    \item \textbf{Oublier la condition $E > E'$ pour le point de fonctionnement :} Calculer $I = (E-E')/(r+r')$ sans vérifier le signe. Si $E < E'$, $I$ est négatif : le « générateur » se fait charger, ce n'est pas le régime demandé. \textbf{Écrire : « Condition $E > E'$ vérifiée, donc $I>0$ »}.
    \item \textbf{Prendre les points de mesure bruts au lieu de la droite tracée :} Pour déterminer $r$ ou $r'$, on doit utiliser la \textbf{droite de régression tracée à la règle} (deux points \emph{sur la droite}, loin l'un de l'autre), pas deux points expérimentaux quelconques (incertitudes).
    \item \textbf{Erreur d'unités / Puissances de 10 :} Mélanger \si{mA} et \si{A}, \si{mV} et \si{V} dans la même formule (ex: $r = \Delta U / \Delta I$ avec $U$ en V et $I$ en mA $\rightarrow$ $r$ en \si{k\ohm} au lieu de \si{\ohm}). \textbf{Toujours convertir en SI (A, V, \si{\ohm}) avant de calculer.}
    \item \textbf{Confondre $r$ (résistance interne) et $R$ (résistance de charge) :} Dans le bilan de puissance, $P_{perdues} = r I^2$ (dans le générateur) et $P_{utiles} = R I^2$ (dans la charge) ou $E'I$ (moteur). Ne pas écrire $P_{charge} = E I$. La puissance fournie par le générateur est $P_{gen} = E I$ ; la puissance absorbée par le récepteur est $P_{rec} = U I$.
\end{enumerate}
\end{attention}

\newpage

\coursec{Exercices}

\courssub{Je vérifie mes connaissances}

\begin{exercice}[QCM : générateur ou récepteur ?]
Pour chaque affirmation, choisir la (ou les) bonne(s) réponse(s) en justifiant brièvement.
\begin{enumerate}
\item Un générateur électrique est un dipôle qui :
\begin{enumerate}
\item transforme toute l'énergie électrique en chaleur ;
\item transforme une forme d'énergie (chimique, mécanique, lumineuse...) en énergie électrique ;
\item consomme de l'énergie électrique pour la transformer en une autre forme d'énergie.
\end{enumerate}
\item La caractéristique intensité--tension d'un générateur de f.é.m. $E$ et de résistance interne $r$ s'écrit, en convention générateur :
\begin{enumerate}
\item $U = E + rI$ ; \item $U = E - rI$ ; \item $U = rI - E$.
\end{enumerate}
\item Un récepteur actif (moteur, électrolyseur, accumulateur en charge) est caractérisé par :
\begin{enumerate}
\item une f.c.é.m. $E'$ et une résistance interne $r'$ ;
\item uniquement sa résistance interne ;
\item sa f.é.m. $E$.
\end{enumerate}
\item L'intensité de court-circuit d'un générateur $(E, r)$ vaut :
\begin{enumerate}
\item $I_{cc} = \dfrac{E}{r}$ ; \item $I_{cc} = Er$ ; \item $I_{cc} = \dfrac{r}{E}$.
\end{enumerate}
\end{enumerate}
\end{exercice}

\begin{exercice}[Vrai ou faux ?]
Répondre par vrai ou faux en justifiant chaque réponse par une phrase ou une formule.
\begin{enumerate}
\item La tension aux bornes d'un générateur réel diminue quand l'intensité débitée augmente.
\item Deux générateurs identiques montés en série ont une f.é.m. équivalente égale à celle d'un seul générateur.
\item Un conducteur ohmique est un récepteur passif : il ne possède ni f.c.é.m. ni résistance interne au sens des récepteurs actifs.
\item Lorsqu'un accumulateur est rechargé, il se comporte comme un récepteur actif de f.c.é.m. $E'$.
\end{enumerate}
\end{exercice}

\begin{exercice}[Définitions et conventions]
\begin{enumerate}
\item Définir : générateur électrique ; récepteur actif ; force électromotrice ; force contre-électromotrice.
\item Donner le symbole normalisé d'un générateur idéal de tension et celui d'un récepteur actif $(E', r')$, en précisant le sens de la f.c.é.m. par rapport au courant.
\item Citer deux récepteurs passifs et deux récepteurs actifs utilisés dans la vie courante au Cameroun.
\end{enumerate}
\end{exercice}

\begin{exercice}[Calculs directs]
\begin{enumerate}
\item Une pile de f.é.m. $E = 4,5~\text{V}$ et de résistance interne $r = 1,5~\Omega$ débite un courant d'intensité $I = 0,40~\text{A}$. Calculer la tension $U$ entre ses bornes.
\item Un moteur de f.c.é.m. $E' = 6,0~\text{V}$ et de résistance interne $r' = 2,0~\Omega$ est traversé par un courant $I = 0,50~\text{A}$. Calculer la tension à ses bornes.
\item Calculer l'intensité de court-circuit de la pile de la question 1. Pourquoi faut-il éviter le court-circuit d'une batterie de mototaxi ?
\end{enumerate}
\end{exercice}

\courssub{Je m'entraîne}

\begin{exercice}[Association série d'accumulateurs]
Quatre accumulateurs identiques, chacun de f.é.m. $E_0 = 2,0~\text{V}$ et de résistance interne $r_0 = 0,10~\Omega$, sont montés en série (dans le même sens) pour alimenter une lampe de résistance $R = 12~\Omega$.
\begin{enumerate}
\item Déterminer la f.é.m. $E$ et la résistance interne $r$ du générateur équivalent.
\item Schématiser le circuit équivalent.
\item Calculer l'intensité $I$ du courant dans le circuit et la tension $U$ aux bornes de la lampe.
\end{enumerate}
\end{exercice}

\begin{exercice}[Générateur et moteur]
Un générateur $(E = 12~\text{V} \,;\, r = 0,80~\Omega)$ alimente un moteur électrique $(E' = 9,0~\text{V} \,;\, r' = 1,2~\Omega)$ monté en série avec lui.
\begin{enumerate}
\item Faire le schéma du circuit en fléchant les tensions.
\item Établir l'expression de l'intensité $I$ du courant, puis calculer sa valeur.
\item Calculer la tension $U$ aux bornes du moteur.
\item Calculer la puissance utile du moteur, la puissance dissipée par effet Joule dans tout le circuit, puis le rendement du transfert d'énergie du générateur vers le moteur.
\end{enumerate}
\end{exercice}

\begin{exercice}[Électrolyseur]
Un électrolyseur de f.c.é.m. $E' = 1,8~\text{V}$ et de résistance interne $r' = 0,60~\Omega$ est alimenté sous une tension constante $U = 3,0~\text{V}$.
\begin{enumerate}
\item Rappeler la caractéristique $U = f(I)$ d'un récepteur actif et calculer l'intensité $I$ qui traverse l'électrolyseur.
\item Calculer la puissance convertie en énergie chimique.
\item Calculer la puissance dissipée par effet Joule.
\item Vérifier la conservation de la puissance et calculer le rendement de la conversion électrochimique.
\end{enumerate}
\end{exercice}

\begin{exercice}[Piles en opposition]
Deux piles $P_1$ $(E_1 = 4,5~\text{V} \,;\, r_1 = 1,5~\Omega)$ et $P_2$ $(E_2 = 3,0~\text{V} \,;\, r_2 = 2,0~\Omega)$ sont montées en opposition (pôles positifs reliés entre eux) dans un circuit fermé ne comportant qu'elles.
\begin{enumerate}
\item Faire le schéma du montage.
\item Déterminer la f.é.m. et la résistance interne du dipôle équivalent.
\item Calculer l'intensité du courant qui circule. Quel est le rôle de chaque pile (générateur ou récepteur) ?
\item Ce montage est-il utile en pratique ? Conclure.
\end{enumerate}
\end{exercice}

\courssub{Je me prépare au Probatoire}

\begin{exercice}[Caractéristique d'une pile — type Probatoire]
Au cours d'une séance de travaux pratiques, un élève de Première C du lycée de Ngaoundéré mesure la tension $U$ aux bornes d'une pile pour différentes valeurs de l'intensité $I$ débitée, réglée à l'aide d'un rhéostat. Il obtient le tableau suivant :
\begin{center}
\begin{tabular}{|l|c|c|c|c|c|c|}
\hline
$I$ (A) & 0 & 0,10 & 0,20 & 0,30 & 0,40 & 0,50 \\
\hline
$U$ (V) & 4,50 & 4,20 & 3,90 & 3,60 & 3,30 & 3,00 \\
\hline
\end{tabular}
\end{center}
\begin{enumerate}
\item Faire le schéma annoté du montage ayant permis ces mesures (pile, rhéostat, ampèremètre, voltmètre, interrupteur).
\item Tracer la caractéristique $U = f(I)$ de la pile sur papier millimétré. Échelles : $1~\text{cm}$ pour $0,05~\text{A}$ et $1~\text{cm}$ pour $0,50~\text{V}$.
\item Justifier que la pile peut être modélisée par un générateur linéaire $(E, r)$ et écrire sa loi d'Ohm.
\item Déduire du graphe la f.é.m. $E$ et la résistance interne $r$ de la pile.
\item Calculer l'intensité de court-circuit $I_{cc}$ de cette pile. Placer le point de court-circuit sur le graphe.
\item Cette pile alimente un conducteur ohmique de résistance $R = 12~\Omega$. Déterminer le point de fonctionnement du circuit graphiquement, puis vérifier par le calcul les valeurs de $I$ et de $U$.
\end{enumerate}
\end{exercice}

\begin{exercice}[Recharge d'une batterie de téléphone solaire — type Probatoire]
Dans un village de la région de l'Extrême-Nord, un petit panneau solaire sert de chargeur pour recharger une batterie. Le chargeur est modélisé par un générateur de f.é.m. $E = 6,0~\text{V}$ et de résistance interne $r = 2,0~\Omega$. La batterie, partiellement déchargée, se comporte comme un récepteur actif de f.c.é.m. $E' = 3,7~\text{V}$ et de résistance interne $r' = 0,50~\Omega$.
\begin{enumerate}
\item Schématiser le circuit de charge en respectant les polarités et en fléchant $E$, $E'$, $U$ et $I$.
\item Montrer que l'intensité du courant de charge s'écrit $I = \dfrac{E - E'}{r + r'}$. Calculer sa valeur.
\item Calculer la tension $U$ aux bornes de la batterie pendant la charge.
\item Calculer :
\begin{enumerate}
\item la puissance totale fournie par le chargeur ;
\item la puissance chimiquement stockée par la batterie ;
\item la puissance totale perdue par effet Joule dans le circuit.
\end{enumerate}
\item Vérifier le bilan de puissance et calculer le rendement de la charge.
\item Au bout de combien de temps la batterie aura-t-elle stocké une énergie chimique de $2,0~\text{Wh}$, en supposant le courant constant ?
\end{enumerate}
\end{exercice}

\begin{exercice}[Étude graphique d'un circuit générateur--moteur — type Probatoire]
Un atelier de réparation à Douala utilise un moteur électrique alimenté par un générateur. On a tracé sur un même graphe les caractéristiques intensité--tension des deux dipôles :
\begin{itemize}
\item générateur : droite passant par les points $(I = 0 \,;\, U = 24~\text{V})$ et $(I = 4,0~\text{A} \,;\, U = 16~\text{V})$ ;
\item moteur : droite passant par les points $(I = 0 \,;\, U = 12~\text{V})$ et $(I = 4,0~\text{A} \,;\, U = 20~\text{V})$.
\end{itemize}
\begin{enumerate}
\item Reproduire le graphe et identifier, en justifiant, la caractéristique du générateur et celle du moteur.
\item Déduire du graphe la f.é.m. $E$ et la résistance interne $r$ du générateur, ainsi que la f.c.é.m. $E'$ et la résistance interne $r'$ du moteur.
\item Définir le point de fonctionnement du circuit générateur--moteur. Le déterminer graphiquement (coordonnées $I_F$ et $U_F$).
\item Retrouver $I_F$ et $U_F$ par le calcul à partir des lois d'Ohm des deux dipôles.
\item Calculer la puissance mécanique utile du moteur et le rendement global du circuit.
\item Le moteur se bloque accidentellement : sa f.c.é.m. s'annule ($E' = 0$). Calculer la nouvelle intensité dans le circuit. Pourquoi cette situation est-elle dangereuse ?
\end{enumerate}
\end{exercice}



\newpage

\coursec{Corrigés des exercices}

\begin{corrige}[Exercice 1 — QCM : générateur ou récepteur ?]
\begin{enumerate}
\item Réponse \textbf{b} : un générateur convertit une énergie non électrique (chimique pour une pile, mécanique pour un alternateur, lumineuse pour une cellule photovoltaïque) en énergie électrique. La proposition a décrit un conducteur ohmique (récepteur passif), la proposition c un récepteur.
\item Réponse \textbf{b} : $U = E - rI$. La chute de tension interne $rI$ diminue la tension disponible aux bornes quand le générateur débite un courant.
\item Réponse \textbf{a} : un récepteur actif est modélisé par une f.c.é.m. $E'$ (qui rend compte de la conversion d'énergie en une forme non thermique) en série avec une résistance interne $r'$.
\item Réponse \textbf{a} : en court-circuit $U = 0$, donc $0 = E - rI_{cc}$, d'où $I_{cc} = \dfrac{E}{r}$. Vérification dimensionnelle : $\dfrac{[\text{V}]}{[\Omega]} = [\text{A}]$ : l'expression est homogène à une intensité.
\end{enumerate}
\end{corrige}

\begin{corrige}[Exercice 2 — Vrai ou faux ?]
\begin{enumerate}
\item \textbf{Vrai} : $U = E - rI$ est une fonction affine décroissante de $I$ (pente $-r < 0$) : plus le générateur débite, plus la tension à ses bornes chute.
\item \textbf{Faux} : en série et dans le même sens, les f.é.m. s'ajoutent : $E_{\text{éq}} = 2E_0$ (et $r_{\text{éq}} = 2r_0$). La f.é.m. équivalente est donc \emph{double} de celle d'un seul générateur.
\item \textbf{Vrai} : un conducteur ohmique convertit \emph{toute} l'énergie électrique reçue en chaleur (effet Joule) ; sa caractéristique est $U = RI$, sans f.c.é.m. : c'est un récepteur passif.
\item \textbf{Vrai} : en charge, l'accumulateur reçoit de l'énergie électrique et la convertit en énergie chimique ; on le modélise alors par $(E', r')$ avec $U = E' + r'I$ : il se comporte en récepteur actif.
\end{enumerate}
\end{corrige}

\begin{corrige}[Exercice 3 — Définitions et conventions]
\begin{enumerate}
\item \textbf{Générateur électrique} : dipôle qui convertit une forme d'énergie non électrique en énergie électrique et qui maintient le courant dans un circuit. \textbf{Récepteur actif} : dipôle qui convertit une partie de l'énergie électrique reçue en une forme d'énergie autre que la chaleur (mécanique, chimique, rayonnement). \textbf{Force électromotrice} $E$ : tension aux bornes du générateur lorsqu'il ne débite aucun courant ($I = 0$) ; elle s'exprime en volts. \textbf{Force contre-électromotrice} $E'$ : tension minimale qu'il faut appliquer à un récepteur actif pour qu'il soit traversé par un courant ; elle rend compte de la conversion d'énergie non thermique.
\item Générateur idéal de tension : cercle avec les bornes $+$ et $-$ (ou symbole de la pile : trait long $=$ pôle $+$, trait court $=$ pôle $-$) ; la flèche tension $E$ est orientée du pôle $-$ vers le pôle $+$, dans le \emph{même sens} que le courant débité (convention générateur). Récepteur actif : f.c.é.m. $E'$ en série avec $r'$ ; la flèche $E'$ est orientée en sens \emph{opposé} au courant qui le traverse (convention récepteur), d'où $U = E' + r'I$.
\item Récepteurs passifs : lampe à incandescence, fer à repasser, résistance de cuisinière électrique. Récepteurs actifs : moteur de moulin à maïs, ventilateur, batterie de téléphone en charge, électrolyseur.
\end{enumerate}
\end{corrige}

\begin{corrige}[Exercice 4 — Calculs directs]
\begin{enumerate}
\item Loi d'Ohm du générateur : $U = E - rI = 4,5 - 1,5 \times 0,40 = 4,5 - 0,60 = 3,9~\text{V}$.
\item Loi d'Ohm du récepteur actif : $U = E' + r'I = 6,0 + 2,0 \times 0,50 = 7,0~\text{V}$.
\item $I_{cc} = \dfrac{E}{r} = \dfrac{4,5}{1,5} = 3,0~\text{A}$. Pour une batterie de mototaxi, $r$ est très faible (quelques centièmes d'ohm) : $I_{cc}$ atteint des centaines d'ampères, d'où un dégagement de chaleur considérable (puissance Joule $rI^2$) pouvant provoquer un incendie, la dégradation des électrodes ou l'explosion de la batterie.
\end{enumerate}
\end{corrige}

\begin{corrige}[Exercice 5 — Association série d'accumulateurs]
\begin{enumerate}
\item En série et dans le même sens, les f.é.m. et les résistances internes s'ajoutent :
\[ E = 4E_0 = 4 \times 2,0 = 8,0~\text{V} \qquad r = 4r_0 = 4 \times 0,10 = 0,40~\Omega. \]
\item Circuit équivalent : générateur unique $(E, r)$ en série avec la lampe $R$ ; ampèremètre en série pour mesurer $I$, voltmètre en dérivation aux bornes de la lampe.
\item Loi de Pouillet :
\[ I = \dfrac{E}{r + R} = \dfrac{8,0}{0,40 + 12} = \dfrac{8,0}{12,4} \approx 0,65~\text{A}. \]
Tension aux bornes de la lampe : $U = RI = 12 \times 0,645 \approx 7,7~\text{V}$. Vérification croisée : $U = E - rI = 8,0 - 0,40 \times 0,645 \approx 7,7~\text{V}$. Les deux calculs concordent.
\end{enumerate}
\end{corrige}

\begin{corrige}[Exercice 6 — Générateur et moteur]
\begin{enumerate}
\item Schéma : générateur $(E, r)$ et moteur $(E', r')$ en série ; le courant $I$ sort du pôle $+$ du générateur et entre par la borne $+$ du moteur ; flèche $E$ dans le sens de $I$ (convention générateur), flèche $E'$ opposée à $I$ (convention récepteur).
\item Loi des mailles : $E - rI = E' + r'I$, d'où
\[ I = \dfrac{E - E'}{r + r'} = \dfrac{12 - 9,0}{0,80 + 1,2} = \dfrac{3,0}{2,0} = 1,5~\text{A}. \]
\item $U = E' + r'I = 9,0 + 1,2 \times 1,5 = 10,8~\text{V}$.
\item Puissance utile (mécanique) du moteur : $P_u = E'I = 9,0 \times 1,5 = 13,5~\text{W}$. Puissance dissipée par effet Joule dans tout le circuit : $P_J = (r + r')I^2 = 2,0 \times (1,5)^2 = 4,5~\text{W}$. Puissance fournie par le générateur : $P_g = EI = 12 \times 1,5 = 18~\text{W}$. Bilan : $P_u + P_J = 13,5 + 4,5 = 18~\text{W} = P_g$ : la conservation de la puissance est vérifiée. Rendement du transfert :
\[ \eta = \dfrac{P_u}{P_g} = \dfrac{13,5}{18} = 0,75, \quad \text{soit } 75~\%. \]
\end{enumerate}
\end{corrige}

\begin{corrige}[Exercice 7 — Électrolyseur]
\begin{enumerate}
\item Caractéristique d'un récepteur actif : $U = E' + r'I$. On en tire :
\[ I = \dfrac{U - E'}{r'} = \dfrac{3,0 - 1,8}{0,60} = \dfrac{1,2}{0,60} = 2,0~\text{A}. \]
\item Puissance convertie en énergie chimique : $P_{\text{chim}} = E'I = 1,8 \times 2,0 = 3,6~\text{W}$.
\item Puissance dissipée par effet Joule : $P_J = r'I^2 = 0,60 \times (2,0)^2 = 0,60 \times 4,0 = 2,4~\text{W}$.
\item Puissance électrique reçue : $P = UI = 3,0 \times 2,0 = 6,0~\text{W}$. Or $P_{\text{chim}} + P_J = 3,6 + 2,4 = 6,0~\text{W} = P$ : la conservation de la puissance est vérifiée. Rendement de la conversion électrochimique :
\[ \eta = \dfrac{E'I}{UI} = \dfrac{E'}{U} = \dfrac{1,8}{3,0} = 0,60, \quad \text{soit } 60~\%. \]
\end{enumerate}
\end{corrige}

\begin{corrige}[Exercice 8 — Piles en opposition]
\begin{enumerate}
\item Schéma : les deux pôles $+$ sont reliés entre eux, les deux pôles $-$ sont reliés par le reste du circuit ; dans la maille, les flèches $E_1$ et $E_2$ sont en sens opposés.
\item En opposition, la f.é.m. la plus grande impose le sens du courant : $E_{\text{éq}} = E_1 - E_2 = 4,5 - 3,0 = 1,5~\text{V}$ (sens imposé par $P_1$). Les résistances internes restent en série : $r_{\text{éq}} = r_1 + r_2 = 1,5 + 2,0 = 3,5~\Omega$.
\item $I = \dfrac{E_1 - E_2}{r_1 + r_2} = \dfrac{1,5}{3,5} \approx 0,43~\text{A}$. Le courant sort du pôle $+$ de $P_1$ : $P_1$ fonctionne en \textbf{générateur}. Le courant entre par le pôle $+$ de $P_2$ : $P_2$ fonctionne en \textbf{récepteur} (traversée à contre-sens, elle s'échauffe).
\item Ce montage n'a aucune utilité pratique : la f.é.m. disponible est faible, l'énergie de $P_1$ est en grande partie dissipée en chaleur dans les résistances internes, et une pile non rechargeable traversée à contre-sens peut être détruite (dégagement de gaz, fuite d'électrolyte). Il ne sert qu'à illustrer la notion de f.é.m. équivalente et le double rôle possible d'un dipôle.
\end{enumerate}
\end{corrige}

\begin{corrige}[Exercice 9 — Caractéristique d'une pile — type Probatoire]
\textbf{Barème indicatif (sur 10)} : schéma 1,5 pt ; tracé 2 pts ; justification de la modélisation 1 pt ; détermination de $E$ et $r$ 2 pts ; $I_{cc}$ 1 pt ; point de fonctionnement (graphique + calcul) 2,5 pts.
\begin{enumerate}
\item Schéma : pile, interrupteur, rhéostat et ampèremètre en série ; voltmètre en dérivation aux bornes de la pile (bornes V et COM respectées). \emph{Point de vigilance : l'ampèremètre se branche en série, le voltmètre en dérivation ; inverser les deux est une erreur classique sanctionnée.}
\item Les points sont alignés sur une droite décroissante coupant l'axe des ordonnées en $U = 4,50~\text{V}$. \emph{Point de vigilance : tracer la droite moyenne passant au plus près des points, et non relier les points entre eux.}
\item Les points expérimentaux sont alignés : $U$ est une fonction affine décroissante de $I$, de la forme $U = E - rI$. C'est la loi d'Ohm d'un générateur linéaire : la pile peut être modélisée par $(E, r)$.
\item Ordonnée à l'origine : $E = 4,50~\text{V}$. Pente de la droite :
\[ -r = \dfrac{\Delta U}{\Delta I} = \dfrac{3,00 - 4,50}{0,50 - 0} = \dfrac{-1,50}{0,50} = -3,0~\text{V}\cdot\text{A}^{-1}, \quad \text{donc } r = 3,0~\Omega. \]
\emph{Point de vigilance : la pente se calcule avec deux points éloignés de la droite tracée, pas nécessairement des points du tableau ; ne pas oublier le signe.}
\item Court-circuit : $U = 0 \Rightarrow I_{cc} = \dfrac{E}{r} = \dfrac{4,50}{3,0} = 1,5~\text{A}$. Le point $(1,5~\text{A} \,;\, 0~\text{V})$ se place sur le prolongement de la droite, sur l'axe des abscisses.
\item On trace sur le même graphe la droite $U = RI = 12I$, passant par l'origine et par exemple par le point $(0,30~\text{A} \,;\, 3,6~\text{V})$. L'intersection avec la caractéristique de la pile donne le point de fonctionnement : $I_F \approx 0,30~\text{A}$ et $U_F \approx 3,6~\text{V}$. Vérification par le calcul :
\[ I = \dfrac{E}{r + R} = \dfrac{4,50}{3,0 + 12} = \dfrac{4,50}{15} = 0,30~\text{A}, \qquad U = RI = 12 \times 0,30 = 3,6~\text{V}. \]
Les deux méthodes concordent parfaitement. \emph{Point de vigilance : au point de fonctionnement, le générateur et le récepteur ont même tension et même courant : c'est cette double condition qui justifie l'intersection des caractéristiques.}
\end{enumerate}
\end{corrige}

\begin{corrige}[Exercice 10 — Recharge d'une batterie de téléphone solaire — type Probatoire]
\textbf{Barème indicatif (sur 10)} : schéma fléché 1,5 pt ; démonstration de l'expression de $I$ 2 pts ; calcul de $U$ 1 pt ; puissances 3 pts ; bilan et rendement 1,5 pt ; durée de charge 1 pt.
\begin{enumerate}
\item Schéma : pôle $+$ du chargeur relié au pôle $+$ de la batterie ; le courant $I$ sort du $+$ du chargeur et entre par le $+$ de la batterie ; flèche $E$ dans le sens de $I$ (convention générateur), flèche $E'$ en sens opposé (convention récepteur), tension $U$ fléchée aux bornes de la batterie en convention récepteur. \emph{Point de vigilance : citer explicitement la condition $E > E'$ pour que le courant circule dans le sens de la charge.}
\item Loi des mailles dans le sens du courant : $E - rI - r'I - E' = 0$, soit $E - E' = (r + r')I$, d'où
\[ I = \dfrac{E - E'}{r + r'} = \dfrac{6,0 - 3,7}{2,0 + 0,50} = \dfrac{2,3}{2,5} = 0,92~\text{A}. \]
\item Tension aux bornes de la batterie : $U = E' + r'I = 3,7 + 0,50 \times 0,92 = 3,7 + 0,46 = 4,16~\text{V} \approx 4,2~\text{V}$.
\item \begin{enumerate}
\item Puissance totale fournie par le chargeur : $P_{\text{chargeur}} = EI = 6,0 \times 0,92 = 5,5~\text{W}$ ;
\item Puissance chimiquement stockée : $P_{\text{chim}} = E'I = 3,7 \times 0,92 \approx 3,4~\text{W}$ ;
\item Puissance perdue par effet Joule : $P_J = (r + r')I^2 = 2,5 \times (0,92)^2 = 2,5 \times 0,846 \approx 2,1~\text{W}$.
\end{enumerate}
\item Bilan : $P_{\text{chim}} + P_J = 3,4 + 2,1 = 5,5~\text{W} = P_{\text{chargeur}}$ : la puissance se conserve. Rendement de la charge :
\[ \eta = \dfrac{E'I}{EI} = \dfrac{E'}{E} = \dfrac{3,7}{6,0} \approx 0,62, \quad \text{soit environ } 62~\%. \]
\item Énergie chimique stockée : $W = E'I\,t$, donc
\[ t = \dfrac{W}{E'I} = \dfrac{2,0~\text{Wh}}{3,4~\text{W}} \approx 0,59~\text{h} \approx 35~\text{min}. \]
\emph{Point de vigilance : utiliser la puissance chimique $E'I$ et non la puissance totale $EI$ ; avec $W$ en Wh et $P$ en W, le temps s'obtient directement en heures.}
\end{enumerate}
\end{corrige}

\begin{corrige}[Exercice 11 — Étude graphique d'un circuit générateur--moteur — type Probatoire]
\textbf{Barème indicatif (sur 10)} : identification des caractéristiques 1,5 pt ; détermination de $E$, $r$, $E'$, $r'$ 2,5 pts ; définition et lecture du point de fonctionnement 2 pts ; vérification par le calcul 1,5 pt ; puissances et rendement 1,5 pt ; moteur bloqué 1 pt.
\begin{enumerate}
\item La caractéristique du générateur est la droite \emph{décroissante} : $U = E - rI$, la tension diminue quand l'intensité augmente. Celle du moteur est la droite \emph{croissante} : $U = E' + r'I$. \emph{Point de vigilance : la justification par le signe de la pente est exigée, pas seulement l'observation du graphe.}
\item Générateur : ordonnée à l'origine $E = 24~\text{V}$ ; pente :
\[ -r = \dfrac{16 - 24}{4,0 - 0} = \dfrac{-8,0}{4,0} = -2,0~\text{V}\cdot\text{A}^{-1}, \quad \text{donc } r = 2,0~\Omega. \]
Moteur : ordonnée à l'origine $E' = 12~\text{V}$ ; pente :
\[ r' = \dfrac{20 - 12}{4,0 - 0} = \dfrac{8,0}{4,0} = 2,0~\Omega. \]
\item Le point de fonctionnement est le point d'intersection des deux caractéristiques : le générateur et le moteur, associés en série, sont soumis à la même tension et traversés par le même courant. Graphiquement, l'intersection se lit à $I_F = 3,0~\text{A}$ et $U_F = 18~\text{V}$.
\item Par le calcul, en égalant les deux expressions de $U$ :
\[ E - rI = E' + r'I \;\Rightarrow\; I_F = \dfrac{E - E'}{r + r'} = \dfrac{24 - 12}{2,0 + 2,0} = \dfrac{12}{4,0} = 3,0~\text{A}, \]
\[ U_F = E - rI_F = 24 - 2,0 \times 3,0 = 18~\text{V}. \]
Accord parfait avec la lecture graphique.
\item Puissance mécanique utile : $P_u = E'I_F = 12 \times 3,0 = 36~\text{W}$. Puissance fournie par le générateur : $P_g = EI_F = 24 \times 3,0 = 72~\text{W}$. Rendement global :
\[ \eta = \dfrac{P_u}{P_g} = \dfrac{36}{72} = 0,50, \quad \text{soit } 50~\%. \]
\item Moteur bloqué : $E' = 0$, d'où
\[ I' = \dfrac{E}{r + r'} = \dfrac{24}{4,0} = 6,0~\text{A}, \]
soit le double du courant normal de fonctionnement. Toute l'énergie électrique est alors convertie en chaleur : $P_J = (r + r')I'^2 = 4,0 \times 36 = 144~\text{W}$, contre $36~\text{W}$ en régime normal. Le moteur chauffe fortement, son bobinage peut être détruit et un incendie peut se déclarer : d'où l'importance des fusibles et disjoncteurs de protection dans les ateliers.
\end{enumerate}
\end{corrige}

\newpage

\coursec{Fiche bilan}

\courssub{Fiche Bilan : Générateur et récepteur électriques}

\begin{popfigure}
\begin{tikzpicture}[
    mindmap,
    grow cyclic,
    every node/.style={concept, font=\small\bfseries, text=white, execute at begin node=\strut},
    concept color=popblue!80,
    level 1/.append style={level distance=4.2cm, sibling angle=60, font=\small\bfseries},
    level 2/.append style={level distance=3.0cm, sibling angle=40, font=\footnotesize}
]
\node[concept] {Générateur \& Récepteur}
    child[concept color=popblue] { node[concept] {Générateur}
        child { node[concept] {Modèle Thévenin} }
        child { node[concept] {Caract. $U=f(I)$} }
        child { node[concept] {$E$, $r$} }
    }
    child[concept color=poporange] { node[concept] {Associations}
        child { node[concept] {Série} }
        child { node[concept] {Dérivation} }
        child { node[concept] {Opposition} }
    }
    child[concept color=popgreen] { node[concept] {Récepteur}
        child { node[concept] {Passif / Actif} }
        child { node[concept] {Caract. $U=f(I)$} }
        child { node[concept] {$E'$, $r'$} }
    }
    child[concept color=poppurple] { node[concept] {Pt Fonctionnement}
        child { node[concept] {Intersection} }
        child { node[concept] {Puissances} }
        child { node[concept] {Rendement $\eta$} }
    }
    child[concept color=poppink] { node[concept] {Méthodes Exp.}
        child { node[concept] {Montage Gén.} }
        child { node[concept] {Montage Récep.} }
    }
    child[concept color=popgold] { node[concept] {Applications CM}
        child { node[concept] {Barrages (Lom Pangar)} }
        child { node[concept] {Réseau SONATREL} }
    };
\end{tikzpicture}
\legende{Carte mentale récapitulative de la leçon 14.}
\end{popfigure}

\begin{bilan}

\textbf{Définitions clés}
\begin{itemize}
    \item \cle{Générateur électrique} : Dipôle actif qui \emph{fournit} de l'énergie électrique au circuit. Modélisé par le \cle{modèle de Thévenin} : une source de tension idéale de f.é.m. $E$ (en V) en série avec une résistance interne $r$ (en $\Omega$).
    \item \cle{Récepteur électrique} : Dipôle qui \emph{reçoit} de l'énergie électrique. \cle{Récepteur passif} : résistance $R$ (effet Joule seul). \cle{Récepteur actif} : modélisé par une f.é.m. de contre-électromotrice $E'$ (en V) en série avec une résistance interne $r'$ (en $\Omega$) — ex. : moteur, batterie en charge.
    \item \cle{Caractéristique $U=f(I)$} : Représentation graphique de la tension $U$ aux bornes du dipôle en fonction du courant $I$ qui le traverse (convention récepteur : $I$ entre par la borne $+$).
    \item \cle{Point de fonctionnement} : Couple $(U_0, I_0)$ solution du système formé par les caractéristiques du générateur et du récepteur. C'est l'intersection des deux droites.
\end{itemize}

\textbf{Formules à connaître par cœur}
\begin{center}
\begin{tabularx}{\linewidth}{|l|X|l|}\hline
\rowcolor{popblueL}
\textbf{Notion} & \textbf{Relation / Formule} & \textbf{Commentaire / Unités} \\ \hline
Générateur (Thévenin) & $U = E - r I$ & $E$ : f.é.m. (V), $r$ : résist. interne ($\Omega$), pente $=-r$ \\ \hline
Récepteur actif & $U = E' + r' I$ & $E'$ : c.-é.m. (V), $r'$ : résist. interne ($\Omega$), pente $=+r'$ \\ \hline
Assoc. Série ($n$ gén.) & $E_{\text{eq}} = \sum E_i$ ; $r_{\text{eq}} = \sum r_i$ & Courant $I$ identique, tensions s'additionnent \\ \hline
Assoc. Dérivation ($n$ gén. identiques) & $E_{\text{eq}} = E$ ; $r_{\text{eq}} = \dfrac{r}{n}$ & Tension $U$ identique, courants se partagent \\ \hline
Assoc. Opposition ($E_1 > E_2$) & $E_{\text{eq}} = E_1 - E_2$ ; $r_{\text{eq}} = r_1 + r_2$ & Générateurs tête-bêche \\ \hline
Puissance fournie (Gén.) & $P_f = E I = U I + r I^2$ & $P_f = P_u + P_J$ (W) \\ \hline
Puissance utile (Récep.) & $P_u = U I = E' I + r' I^2$ & Énergie convertie (mécanique, chimique...) \\ \hline
Rendement & $\eta = \dfrac{P_u}{P_f} = \dfrac{U}{E} = 1 - \dfrac{r I}{E}$ & Sans unité, $0 \le \eta \le 1$ (souvent en \%) \\ \hline
\end{tabularx}
\end{center}

\textbf{Méthodes express}
\begin{enumerate}
    \item \textbf{Tracer la caractéristique d'un générateur} : Monter le circuit (générateur + rhéostat + ampèremètre en série, voltmètre en parallèle aux bornes). Varier le rhéostat, relever $(I_k, U_k)$. Tracer les points $\to$ droite. $E = U(I=0)$ (ordonnée à l'origine), $r = -\text{pente}$.
    \item \textbf{Caractériser une association de générateurs} : Identifier le type (série, dérivation, opposition). Appliquer les formules du tableau ci-dessus pour $E_{\text{eq}}$ et $r_{\text{eq}}$. Tracer la droite $U = E_{\text{eq}} - r_{\text{eq}} I$.
    \item \textbf{Tracer la caractéristique d'un récepteur actif} : Utiliser un générateur auxiliaire variable en série avec le récepteur et un ampèremètre. Varier $E_{\text{aux}}$, relever $(I_k, U_k)$. Tracer $\to$ droite. $E' = U(I=0)$, $r' = +\text{pente}$.
    \item \textbf{Déterminer le point de fonctionnement} :
    \begin{itemize}
        \item \emph{Calcul} : Résoudre $\begin{cases} U = E - r I \\ U = E' + r' I \end{cases} \Rightarrow I_0 = \dfrac{E - E'}{r + r'},\quad U_0 = E - r I_0$.
        \item \emph{Graphique} : Tracer les deux droites sur le même repère $(O; I, U)$. L'intersection donne $(I_0, U_0)$.
    \end{itemize}
    \item \textbf{Bilan énergétique} : Calculer $P_f = E I_0$, $P_u = U_0 I_0$, $P_{J,\text{tot}} = (r+r') I_0^2$. Vérifier $P_f = P_u + P_{J,\text{tot}}$. Calculer $\eta = P_u / P_f$.
\end{enumerate}

\end{bilan}

\begin{aretenir}[Checklist avant le Probatoire]
$\square$ Je sais définir un générateur et un récepteur (passif/actif) et donner leur rôle énergétique.\par
$\square$ Je connais par cœur le modèle de Thévenin (schéma + équation $U = E - rI$) et je sais identifier $E$ et $r$ sur une caractéristique.\par
$\square$ Je sais extraire $E$ (ordonnée à l'origine) et $r$ (valeur absolue de la pente) d'un graphique $U=f(I)$ d'un générateur.\par
$\square$ Je maîtrise les formules des associations de générateurs (série, dérivation identiques, opposition) et je sais calculer $E_{\text{eq}}$ et $r_{\text{eq}}$.\par
$\square$ Je différencie récepteur passif ($U=RI$) et récepteur actif ($U = E' + r'I$) ; je connais l'équation du récepteur actif.\par
$\square$ Je sais déterminer le point de fonctionnement $(U_0, I_0)$ par le calcul (système de 2 équations) ET par la méthode graphique (intersection).\par
$\square$ Je sais calculer les puissances : fournie $P_f = EI$, utile $P_u = UI$, dissipée par effet Joule $P_J = rI^2$ (générateur) et $P_{J}' = r'I^2$ (récepteur).\par
$\square$ Je sais calculer le rendement $\eta = P_u/P_f = U/E = 1 - rI/E$ et l'exprimer en \%. Je sais qu'il est maximal pour $I \to 0$ (circuit ouvert).\par
$\square$ Je sais décrire le montage expérimental pour tracer la caractéristique d'un générateur (rhéostat, A, V) et celle d'un récepteur actif (générateur auxiliaire, A, V).\par
$\square$ Je vérifie systématiquement l'homogénéité de mes formules (analyse dimensionnelle) et j'écris les unités SI (V, A, $\Omega$, W) à chaque résultat numérique.\par
$\square$ Je respecte la convention récepteur pour les schémas et les équations (flèche $I$ entrant par la borne $+$ de la tension $U$).\par
$\square$ Je relis l'énoncé pour vérifier le sens des f.é.m. (opposition ou non) et la nature des dipôles (générateur/récepteur) avant de calculer.
\end{aretenir}

\findecours

\end{document}
